\documentclass{article}
\usepackage{amsmath}
\usepackage{amssymb}
\usepackage{amsthm}
\usepackage{geometry}
\usepackage{enumitem}
\usepackage{xcolor}
\usepackage{fancyhdr}
\usepackage[hidelinks]{hyperref}

\geometry{
    a4paper,
    margin=1in,
}

% Package for theorems
\newtheoremstyle{mythm} % name
  {\topsep} % above space
  {\topsep} % below space
  {\itshape} % body font
  {} % indent
  {\bfseries} % head font
  {.} % punctuation
  {.5em} % after head space
  {} % head spec

\theoremstyle{mythm}
% This ensures numbering follows the section (e.g., 1.1, 1.2)
\newtheorem{theorem}{Theorem}[section] 
\newtheorem{corollary}{Corollary}[theorem]
\newtheorem{definition}[theorem]{Definition}
\newtheorem{remark}[theorem]{Remark}
\newtheorem{example}[theorem]{Example}
\newtheorem{lemma}[theorem]{Lemma}
\newtheorem{exercise}{Exercise}

% --- FANCY HEADER SETUP ---
\pagestyle{fancy}
\fancyhead{} % Clear all header fields
\fancyhead[R]{DMS KU - AVP} % Right header content
\fancyhead[L]{Real Analysis Notes} % Left header content
\fancyfoot[C]{\thepage} % Centered footer with page number
\renewcommand{\headrulewidth}{0.4pt} % Line under the header
% --------------------------

\begin{document}
\begin{center}
\textbf{\Large Real Analysis Notes} \\[0.5em]
\textbf{Course Code: MSMAT01DSC04} \\[0.5em]
\textbf{Anoop V. P.}
\end{center}

\vspace{1em}

\section{Basic Topology}

\subsection{Finite, Countable and Uncountable Sets}

\begin{definition}
Let $A$ and $B$ be any two sets. Suppose there is a rule or correspondence that assigns to each element $x \in A$ exactly one element in $B$, which we denote by $f(x)$. We call this correspondence $f$ a \textit{function} from $A$ to $B$, or a \textit{mapping} of $A$ into $B$. 

In this context, the set $A$ is referred to as the \textit{domain} of $f$ (meaning $f$ is defined on $A$), and the individual elements $f(x)$ are the \textit{values} taken by $f$. The collection of all such values is called the \textit{range} of $f$.
\end{definition}

\begin{definition}
Suppose we have a mapping $f$ from a set $A$ into a set $B$. For any subset $E \subseteq A$, we define $f(E)$ as the set of all elements of the form $f(x)$, where $x \in E$. We refer to $f(E)$ as the \textit{image} of $E$ under the function $f$. 

Notice that the entire range of $f$ can be written as $f(A)$, and it naturally follows that $f(A) \subseteq B$. In the special case where $f(A) = B$ (meaning every element in the codomain $B$ is mapped to by some element in $A$), we say that $f$ maps $A$ \textit{onto} $B$. Such a function is also commonly called \textit{surjective}.
\end{definition}
\begin{definition}
Let $f$ be a mapping from a set $A$ into a set $B$. For any subset $E \subseteq B$, we use the notation $f^{-1}(E)$ to denote the set of all elements $x \in A$ which map into $E$ (that is, where $f(x) \in E$). This set $f^{-1}(E)$ is called the \textit{inverse image} (or preimage) of $E$ under $f$. 
\end{definition}

\begin{definition}
For a specific element $y \in B$, $f^{-1}(y)$ represents the set of all $x \in A$ such that $f(x) = y$. If, for every $y \in B$, the set $f^{-1}(y)$ contains at most one element from $A$, we say that $f$ is a \textit{1-1 (one-to-one)} mapping from $A$ into $B$. 

An equivalent and often more practical way to state this is: a function $f$ is 1-1 (or injective) provided that distinct inputs yield distinct outputs. Formally, $f(x_1) \neq f(x_2)$ whenever $x_1 \neq x_2$ for any $x_1, x_2 \in A$.
\end{definition}
\begin{definition}
When there exists a 1-1 mapping from a set $A$ onto a set $B$ (a bijection), we say that $A$ and $B$ can be placed in a 1-1 \textit{correspondence}. We also express this by saying $A$ and $B$ share the same \textit{cardinal number}, or simply that they are \textit{equivalent}, denoted by $A \sim B$. 

This equivalence behaves as an \textit{equivalence relation} because it naturally satisfies three key properties:
\begin{enumerate}
    \item \textbf{Reflexivity:} $A \sim A$ for any set $A$.
    \item \textbf{Symmetry:} If $A \sim B$, then $B \sim A$.
    \item \textbf{Transitivity:} If $A \sim B$ and $B \sim C$, then $A \sim C$.
\end{enumerate}
\end{definition}

\begin{definition}
Let $J$ denote the set of all positive integers, and for any given positive integer $n$, let $J_n = \{1, 2, \dots, n\}$ represent the set of the first $n$ positive integers. For an arbitrary set $A$, we define its size as follows:
\begin{enumerate}[label=(\alph*)]
    \item $A$ is \textit{finite} if there exists some $n$ such that $A \sim J_n$. (By convention, the empty set is also considered finite).
    \item $A$ is \textit{infinite} if it is not a finite set.
    \item $A$ is \textit{countable} (sometimes referred to as \textit{enumerable} or \textit{denumerable}) if $A \sim J$.
    \item $A$ is \textit{uncountable} if it is neither finite nor countable.
    \item $A$ is \textit{at most countable} if it is either finite or countable.
\end{enumerate}
\end{definition}
\begin{remark}
For two finite sets $A$ and $B$, the equivalence $A \sim B$ clearly holds if and only if $A$ and $B$ contain the exact same number of elements. However, when dealing with infinite sets, simply "counting" elements to establish cardinality does not work, and the intuitive idea of "having the same number of elements" becomes quite vague. This is why the rigorous notion of a 1-1 correspondence is crucial; it allows us to mathematically establish when two infinite sets have the "same size" or "same number of elements" while retaining complete clarity.
\end{remark}
\begin{example}
Let $A$ be the set of all integers. We can show that $A$ is countable by setting up a 1-1 correspondence with the set of positive integers $J$. Consider the following arrangement of the sets $A$ and $J$:
\begin{align*}
    A &: \quad 0, \quad 1, \quad -1, \quad 2, \quad -2, \quad 3, \quad -3, \dots \\
    J &: \quad 1, \quad 2, \quad \phantom{-}3, \quad 4, \quad \phantom{-}5, \quad 6, \quad \phantom{-}7, \dots
\end{align*}

In this case, we can provide an explicit formula for a function $f$ from $J$ to $A$ which establishes this 1-1 correspondence:
\[
f(n) = \begin{cases} 
      \dfrac{n}{2} & \text{if } n \text{ is even}, \\[1em]
      -\dfrac{n-1}{2} & \text{if } n \text{ is odd}.
   \end{cases}
\]
\end{example}
\begin{example}
Any two open intervals $(a, b)$ and $(c, d)$ have the same cardinality. We can explicitly construct a bijection $f: (a, b) \to (c, d)$ using a linear transformation:
\[
f(x) = c + \frac{d-c}{b-a}(x-a)
\]
This exact same mapping provides a 1-1 correspondence between any two closed intervals $[a, b]$ and $[c, d]$. 
(Note: For students' understanding, observe that these linear maps are continuous; they simply stretch and shift one interval to perfectly cover the other).
\end{example}

\begin{example}
Any open interval has the same cardinality as the entire set of real numbers $\mathbb{R}$. For instance, consider the open interval $(-\frac{\pi}{2}, \frac{\pi}{2})$ and $\mathbb{R}$. A natural bijection is the tangent function:
\[
f(x) = \tan(x) \quad \text{for } x \in \left(-\frac{\pi}{2}, \frac{\pi}{2}\right)
\]
(Note: This mapping is also continuous, extending the bounded interval to cover all of $\mathbb{R}$).
\end{example}

\begin{example}
An open interval and a closed interval have the same cardinality, for instance $(0, 1) \sim [0, 1]$. An elegant way to prove this is by using the Schröder-Bernstein Theorem.

\textbf{Schröder-Bernstein Theorem (without proof):} If there exists a 1-1 mapping from a set $A$ into a set $B$, and a 1-1 mapping from $B$ into $A$, then there exists a bijection between $A$ and $B$ (i.e., $A \sim B$).

To apply this to $(0, 1)$ and $[0, 1]$, we just need two 1-1 functions:
\begin{align*}
    f &: (0, 1) \to [0, 1] \quad \text{defined by } f(x) = x \\
    g &: [0, 1] \to (0, 1) \quad \text{defined by } g(x) = \frac{1}{4} + \frac{1}{2}x
\end{align*}
Since both $f$ and $g$ are 1-1, the theorem guarantees $(0,1) \sim [0,1]$.

Alternatively, we can construct an explicit bijection. Consider the sequence of points $1/2, 1/3, 1/4, \dots$ inside $(0,1)$. We can define a function $h: (0, 1) \to [0, 1]$ that maps the first two terms of this sequence to the endpoints $0$ and $1$, shifts the rest of the sequence, and leaves everything else unchanged:
\[
h(x) = \begin{cases} 
      0 & \text{if } x = 1/2 \\
      1 & \text{if } x = 1/3 \\
      \frac{1}{n-2} & \text{if } x = 1/n \text{ for } n \geq 4 \text{ ($n \in \mathbb{N}$)} \\
      x & \text{if } x \neq 1/n \text{ for any integer } n \geq 2
   \end{cases}
\]
(Observe that unlike the previous examples, this piecewise explicit bijection is not continuous).
\end{example}

\begin{remark}
In the first two examples, the bijections and their inverses are continuous. You will learn in topology that a continuous bijection with a continuous inverse is called a \textit{homeomorphism}, which effectively identifies two spaces as having the "same shape". 

However, when mapping a closed interval to an open interval (as in our third example), you cannot get a homeomorphism. In fact, no bijection between an open and closed interval can be continuous in both directions. This is due to a topological property you will encounter later: continuous functions map \textit{compact} sets (like closed, bounded intervals) to compact sets. Since an open interval is not compact, a continuous bijection from a closed interval to an open interval is impossible.
\end{remark}

\begin{remark}
Using the sequence-shifting technique or the Schröder-Bernstein theorem from the third example, 1-1 correspondences between open intervals and half-open intervals, or half-open intervals and closed intervals, can be established quite similarly. 

The ultimate takeaway from these examples is profound: the cardinality of the real numbers $\mathbb{R}$ is exactly the same as the cardinality of any interval (open, closed, or half-open), provided the interval has a non-zero length.
\end{remark}
\begin{remark}
A finite set cannot be equivalent to one of its proper subsets. That this is, however, possible for infinite sets, is demonstrated by our previous example establishing a 1-1 correspondence between the positive integers $J$ and all integers $A$, even though $J$ is a proper subset of $A$.

In fact, we could replace our earlier definition of an infinite set with the following statement: A set $A$ is infinite if $A$ is equivalent to one of its proper subsets.
\end{remark}
\begin{definition}
By a \textit{sequence}, we mean a function $f$ defined on the set $J$ of all positive integers. If we evaluate this function at an integer $n \in J$ and let $f(n) = x_n$, it is customary to denote the sequence by the symbol $\{x_n\}$, or simply by listing its outputs: $x_1, x_2, x_3, \dots$. 

The values of $f$, that is, the elements $x_n$, are called the \textit{terms} of the sequence. If $A$ is a set and if $x_n \in A$ for all $n \in J$, then $\{x_n\}$ is said to be a \textit{sequence in $A$}, or a sequence of elements of $A$. Note that unlike sets, the terms $x_1, x_2, x_3, \dots$ of a sequence need not be distinct.
\end{definition}

\begin{remark}
The concept of a sequence gives us a highly intuitive way to understand countable sets. Since a countable set has a 1-1 correspondence with the positive integers $J$, every countable set can be viewed as the range of a sequence where all terms are distinct. 

In simple words, if a set is countable, its elements can be \textit{enumerated}. This means they can be "arranged in a sequence." You can literally line the elements up and label them: "Here is the first element ($x_1$), here is the second element ($x_2$), here is the third ($x_3$)," and so on. You are guaranteed that every single element in the set will eventually appear on your list exactly once. 

(As a minor note, sometimes it is more convenient to replace $J$ in our definition with the set of all non-negative integers, meaning we simply start our sequence index with $0$ rather than $1$).
\end{remark}
\begin{theorem}
Every infinite subset of a countable set $A$ is countable.
\end{theorem}

\begin{proof}
Suppose $E \subset A$, and $E$ is an infinite set. Our goal is to show that we can arrange the elements of $E$ into a sequence, proving it is countable.

\textbf{Step 1: List the elements of the parent set.}
Since $A$ is countable, we know we can arrange all of its elements into a single sequence of distinct terms: $x_1, x_2, x_3, \dots$. 

\textbf{Step 2: Pick out the elements of $E$ in order.}
We will build a sequence just for $E$ by scanning through our list of $A$ and picking out the elements that belong to $E$. We do this by keeping track of their index positions:
\begin{itemize}
    \item Let $n_1$ be the smallest positive integer such that $x_{n_1} \in E$. This is the first element of $E$ we encounter.
    \item Having found the first element, let $n_2$ be the next smallest integer (greater than $n_1$) such that $x_{n_2} \in E$.
    \item In general, once we have found $k-1$ elements (at positions $n_1, \dots, n_{k-1}$), we find the next element by letting $n_k$ be the smallest integer greater than $n_{k-1}$ such that $x_{n_k} \in E$.
\end{itemize}

\textbf{Step 3: Establish the 1-1 correspondence.}
Because $E$ is an infinite set, this scanning process will never stop; we will never run out of elements. We can define a function $f$ from the set of positive integers $J$ to $E$ by $f(k) = x_{n_k}$ for $k = 1, 2, 3, \dots$. 

This function matches every positive integer $k$ to exactly one distinct element of $E$, exhausting all elements of $E$ in the process. This establishes a perfect 1-1 correspondence between $J$ and $E$, meaning $E$ is countable.
\end{proof}

\begin{remark}
This theorem reveals a fundamental property of infinite sets: roughly speaking, countable sets represent the "smallest" type of infinity. Because any infinite subset of a countable set simply drops down to being countable again, it is impossible for an uncountable set to exist as a subset of a countable set.
\end{remark}
\begin{definition}
Let $A$ and $\Omega$ be sets, and suppose that with each element $\alpha$ of $A$ there is associated a subset of $\Omega$ which we denote by $E_\alpha$. The set whose elements are these subsets $E_\alpha$ is denoted by $\{E_\alpha\}$. Instead of confusingly saying "a set of sets," we generally refer to this as a \textit{collection of sets} or a \textit{family of sets}.

For any such family of sets, we can generalize the operations of union and intersection:
\begin{itemize}
    \item \textbf{Arbitrary Union:} The union of the sets $E_\alpha$ is defined to be the set $S$ such that $x \in S$ if and only if $x \in E_\alpha$ for \textit{at least one} $\alpha \in A$. We use the notation:
    \[ S = \bigcup_{\alpha \in A} E_\alpha \]
    \item \textbf{Arbitrary Intersection:} The intersection of the sets $E_\alpha$ is defined to be the set $P$ such that $x \in P$ if and only if $x \in E_\alpha$ for \textit{every} $\alpha \in A$. We use the notation:
    \[ P = \bigcap_{\alpha \in A} E_\alpha \]
\end{itemize}

Finally, if $A \cap B$ is not empty, we say that $A$ and $B$ \textit{intersect}; if it is empty, they are \textit{disjoint}.
\end{definition}

\begin{remark}
To make sense of this notation, it helps to think of the set $A$ as an "index set" (like a collection of name tags). The way we write unions and intersections changes depending on how large this index set is:

\begin{itemize}
    \item \textbf{Finite Collections:} If $A$ just consists of the integers $1, 2, \dots, n$, we are dealing with a finite number of sets. We usually write the union explicitly as $S = E_1 \cup E_2 \cup \dots \cup E_n$, or compactly using a limit: $S = \bigcup_{m=1}^n E_m$. The exact same logic and notation apply to finite intersections, replacing the $U$ with $\cap$.
    
    \item \textbf{Countable Infinite Collections:} If $A$ is the set of all positive integers, we have an infinite sequence of sets. The usual notation for their union becomes $S = \bigcup_{m=1}^\infty E_m$. It is important to note that the $\infty$ symbol here simply indicates that we are taking the union of a \textit{countable} collection of sets; it should not be confused with the literal numbers $+\infty$ or $-\infty$.
    
    \item \textbf{Uncountable Arbitrary Collections:} The generalized notation $\bigcup_{\alpha \in A} E_\alpha$ is incredibly powerful because it works even if $A$ is an \textit{uncountable} set (like the real numbers). It guarantees that no matter how unimaginably large the family of sets is, we can rigorously define what it means to combine them all together or find what they all share.
\end{itemize}
\end{remark}
\begin{example}
Let $A$ be the set of real numbers $x$ such that $0 < x \le 1$. We can write this index set as $A = (0, 1]$. 

For every $x \in A$, let $E_x$ be the set of real numbers $y$ such that $0 < y < x$. In other words, each set $E_x$ is the open interval $(0, x)$. 

Then this family of sets has the following properties:
\begin{enumerate}[label=(\roman*)]
    \item $E_x \subset E_z$ if and only if $0 < x \le z \le 1$;
    \item $\displaystyle\bigcup_{x \in A} E_x = E_1$;
    \item $\displaystyle\bigcap_{x \in A} E_x$ is empty.
\end{enumerate}

\textbf{Explanation of Properties:}
\begin{itemize}
    \item \textbf{Property (i)} is visually clear. The open interval $(0, x)$ is completely contained inside the interval $(0, z)$ if and only if the upper bound $x$ is less than or equal to $z$. This shows our sets are "nested" inside one another.
    
    \item \textbf{Property (ii)} is also straightforward. Since the sets are nested, as $x$ gets larger, the interval $E_x$ gets larger. The union of all these intervals is simply the largest interval in the collection. Because $x$ can take a maximum value of $1$ (since $x \in A$), the largest set is $E_1 = (0, 1)$. Every element in the union will fall into this maximal set.
    
    \item \textbf{To prove Property (iii)}, we must show that no single real number belongs to \textit{every} set $E_x$. Suppose, for the sake of contradiction, that there is some real number $y$ in the intersection. For $y$ to be in any $E_x$, it must be strictly positive ($y > 0$). 
    
    However, no matter how tiny this positive number $y$ is, we can always find a real number $x \in A$ that is strictly smaller than $y$ (for instance, let $x = y/2$). For this specific choice of $x$, $y$ is not in $E_x$ because $y$ is not less than $x$. Since $y \notin E_x$ for this particular $x$, it cannot be in the intersection of \textit{all} $E_x$. Thus, the intersection must be empty.
\end{itemize}
\end{example}
\begin{theorem}
Let $\{E_n\}$, $n = 1, 2, 3, \dots$, be a sequence of countable sets, and put
\[
S = \bigcup_{n=1}^\infty E_n.
\]
Then $S$ is countable.
\end{theorem}

Before diving into the formal proof, it is helpful to state what this theorem means in simple terms: \textbf{a countable union of countable sets is countable.}

\begin{proof}
Since each individual set $E_n$ is countable, its elements can be arranged in a sequence. Let us denote the elements of the first set $E_1$ as $x_{11}, x_{12}, x_{13}, \dots$, the elements of the second set $E_2$ as $x_{21}, x_{22}, x_{23}, \dots$, and in general, the elements of the $n$-th set $E_n$ as $\{x_{nk}\}$ for $k = 1, 2, 3, \dots$.

\textbf{Step 1: Construct an infinite grid.} \\
We can arrange all the elements of all the sets into an infinite two-dimensional array, where the $n$-th row contains all the elements of the set $E_n$:
\begin{center}
$
\begin{array}{ccccc}
x_{11} & x_{12} & x_{13} & x_{14} & \dots \\
x_{21} & x_{22} & x_{23} & x_{24} & \dots \\
x_{31} & x_{32} & x_{33} & x_{34} & \dots \\
x_{41} & x_{42} & x_{43} & x_{44} & \dots \\
\vdots & \vdots & \vdots & \vdots & \ddots 
\end{array}
$
\end{center}
This grid contains every single element of the union set $S$.

\textbf{Step 2: Traverse the grid to form a single sequence.} \\
To show that $S$ is countable, we must show that all elements in this grid can be listed in a single, one-dimensional sequence without missing any. We cannot just read across the first row, because it is infinitely long and we would never reach the second row. 

Instead, we read the grid by tracing finite diagonals (moving from bottom-left to top-right along diagonals). This gives us the sequence:
\[
x_{11}; \quad x_{21}, x_{12}; \quad x_{31}, x_{22}, x_{13}; \quad x_{41}, x_{32}, x_{23}, x_{14}; \quad \dots
\]

\textbf{Step 3: Address duplicate elements.} \\
The sets $E_n$ might not be mutually disjoint; they might share elements. If they do, some elements will appear more than once in our new diagonal sequence. 

If we move through our sequence and simply skip any element we have already listed, we create a 1-1 correspondence between the unique elements of $S$ and a subset (let's call it $T$) of the positive integers. This establishes that $S$ is \textit{at most countable}.

\textbf{Step 4: Conclude using previous results.} \\
We know that the first set $E_1$ is completely contained within $S$ (that is, $E_1 \subset S$). Since $E_1$ is an infinite set, $S$ must also be an infinite set. 

From Step 3, $S$ is equivalent to $T$, which is a subset of the positive integers. Since $S$ is infinite, $T$ must be an \textit{infinite} subset of the positive integers. By our previous theorem, \textbf{every infinite subset of a countable set is countable}. Because $T$ is an infinite subset of the countable set of positive integers, $T$ is countable, and therefore $S$ is countable.
\end{proof}
\begin{corollary}
Suppose $A$ is at most countable, and, for every $\alpha \in A$, $B_\alpha$ is at most countable. Put
\[ T = \bigcup_{\alpha \in A} B_\alpha. \]
Then $T$ is at most countable.
\end{corollary}

\begin{proof}
The text provides a very compact proof for this, stating that $T$ is equivalent to a subset of the previously defined set $S$ from equation 15. Here is how that deduction works.

The previous theorem proved that a strictly countable union of strictly countable sets is countable. This corollary relaxes those conditions to "at most countable," meaning the index set $A$ might be finite, and the individual sets $B_\alpha$ might also be finite or even empty. 

Because $A$ is at most countable, we can arrange its elements into a sequence $\alpha_1, \alpha_2, \dots$. This allows us to list our sets as $B_{\alpha_1}, B_{\alpha_2}, \dots$. If $A$ is finite, this list eventually stops. We can easily turn it into an infinite sequence by simply defining all the remaining sets in the sequence as empty sets. Adding empty sets to a union does not change the final set $T$.

Similarly, because each $B_\alpha$ is at most countable, it is either equivalent to a fully countable set or it is a subset of one. Therefore, when we take the union of all these $B_\alpha$ sets to form $T$, we are essentially building a set that is contained entirely within the union $S$ from equation 15. 

Because the previous theorem definitively proved that the larger set $S$ is countable, and $T$ is merely a subset of $S$, it logically follows that $T$ must be at most countable.
\end{proof}
\begin{theorem}
Let $A$ be a countable set, and let $B_n$ be the set of all $n$-tuples $(a_1, \dots, a_n)$, where $a_k \in A$ ($k = 1, \dots, n$), and the elements $a_1, \dots, a_n$ need not be distinct. Then $B_n$ is countable.
\end{theorem}

\begin{proof}
We can prove this theorem using the principle of mathematical induction. 

First, consider the base case where $n=1$. The set $B_1$ consists of all 1-tuples, which are simply the individual elements of $A$. Since $B_1$ is exactly the same set as $A$, and we are given that $A$ is countable, it is evident that $B_1$ must also be countable.

Next, we move to the inductive step. We assume that the theorem holds true for some integer $n-1$. This means we assume the set $B_{n-1}$, which contains all $(n-1)$-tuples, is a countable set. We must show that this assumption implies the set $B_n$ is also countable.

Any $n$-tuple in $B_n$ can be constructed by taking an $(n-1)$-tuple and appending one more element to it. Therefore, an element in $B_n$ can be written in the form $(b, a)$, where $b$ represents an entire $(n-1)$-tuple from the set $B_{n-1}$, and $a$ is a single element from our original set $A$.

Let us fix one specific $(n-1)$-tuple, which we will call $b$. We can form a collection of pairs by keeping this $b$ constant and pairing it with every possible element $a$ from the set $A$. Because the set $A$ is countable, this specific collection of pairs $(b, a)$ is equivalent in size to $A$, making it a countable set itself.

To construct the entire set $B_n$, we must perform this pairing process for every single $(n-1)$-tuple $b$ that exists in $B_{n-1}$. We are essentially taking a union of all these individual, countable collections of pairs. The number of these collections is determined by the number of elements in $B_{n-1}$. Since our induction hypothesis established that $B_{n-1}$ is a countable set, we are combining a countable number of these collections.

Consequently, $B_n$ is formed by taking a countable union of countable sets. As established by our previous theorem, a countable union of countable sets is always countable. Thus, $B_n$ is countable, and the proof follows by induction.
\end{proof}
\begin{corollary}
The set of all rational numbers is countable.
\end{corollary}

\begin{proof}
The text provides a concise proof by referencing our previous theorem about $n$-tuples. Here is how that connection works in detail.

Recall that a rational number is defined as any number that can be expressed as a fraction $b/a$, where both $b$ and $a$ are integers (with $a$ not equal to zero). This means every single rational number can be represented by a pair of integers $(b, a)$. 

Our previous theorem established that if a base set is countable, then the set of all $n$-tuples formed from that base set is also countable. If we let our base set be the set of all integers, which we already proved is a countable set, and we set $n = 2$, the theorem guarantees that the set of all possible integer pairs $(b, a)$ is countable.

Because every rational number is formed by one of these integer pairs, the set of all fractions $b/a$ is essentially a subset of this larger set of pairs. We simply discard the pairs where the denominator $a$ is zero, and we recognize that some pairs represent the same fraction (like $1/2$ and $2/4$). Since the entire collection of integer pairs is countable, this sub-collection of valid fractions must also be countable. Therefore, the set of all rational numbers is countable.
\end{proof}
% If you haven't already, make sure to add this to your preamble:
% \newtheorem{exercise}{Exercise}

\begin{exercise}
Prove that the set of all algebraic numbers is countable.
\end{exercise}

\begin{proof}
Before proving the statement, we must first define what an algebraic number is. An algebraic number is any complex number $x$ that is a root of a non-zero polynomial equation with integer coefficients. In other words, $x$ is a solution to an equation of the form 
\[ a_n x^n + a_{n-1} x^{n-1} + \dots + a_1 x + a_0 = 0 \]
where the degree $n \ge 1$, all the coefficients $a_0, a_1, \dots, a_n$ are integers, and the leading coefficient $a_n$ is not zero.

To prove that the set of all such numbers is countable, we can organize them by the degree $n$ of their corresponding polynomials. Let us fix a specific degree $n$. A polynomial of this degree is completely determined by its coefficients, which we can list as an $(n+1)$-tuple of integers: $(a_0, a_1, \dots, a_n)$. 

Because the set of all integers is countable, our earlier theorem regarding tuples guarantees that the set of all possible $(n+1)$-tuples of integers is also countable. Therefore, there are only countably many polynomials of degree $n$ with integer coefficients.

Furthermore, a fundamental theorem of algebra states that any polynomial of degree $n$ has at most $n$ distinct roots. This means each polynomial only yields a finite set of algebraic numbers. To gather all algebraic numbers associated with degree $n$, we take the union of the roots from each of these polynomials. Since we have a countable number of polynomials, and each has a finite number of roots, we are taking a countable union of finite sets. This guarantees that the set of all algebraic numbers of a fixed degree $n$ is countable.

Finally, to construct the entire set of all algebraic numbers, we must combine the sets of roots for every possible degree $n = 1, 2, 3, \dots$. Since the set of algebraic numbers for each individual degree $n$ is countable, combining all of them together means we are taking a countable union of countable sets. As established by our previous theorem, a countable union of countable sets is always countable. Thus, the set of all algebraic numbers is countable.
\end{proof}
\begin{theorem}
Let $A$ be the set of all sequences whose elements are the digits $0$ and $1$. This set $A$ is uncountable. 
(Note: The elements of $A$ are sequences like $1, 0, 0, 1, 0, 1, 1, 1, \dots$).
\end{theorem}

\begin{proof}
We will prove this by showing that no countable list can possibly contain every sequence in $A$.

\textbf{Step 1: Assume a countable subset.} \\
Let $E$ be any countable subset of $A$. Because $E$ is countable, we can enumerate its elements in a sequence: $s_1, s_2, s_3, \dots$. Keep in mind that each element $s_n$ in this list is itself an infinite sequence of 0s and 1s.

\textbf{Step 2: Construct a new sequence (Cantor's Diagonal Process).} \\
We will now construct a brand new sequence, which we will call $s$. We determine the digits of $s$ one by one using a famous technique known as \textbf{Cantor's diagonal process}:
\begin{itemize}
    \item Look at the $1$st digit of the first sequence $s_1$. If it is a $1$, make the $1$st digit of $s$ a $0$. If it is a $0$, make it a $1$.
    \item Look at the $2$nd digit of the second sequence $s_2$. If it is a $1$, make the $2$nd digit of $s$ a $0$, and vice versa.
    \item In general, for any integer $n$, look at the $n$-th digit of the sequence $s_n$. If that digit is a $1$, we define the $n$-th digit of our new sequence $s$ to be $0$, and vice versa.
\end{itemize}

\textbf{Step 3: Show the new sequence is missing from the list.} \\
Because our new sequence $s$ consists entirely of 0s and 1s, it is clearly a valid element of the overall set $A$. But is it in our countable subset $E$? 
By our construction, $s$ differs from the first sequence $s_1$ in the first digit. It differs from the second sequence $s_2$ in the second digit. In general, it differs from every sequence $s_n$ in our list at the $n$-th digit. Therefore, $s$ is not anywhere in the set $E$. 

\textbf{Step 4: Conclusion.} \\
Since $s \in A$ but $s \notin E$, it means $E$ is missing at least one element of $A$. Thus, $E$ is a \textit{proper subset} of $A$. 

We have just shown that \textit{every} countable subset of $A$ is merely a proper subset of $A$. It logically follows that $A$ itself must be uncountable. (If $A$ were countable, we could let our list $E$ be the entirety of $A$. But that would mean $A$ is a proper subset of itself, which is absurd).
\end{proof}
\begin{exercise}
Prove that the set of all real numbers $\mathbb{R}$ is uncountable.
\end{exercise}

\begin{proof}
We will prove this by contradiction, relying on a principle established earlier: every infinite subset of a countable set must itself be countable.

Assume, for the sake of contradiction, that the entire set of real numbers $\mathbb{R}$ is countable. 

Let us look at a very specific subset of the real numbers. Let $A$ be the set of all real numbers between $0$ and $1$ whose decimal representations consist exclusively of the digits $0$ and $1$ (for example, numbers like $0.1011001\dots$ or $0.0101010\dots$). Clearly, $A$ is a subset of $\mathbb{R}$.

Notice that there is a perfect 1-1 correspondence between the numbers in set $A$ and the set of all infinite sequences made up of the digits $0$ and $1$. We proved in our previous theorem (using Cantor's diagonal process) that the set of all such $0$ and $1$ sequences is uncountable. Because they map perfectly to one another, our subset $A$ must also be an uncountable set.

This creates a fatal logical flaw. Under our assumption that $\mathbb{R}$ is countable, every infinite subset of $\mathbb{R}$ must also be countable. However, we just constructed a subset $A$ that is strictly uncountable. This contradiction proves that our initial assumption was false. 

Therefore, the set of all real numbers $\mathbb{R}$ is uncountable.
\end{proof}

\begin{remark}
This profound result leads to several important consequences regarding the "size" of different collections of numbers on the real line.

\textbf{Every interval is uncountable:} 
It is a known mathematical fact that any continuous interval of real numbers (such as the open interval $(0, 1)$ or any interval $(a, b)$) has the exact same cardinality as the entire set of real numbers $\mathbb{R}$. Since we just proved that $\mathbb{R}$ is uncountable, it immediately shows that every interval on the number line is also uncountable.

\textbf{The irrational numbers are uncountable:} 
The set of real numbers $\mathbb{R}$ is simply the union of the rational numbers (which we denote as $\mathbb{Q}$) and the irrational numbers. We have already proven that the rational numbers $\mathbb{Q}$ are countable. 

If the set of irrational numbers were also countable, then $\mathbb{R}$ would be formed by taking the union of two countable sets. We know that a union of countable sets is always countable, which would mean $\mathbb{R}$ is countable. But $\mathbb{R}$ is uncountable! The only way to avoid this contradiction is if the set of irrational numbers is uncountable.

\textbf{Intervals are dominated by irrationals:} 
We can apply this exact same logic to any specific segment of the number line. Take any interval $I$. The rational numbers that happen to fall inside this interval, denoted as $I \cap \mathbb{Q}$, simply form a subset of all rational numbers, meaning they are countable. 

The rest of the numbers in the interval are irrational. If the irrationals in the interval were countable, the entire interval $I$ would be the union of two countable sets, making the interval itself countable. Since we know every interval is uncountable, its irrational part must be uncountable. In simple terms, within any interval, the rational numbers are merely a countable dusting, while the irrational numbers make up the uncountable bulk.
\end{remark}
\subsection{Metric Spaces}

\begin{definition}
A set $X$, whose elements we shall call \textit{points}, is said to be a \textit{metric space} if with any two points $p$ and $q$ of $X$ there is associated a real number $d(p, q)$, called the \textit{distance} from $p$ to $q$, such that
\begin{enumerate}[label=(\alph*)]
    \item $d(p,q) > 0$ if $p \neq q$; $d(p,p) = 0$;
    \item $d(p,q) = d(q,p)$;
    \item $d(p,q) \le d(p,r) + d(r,q)$, for any $r \in X$.
\end{enumerate}
Any function with these three properties is called a \textit{distance function}, or a \textit{metric}.
\end{definition}

\begin{example}
The most important examples of metric spaces, from our standpoint, are the euclidean spaces $\mathbb{R}^k$, especially $\mathbb{R}^1$ (the real line) and $\mathbb{R}^2$ (the complex plane); the distance in $\mathbb{R}^k$ is defined by
\[
d(\mathbf{x}, \mathbf{y}) = |\mathbf{x} - \mathbf{y}| \quad (\mathbf{x}, \mathbf{y} \in \mathbb{R}^k).
\]
By the properties of the euclidean norm (in particular the triangle inequality $|\mathbf{x} - \mathbf{z}| \leq |\mathbf{x} - \mathbf{y}| + |\mathbf{y} - \mathbf{z}|$), the conditions of Definition 1.26 are satisfied by this formula.
\end{example}
\begin{remark}
In Rudin, an open interval $(a, b)$ in $\mathbb{R}^1$ is called a \textit{segment}, while the term \textit{interval} is reserved for a closed interval $[a, b]$. Throughout these notes, we explicitly refer to them as \textit{open intervals}, \textit{closed intervals}, and \textit{half-open intervals}.
\end{remark}
\begin{definition}
If $a_i < b_i$ for $i = 1, \dots, k$, the set of all points $\mathbf{x} = (x_1, \dots, x_k) \in \mathbb{R}^k$ whose coordinates satisfy the inequalities $a_i \le x_i \le b_i$ ($1 \le i \le k$) is called a \textit{$k$-cell}. Thus, a 1-cell is a closed interval, a 2-cell is a rectangle, etc.
\end{definition}

\begin{definition}
If $\mathbf{x} \in \mathbb{R}^k$ and $r > 0$, the \textit{open ball} (or simply \textit{ball}) with center at $\mathbf{x}$ and radius $r$ is defined to be the set of all $\mathbf{y} \in \mathbb{R}^k$ such that $|\mathbf{y} - \mathbf{x}| < r$. 

Similarly, the \textit{closed ball} with center at $\mathbf{x}$ and radius $r$ is defined as the set of all $\mathbf{y} \in \mathbb{R}^k$ such that $|\mathbf{y} - \mathbf{x}| \le r$.
\end{definition}

\begin{definition}
A set $E \subset \mathbb{R}^k$ is said to be \textit{convex} if
\[
\lambda \mathbf{x} + (1 - \lambda)\mathbf{y} \in E
\]
whenever $\mathbf{x} \in E$, $\mathbf{y} \in E$, and $0 < \lambda < 1$.
\end{definition}

\begin{example}
Open and closed balls in $\mathbb{R}^k$ are convex. 

To prove this for an open ball $B$ centered at $\mathbf{x}$ with radius $r$, let $\mathbf{y}, \mathbf{z} \in B$ and choose $0 < \lambda < 1$. By definition of the open ball, we have $|\mathbf{y} - \mathbf{x}| < r$ and $|\mathbf{z} - \mathbf{x}| < r$. Applying the triangle inequality yields:
\begin{align*}
|\lambda \mathbf{y} + (1 - \lambda)\mathbf{z} - \mathbf{x}| &= |\lambda(\mathbf{y} - \mathbf{x}) + (1 - \lambda)(\mathbf{z} - \mathbf{x})| \\
&\le \lambda |\mathbf{y} - \mathbf{x}| + (1 - \lambda)|\mathbf{z} - \mathbf{x}| \\
&< \lambda r + (1 - \lambda)r \\
&= r.
\end{align*}
Since $|\lambda \mathbf{y} + (1 - \lambda)\mathbf{z} - \mathbf{x}| < r$, the point $\lambda \mathbf{y} + (1 - \lambda)\mathbf{z}$ lies inside $B$, proving $B$ is convex.

The exact same proof applies to closed balls (replacing strict inequalities with $\le$). Additionally, $k$-cells are also convex.
\end{example}
\begin{definition}
Let $X$ be a metric space. All points and sets mentioned below are understood to be elements and subsets of $X$.
\begin{enumerate}[label=(\alph*)]
    \item A \textit{neighborhood} of $p$ is a set $N_r(p)$ consisting of all $q$ such that $d(p, q) < r$, for some $r > 0$. The number $r$ is called the \textit{radius} of $N_r(p)$.
    \item A point $p$ is a \textit{limit point} of the set $E$ if every neighborhood of $p$ contains a point $q \neq p$ such that $q \in E$.
    \item If $p \in E$ and $p$ is not a limit point of $E$, then $p$ is called an \textit{isolated point} of $E$.
    \item $E$ is \textit{closed} if every limit point of $E$ is a point of $E$.
    \item A point $p$ is an \textit{interior point} of $E$ if there is a neighborhood $N$ of $p$ such that $N \subset E$.
    \item $E$ is \textit{open} if every point of $E$ is an interior point of $E$.
    \item The \textit{complement} of $E$ (denoted by $E^c$) is the set of all points $p \in X$ such that $p \notin E$.
    \item $E$ is \textit{perfect} if $E$ is closed and if every point of $E$ is a limit point of $E$.
    \item $E$ is \textit{bounded} if there is a real number $M$ and a point $q \in X$ such that $d(p, q) < M$ for all $p \in E$.
    \item $E$ is \textit{dense} in $X$ if every point of $X$ is a limit point of $E$, or a point of $E$ (or both).
\end{enumerate}
\end{definition}
\begin{example}[Neighborhoods]
Consider the set of real numbers $\mathbb{R}$ with the standard metric $d(x,y) = |x - y|$. An $\epsilon$-neighborhood of a point $x \in \mathbb{R}$ is exactly the open interval $(x - \epsilon, x + \epsilon)$. 

Similarly, in $\mathbb{R}^k$ with the standard Euclidean metric, neighborhoods of a point $\mathbf{x}$ are open balls centered at $\mathbf{x}$. 
However, the "shape" of a neighborhood depends on the chosen metric. For instance, consider $\mathbb{R}^2$ equipped with the maximum metric $d(\mathbf{x}, \mathbf{y}) = \max(|x_1 - y_1|, |x_2 - y_2|)$. In this metric space, a neighborhood (an open ball) takes the geometric shape of an open square with sides parallel to the coordinate axes.

The space itself also restricts neighborhoods. Let the metric space be $X = (0, 2]$ with the standard Euclidean distance. The $r$-neighborhood of the point $2$ is defined relative to $X$, meaning $N_r(2) = (2 - r, 2 + r) \cap X$. For example, if $r = 1$, the neighborhood is $N_1(2) = (1, 3) \cap (0, 2] = (1, 2]$. 
\end{example}

\begin{example}[Limit Points]
A limit point of a set $E$ does not actually need to belong to the set $E$. Consider $E = (0, 2]$ as a subset of the metric space $\mathbb{R}$. Every element of $E$ is a limit point of $E$. Furthermore, $0$ is also a limit point of $E$ because every open interval around $0$ will intersect $(0, 2]$, even though $0 \notin E$.

Now, change the context: suppose we consider $X = (0, 2]$ itself as our entire metric space. The candidates for limit points of any subset must come from the metric space $X$ under consideration. In this scenario, $0$ cannot be a limit point of any subset of $X$ simply because $0$ does not exist in the space $X$.
\end{example}

\begin{example}[Isolated Points]
Consider the set $E = (0, 2] \cup \{3\}$ as a subset of $\mathbb{R}$ with the standard metric. The point $3$ is an isolated point of $E$ because we can find a neighborhood around $3$ (e.g., radius $r=1/2$) that contains no other points of $E$. 

If we take the natural numbers $\mathbb{N}$ as a subset of $\mathbb{R}$, every single point in $\mathbb{N}$ is an isolated point, and $\mathbb{N}$ has absolutely no limit points in $\mathbb{R}$.

Consider the set $S = \mathbb{Q} \times (\mathbb{R} \setminus \mathbb{Q})$ (rational $x$-coordinates and irrational $y$-coordinates) as a subset of $\mathbb{R}^2$ with the standard metric. Does $S$ have any isolated points? No. Every neighborhood of any point in $S$ contains infinitely many other points of $S$. Its set of limit points is the entirety of $\mathbb{R}^2$.
\end{example}

\begin{example}[Closed Sets]
The set of natural numbers $\mathbb{N}$ is closed in $\mathbb{R}$. By definition, a set is closed if it contains all of its limit points. Since $\mathbb{N}$ has no limit points, the condition is vacuously satisfied. 

The entire space $\mathbb{R}$ is closed in $\mathbb{R}$. 

The set of rational numbers $\mathbb{Q}$ is not closed in $\mathbb{R}$. Between any two real numbers there exists a rational number, meaning we can construct a sequence of rationals converging to any irrational number. Thus, every irrational number is a limit point of $\mathbb{Q}$, but these limit points are not contained in $\mathbb{Q}$.
\end{example}

\begin{example}[Interior Points and Open Sets]
Let the metric space be $\mathbb{R}$ and consider the subset $E = (0, 2]$. Every point $x \in (0, 2)$ is an interior point. To rigorously see this, for any $x \in (0, 2)$, choose a radius $r = \min(x, 2 - x)$. The neighborhood $N_r(x)$ will lie entirely inside $(0, 2]$. However, the point $2$ is not an interior point because any neighborhood $(2 - r, 2 + r)$ in $\mathbb{R}$ will contain points strictly greater than $2$, which fall outside $E$. 

Now, consider the metric space $X = (0, 2]$ with the standard metric, and let $E = (1, 2]$ be a subset of $X$. In this relative context, $2$ \textit{is} an interior point of $E$. A neighborhood of $2$ in $X$ with radius $1$ is $(1, 2]$, which is entirely contained in $E$. Consequently, while $(1, 2]$ is not an open set in $\mathbb{R}$, it \textit{is} an open set in the metric space $X = (0, 2]$.

In general, within $\mathbb{R}$, open intervals are the most fundamental examples of open sets. Furthermore, any union of open intervals is again an open set.
\end{example}

\begin{example}[Perfect Sets]
Perfect sets are closed sets that have no isolated points (i.e., every point in the set is a limit point of the set). Standard examples in $\mathbb{R}^n$ include closed intervals (or cells) and finite unions of non-degenerate closed intervals. We will explore more intricate examples of perfect sets (such as the Cantor set) later in the course.
\end{example}

\begin{remark}[Bounded Sets]
By definition, a set $E$ in a metric space $X$ is bounded if we can find a neighborhood of some point $q \in X$ with a finite radius $M$ that completely contains $E$. If we are working in $\mathbb{R}$ and we conventionally choose the center point $q = 0$, this translates to $E \subset (-M, M)$, meaning $|x| < M$ for all $x \in E$. This perfectly recovers the standard definition of bounded subsets of real numbers you studied in your previous undergraduate courses.
\end{remark}

\begin{example}[Dense Sets]
A set $E$ is dense in $X$ if every point of $X$ is either a point of $E$ or a limit point of $E$ (or both). Equivalently, every non-empty open set in $X$ intersects $E$. 

Consider the metric space $X = (0, 2] \cup \{3\}$ with the standard metric. Let $E$ be the union of the irrationals in $(0, 2]$ and the singleton set $\{3\}$. $E$ is dense in $X$. Notice that $3 \in E$, but $3$ is not a limit point of $E$ (it is an isolated point). This perfectly illustrates why the definition states "is a limit point \textit{or} a point of $E$."

As another example, the set $\mathbb{Q} \times (\mathbb{R} \setminus \mathbb{Q})$ is dense in $\mathbb{R}^2$ with the standard metric, because every open ball in $\mathbb{R}^2$, no matter how small, will contain points with rational $x$-coordinates and irrational $y$-coordinates.
\end{example}
\begin{theorem}
Every neighborhood is an open set.
\end{theorem}

\begin{proof}
Let $E = N_r(p)$ be a neighborhood of a point $p$ with radius $r > 0$. To prove that $E$ is an open set, we must show that every point inside $E$ is an interior point. 

\textbf{Step 1: Pick an arbitrary point.} \\
Let $q$ be any point in $E$. By the definition of a neighborhood, the distance from $p$ to $q$ is strictly less than $r$, meaning $d(p, q) < r$.

\textbf{Step 2: Find a suitable radius around $q$.} \\
Because $d(p, q) < r$, there is some "leftover" distance between $q$ and the "boundary" of our neighborhood $E$. Let us define this positive leftover distance as a real number $h$:
\[ h = r - d(p, q) \]
Notice that $h > 0$ (since $d(p,q) = r - h$). We claim that the neighborhood of $q$ with radius $h$, denoted $N_h(q)$, is completely contained inside our original set $E$.

\textbf{Step 3: Verify using the triangle inequality.} \\
Take any point $s \in N_h(q)$. By definition, this means $d(q, s) < h$. We want to determine if $s$ is also inside $E$, which requires checking its distance from the original center $p$. 

Using the triangle inequality, we evaluate the distance between $p$ and $s$:
\[ d(p, s) \le d(p, q) + d(q, s) \]
Since we know $d(q, s) < h$, we can strictly bound this sum:
\[ d(p, s) < d(p, q) + h \]
Now, substitute our earlier definition of $h$ into the inequality:
\[ d(p, s) < d(p, q) + (r - d(p, q)) = r \]

    
\textbf{Step 4: Conclusion.} \\
Since $d(p, s) < r$, the point $s$ must belong to $E$. Because this is true for every point $s \in N_h(q)$, we have successfully shown that an entire neighborhood around $q$ is contained in $E$. 

Thus, $q$ is an interior point of $E$. Since $q$ was chosen arbitrarily, every point of $E$ is an interior point, proving that the neighborhood $E$ is an open set.
\end{proof}
\begin{theorem}
If $p$ is a limit point of a set $E$, then every neighborhood of $p$ contains infinitely many points of $E$.
\end{theorem}

\begin{proof}
We will prove this theorem by contradiction, breaking it down logically.

\textbf{Step 1: Assume the contrary.} \\
Suppose $p$ is a limit point of $E$, but there exists some neighborhood $N$ of $p$ which contains only a finite number of points of $E$.

\textbf{Step 2: Isolate the finite points.} \\
Let $q_1, \dots, q_n$ be those specific points in the intersection $N \cap E$ which are distinct from $p$.

\textbf{Step 3: Construct a smaller neighborhood.} \\
We evaluate the distance from our point $p$ to each of these finitely many points $q_m$. Let $r$ be the absolute smallest of these distances:
\[ r = \min_{1 \le m \le n} d(p, q_m) \]
Because we are taking the minimum of a finite set of strictly positive numbers (since $q_m \neq p$, the distance to $p$ is greater than zero), the minimum itself must be strictly positive. Thus, $r > 0$.

\textbf{Step 4: Reach a contradiction.} \\
Now consider the new, smaller neighborhood $N_r(p)$ centered at $p$ with this radius $r$. Because $r$ is less than or equal to the distance to any of the points $q_1, \dots, q_n$, this new neighborhood $N_r(p)$ contains absolutely no point $q$ of $E$ such that $q \neq p$.
    
However, the fundamental definition of a limit point demands that \textit{every} neighborhood of $p$ must contain at least one point of $E$ distinct from $p$. We have just explicitly constructed a neighborhood $N_r(p)$ that fails this requirement. This contradiction establishes the theorem.
\end{proof}
\begin{corollary}
A finite set has no limit points.
\end{corollary}

\begin{proof}
By the preceding theorem, if a point $p$ were a limit point of a set $E$, every neighborhood of $p$ would be required to contain infinitely many distinct points of $E$. However, if $E$ is finite, no neighborhood can possibly contain infinitely many points of $E$. Therefore, a finite set cannot have any limit points.
\end{proof}
\begin{theorem}
Let $\{E_\alpha\}$ be a (finite or infinite) collection of sets $E_\alpha$. Then
\[ \left(\bigcup_\alpha E_\alpha\right)^c = \bigcap_\alpha \left(E_\alpha^c\right). \] 
\end{theorem}

\begin{proof}
To prove that these two sets are equal, we must show mutual inclusion: the set on the left is a subset of the set on the right, and conversely, the set on the right is a subset of the set on the left.

Let $A = \left(\bigcup_\alpha E_\alpha\right)^c$ and let $B = \bigcap_\alpha \left(E_\alpha^c\right)$. 

\textbf{Step 1: Show that $A \subset B$.} \\
Suppose $x$ is an arbitrary element in $A$. By the definition of the complement, this means $x$ does not belong to the union $\bigcup_\alpha E_\alpha$.
For an element to be excluded from the union of a collection of sets, it must be absent from every single set in that collection. Therefore, $x \notin E_\alpha$ for any $\alpha$. 
This implies that $x$ belongs to the complement of each individual set, meaning $x \in E_\alpha^c$ for every $\alpha$. 
Because $x$ is in every such complement, it must belong to their intersection. Thus, $x \in \bigcap_\alpha \left(E_\alpha^c\right)$, which means $x \in B$. This establishes that $A \subset B$. 

\textbf{Step 2: Show that $B \subset A$.} \\
Conversely, suppose $x$ is an arbitrary element in $B$. This means $x \in \bigcap_\alpha \left(E_\alpha^c\right)$. 
By the definition of an intersection, $x$ must belong to $E_\alpha^c$ for every index $\alpha$. 
Consequently, $x$ is not an element of $E_\alpha$ for any $\alpha$.
Since $x$ is in none of the individual sets, it cannot possibly be in their union. Thus, $x \notin \bigcup_\alpha E_\alpha$. 
This directly means $x$ resides in the complement of the union, so $x \in \left(\bigcup_\alpha E_\alpha\right)^c$, which means $x \in A$. This establishes that $B \subset A$. 

\textbf{Step 3: Conclusion.} \\
Since we have rigorously shown both $A \subset B$ and $B \subset A$, it logically follows that the sets are strictly equal, so $A = B$.
\end{proof}
\begin{theorem}
A set $E$ is open if and only if its complement is closed.
\end{theorem}

\begin{proof}
Since this is an "if and only if" statement, we must prove the implication in both directions.

\textbf{Part 1: Assume the complement is closed.} \\
First, let us suppose that the complement $E^c$ is a closed set. Our goal is to prove that $E$ is an open set, which means we must show that every point in $E$ is an interior point.

Choose an arbitrary point $x \in E$. Because $x$ is in $E$, it is clearly not in the complement, so $x \notin E^c$. We are given that $E^c$ is a closed set, which by definition means it contains all of its limit points. Since $x$ is not contained in $E^c$, $x$ cannot possibly be a limit point of $E^c$.

What does it mean for $x$ to \textit{not} be a limit point of $E^c$? It means there must exist at least one neighborhood $N$ of $x$ that does not contain any points of $E^c$. In other words, the intersection $E^c \cap N$ is empty. 

If this neighborhood $N$ contains no points from the complement $E^c$, then every single point inside $N$ must belong to $E$. Therefore, $N \subset E$. Because we found a neighborhood of $x$ completely contained in $E$, $x$ is an interior point of $E$. Since $x$ was chosen arbitrarily, every point in $E$ is an interior point, proving that $E$ is open.

\textbf{Part 2: Assume the set is open.} \\
Next, let us suppose the converse: $E$ is an open set. We must prove that $E^c$ is a closed set. To do this, we need to show that $E^c$ contains all of its limit points.

Let $x$ be an arbitrary limit point of $E^c$. By the fundamental definition of a limit point, \textit{every} neighborhood of $x$ must contain at least one point of $E^c$. 

Because every neighborhood of $x$ contains some point outside of $E$ (since it contains a point of $E^c$), it is impossible to draw a neighborhood around $x$ that is entirely contained within $E$. Consequently, $x$ cannot be an interior point of $E$.

Now we use our initial assumption: $E$ is an open set. This means that \textit{every} element of $E$ is an interior point. Since we just established that $x$ is not an interior point, it is impossible for $x$ to be an element of $E$. Thus, $x \notin E$.

If $x$ is not in $E$, it must belong to its complement, so $x \in E^c$. We have shown that any arbitrary limit point of $E^c$ must belong to $E^c$. Therefore, $E^c$ contains all of its limit points, proving that $E^c$ is closed.
\end{proof}
\begin{corollary}
A set $F$ is closed if and only if its complement $F^c$ is open.
\end{corollary}

\begin{proof}
This follows as an immediate consequence of our previous theorem, utilizing the fundamental fact that the complement of a complement returns the original set (i.e., $(F^c)^c = F$).

\textbf{Part 1: Assume $F$ is closed.} \\
If $F$ is closed, we can think of $F$ as the complement of some set $E$ (specifically, $E = F^c$). Since $F$ is closed, this means $E^c$ is closed. By our previous theorem, a set $E$ is open if and only if $E^c$ is closed. Therefore, $E$ (which is $F^c$) must be open.

\textbf{Part 2: Assume $F^c$ is open.} \\
Conversely, suppose the complement $F^c$ is open. Let us call this open set $E$ (so $E = F^c$). Our previous theorem states that if a set $E$ is open, its complement $E^c$ is closed. Since $E^c = (F^c)^c = F$, it logically follows that $F$ is closed.

Therefore, $F$ is closed if and only if $F^c$ is open.
\end{proof}
\begin{theorem}
The following properties hold for open and closed sets:
\begin{enumerate}[label=(\alph*)]
    \item For any collection $\{G_\alpha\}$ of open sets, the union $\bigcup_\alpha G_\alpha$ is open.
    \item For any collection $\{F_\alpha\}$ of closed sets, the intersection $\bigcap_\alpha F_\alpha$ is closed.
    \item For any finite collection $G_1, \dots, G_n$ of open sets, the intersection $\bigcap_{i=1}^n G_i$ is open.
    \item For any finite collection $F_1, \dots, F_n$ of closed sets, the union $\bigcup_{i=1}^n F_i$ is closed.
\end{enumerate}
\end{theorem}

\begin{proof}
We will prove each part sequentially, using the fundamental definitions of open sets, closed sets, and neighborhoods.

\textbf{Part (a): Arbitrary union of open sets} \\
Let $G = \bigcup_\alpha G_\alpha$. To prove $G$ is open, we must show every point in $G$ is an interior point. 
Suppose $x \in G$. By the definition of a union, $x$ must belong to at least one of the sets in the collection, say $G_\alpha$. Because $G_\alpha$ is given as an open set, $x$ is an interior point of $G_\alpha$. This means there exists a neighborhood $N$ of $x$ such that $N \subset G_\alpha$. 
Since $G_\alpha$ is a subset of the entire union $G$ (i.e., $G_\alpha \subset G$), it directly follows that $N \subset G$. Therefore, $x$ is an interior point of $G$. Since $x$ was arbitrary, $G$ is open.

\textbf{Part (b): Arbitrary intersection of closed sets} \\
We can prove this elegantly using De Morgan's Laws and our previous knowledge about complements. By De Morgan's Law, the complement of an intersection is the union of the complements:
\[ \left(\bigcap_\alpha F_\alpha\right)^c = \bigcup_\alpha (F_\alpha^c) \]
Since each $F_\alpha$ is a closed set, its complement $F_\alpha^c$ must be an open set. The right side of the equation is now an arbitrary union of open sets. By the result we just proved in Part (a), this union must be open. 
Since the complement of $\bigcap_\alpha F_\alpha$ is open, the set $\bigcap_\alpha F_\alpha$ itself must be closed.

\textbf{Part (c): Finite intersection of open sets} \\
Let $H = \bigcap_{i=1}^n G_i$. Let $x \in H$. By the definition of intersection, $x$ belongs to every single $G_i$ for $i = 1, \dots, n$. 
Since each $G_i$ is open, $x$ is an interior point of every $G_i$. Thus, for each $i$, there exists a neighborhood $N_i$ of $x$ with some strictly positive radius $r_i$ such that $N_i \subset G_i$. 
To ensure we find a neighborhood that fits inside \textit{all} $G_i$ simultaneously, we take the smallest of these radii. Let:
\[ r = \min(r_1, \dots, r_n) \]
Crucially, because this is a \textit{finite} collection of strictly positive numbers, their minimum $r$ is guaranteed to be strictly positive ($r > 0$). Let $N$ be the neighborhood of $x$ with radius $r$. Because $r \le r_i$ for every $i$, $N$ is contained in every $N_i$, meaning $N \subset G_i$ for $i = 1, \dots, n$. 
Therefore, $N \subset H$. This shows $x$ is an interior point of $H$, and thus $H$ is open.

\textbf{Part (d): Finite union of closed sets} \\
We apply De Morgan's Laws again, this time for a finite collection. By taking complements, we get:
\[ \left(\bigcup_{i=1}^n F_i\right)^c = \bigcap_{i=1}^n (F_i^c) \]
Because each $F_i$ is closed, its complement $F_i^c$ is open. The right side is a finite intersection of open sets, which we just proved in Part (c) must be open. 
Since the complement of $\bigcup_{i=1}^n F_i$ is open, it follows that the finite union $\bigcup_{i=1}^n F_i$ is closed.
\end{proof}
\begin{example}[Why "Finite" Matters]
In parts (c) and (d) of the previous theorem, we specified that the collections must be finite. This finiteness is absolutely essential, and we can prove why using a simple counterexample. 

Let us define a sequence of open intervals (segments) in the real numbers $\mathbb{R}^1$:
\[ G_n = \left(-\frac{1}{n}, \frac{1}{n}\right) \quad \text{for } n = 1, 2, 3, \dots \]

Each individual set $G_n$ is clearly an open subset of $\mathbb{R}^1$. 
As $n$ gets larger (e.g., $n=1, 2, 10, 100$), the intervals $(-1, 1), (-\frac{1}{2}, \frac{1}{2}), (-\frac{1}{10}, \frac{1}{10})$ shrink tighter and tighter around zero.

Now, let us find the intersection of this entirely infinite collection of sets, which we will call $G$:
\[ G = \bigcap_{n=1}^\infty G_n \]

What points survive this infinite intersection? The only number that belongs to \textit{every single} one of these shrinking intervals is exactly $0$. Therefore, the final set $G$ consists of just a single point: $x = 0$. 

A set consisting of a single point is not an open subset of $\mathbb{R}^1$ (because you cannot draw a neighborhood with a positive radius around $0$ without stepping outside the set $\{0\}$). 

This demonstrates that the intersection of an \textit{infinite} collection of open sets need not be open. By a very similar logic, the union of an \textit{infinite} collection of closed sets need not be closed.
\end{example}
\begin{definition}
If $X$ is a metric space, if $E \subset X$, and if $E'$ denotes the set of all limit points of $E$ in $X$, then the \textit{closure} of $E$ is the set $\overline{E} = E \cup E'$.
\end{definition}

\begin{theorem}
If $X$ is a metric space and $E \subset X$, then:
\begin{enumerate}[label=(\alph*)]
    \item $\overline{E}$ is closed,
    \item $E = \overline{E}$ if and only if $E$ is closed,
    \item $\overline{E} \subset F$ for every closed set $F \subset X$ such that $E \subset F$.
\end{enumerate}
By (a) and (c), $\overline{E}$ is the smallest closed subset of $X$ that contains $E$.
\end{theorem}

\begin{proof}
\textbf{Part (a):} Let $p \in X$ such that $p \notin \overline{E}$. This means $p$ is neither a point of $E$ nor a limit point of $E$. Because $p$ is not a limit point, there exists a neighborhood $N$ of $p$ that does not intersect $E$. Furthermore, $N$ cannot contain any limit points of $E$ (if it did, any neighborhood around that limit point would intersect $E$, which would force $N$ to intersect $E$). Therefore, $N$ does not intersect $E'$ either. Thus, $N \subset \overline{E}^c$, meaning every point in the complement of $\overline{E}$ is an interior point. So $\overline{E}^c$ is open, which means $\overline{E}$ is closed.

\textbf{Part (b):} If $E = \overline{E}$, then $E$ is closed because $\overline{E}$ is closed by part (a). Conversely, if $E$ is closed, it contains all of its limit points by definition. Thus, $E' \subset E$, which implies $\overline{E} = E \cup E' = E$.

\textbf{Part (c):} Let $F$ be a closed set such that $E \subset F$. If $p$ is a limit point of $E$, every neighborhood of $p$ intersects $E$. Since $E \subset F$, every neighborhood of $p$ also intersects $F$, making $p$ a limit point of $F$. Thus, $E' \subset F'$. Because $F$ is closed, it contains all its limit points ($F' \subset F$). Therefore, $E' \subset F$. Since $E \subset F$ and $E' \subset F$, it follows that $\overline{E} = E \cup E' \subset F$.
\end{proof}

\begin{theorem}
Let $E$ be a nonempty set of real numbers which is bounded above. Let $y = \sup E$. Then $y \in \overline{E}$. Hence $y \in E$ if $E$ is closed.
\end{theorem}

\begin{proof}
If $y \in E$, then trivially $y \in \overline{E}$. 

Suppose $y \notin E$. By the definition of the supremum, $y$ is the least upper bound of $E$. For any strictly positive number $h > 0$, the value $y - h$ is strictly less than $y$, and therefore $y - h$ cannot be an upper bound for $E$. This means there must exist some point $x \in E$ such that $y - h < x < y$. 

Since we can find such an $x \in E$ for any $h > 0$, every neighborhood of $y$ contains a point of $E$ distinct from $y$. Therefore, $y$ is a limit point of $E$ (i.e., $y \in E'$). Since $y \in E'$, it follows that $y \in \overline{E}$. If $E$ is a closed set, $E = \overline{E}$, and therefore $y \in E$.
\end{proof}

\begin{example}[Examples of Closure]
\begin{enumerate}
    \item \textbf{Rationals and Irrationals:} Consider the rational numbers $\mathbb{Q}$ and the irrational numbers $\mathbb{R} \setminus \mathbb{Q}$ as subsets of $\mathbb{R}$. Both sets are dense in $\mathbb{R}$. Every real number is a limit point for both sets. Therefore, $\overline{\mathbb{Q}} = \mathbb{R}$ and $\overline{\mathbb{R} \setminus \mathbb{Q}} = \mathbb{R}$.
    
    \item \textbf{Relative Closure:} Consider the metric space $X = (0, 2]$ equipped with the standard metric. Let $E = (0, 1)$ be a subset of $X$. The only limit point of $E$ that actually belongs to the metric space $X$ is the point $1$. The point $0$ is not in $X$, so it cannot be a limit point in this space. Therefore, the closure of $E$ in $X$ is $\overline{E} = (0, 1]$.
    
    \item \textbf{Sets without Limit Points:} Consider the natural numbers $\mathbb{N}$ as a subset of $\mathbb{R}$. $\mathbb{N}$ has no limit points, so its set of limit points is empty ($\mathbb{N}' = \emptyset$). Therefore, $\overline{\mathbb{N}} = \mathbb{N} \cup \emptyset = \mathbb{N}$. The exact same logic applies to any finite set or any set consisting entirely of isolated points: they contain no limit points, so their closure is simply the set itself.
\end{enumerate}
\end{example}
\begin{definition}
Suppose $E \subset Y \subset X$, where $X$ is a metric space. We say that $E$ is \textit{open relative to $Y$} if to each $p \in E$ there is associated an $r > 0$ such that $q \in E$ whenever $d(p, q) < r$ and $q \in Y$. As a consequence, a set may be open relative to $Y$ without being an open subset of the larger space $X$.
\end{definition}

\begin{example}[Relative Openness in Intervals]
Let $X = \mathbb{R}$ with the standard metric, and consider the subspace $Y = (0, 2]$. 

First, consider the subset $E = (1, 2]$. In the parent space $X = \mathbb{R}$, $E$ is not open because no standard open interval around the point $2$ is entirely contained in $E$ (any interval $(2-r, 2+r)$ contains points strictly greater than $2$). However, $E$ is open relative to $Y$. For the point $p = 2 \in E$, we can choose a radius $r = 1$. The relative neighborhood consists of all $q \in Y$ such that $d(2, q) < 1$. These points form the interval $(1, 2]$, which is entirely contained inside $E$. Thus, $(1, 2]$ is an open set in the relative topology of $Y$.

Next, consider $E = (0, 1)$. This set is open as a standard subset of $\mathbb{R}$, and it also remains open relative to $Y$, since restricting its standard neighborhoods to $Y$ keeps them fully contained within $E$.
\end{example}

\begin{example}[Relative Openness of Natural Numbers]
Consider $Y = \mathbb{N}$ as a subspace of the metric space $X = \mathbb{R}$ with the standard metric. Let $E = \{n\}$ be a singleton set where $n \in \mathbb{N}$. In $\mathbb{R}$, a single point is never an open set. 

However, $E$ is open relative to $\mathbb{N}$. If we choose a radius $r = \frac{1}{2}$, the only point $q \in \mathbb{N}$ satisfying $d(n, q) < \frac{1}{2}$ is the point $n$ itself. Thus, the relative neighborhood around $n$ is just $\{n\}$, which is entirely contained in $E$. Because every singleton is open relative to $\mathbb{N}$, any subset of $\mathbb{N}$ is open relative to $\mathbb{N}$.
\end{example}

\begin{remark}[Subspace Topology]
This concept of relative openness formally describes the \textit{subspace topology}. When a subset $Y$ inherits its metric from a parent space $X$, the open sets of $Y$ are exactly the intersections of the open sets of $X$ with $Y$. As demonstrated in the previous example, the natural numbers $\mathbb{N}$ inherit the \textit{discrete topology} from the usual topology of $\mathbb{R}$, which fundamentally means that every subset of $\mathbb{N}$ (including singletons) is an open set in that relative space.
\end{remark}
\begin{theorem}
Suppose $Y \subset X$. A subset $E$ of $Y$ is open relative to $Y$ if and only if $E = Y \cap G$ for some open subset $G$ of $X$.
\end{theorem}

\begin{proof}
Because this is an "if and only if" statement, we divide the proof into two parts to show both directions of the implication.

\textbf{Part 1: Assume $E$ is open relative to $Y$.} \\
Suppose $E$ is open relative to $Y$. Our goal is to construct an open set $G$ in the larger space $X$ such that $E = G \cap Y$.

By the definition of relative openness, for every point $p \in E$, there exists a positive radius $r_p > 0$ such that any point $q$ satisfying $d(p, q) < r_p$ and $q \in Y$ must belong to $E$. 

Let us construct a standard open neighborhood in the entire space $X$ around each point $p \in E$ using this specific radius $r_p$. Let this neighborhood be $V_p$:
\[ V_p = \{ q \in X \mid d(p, q) < r_p \} \]
Next, we define our set $G$ as the union of all these individual neighborhoods for every point in $E$:
\[ G = \bigcup_{p \in E} V_p \]
Because every neighborhood $V_p$ is an open set in $X$, and any arbitrary union of open sets is open, $G$ is guaranteed to be an open subset of $X$.

Now we must verify that $E = G \cap Y$. 
First, since every point $p \in E$ is the center of its own neighborhood $V_p$, it is clear that $p \in G$. Since $E$ is initially given as a subset of $Y$, it follows that $E \subset G \cap Y$.
Second, by our specific choice of the radius $r_p$, the portion of any neighborhood $V_p$ that actually lies in $Y$ is completely contained within $E$. Mathematically, $V_p \cap Y \subset E$ for every $p \in E$. Therefore, the union of all these intersections, which represents $G \cap Y$, must also be a subset of $E$. So, $G \cap Y \subset E$.

Since $E \subset G \cap Y$ and $G \cap Y \subset E$, we conclude that $E = G \cap Y$.

\textbf{Part 2: Assume $E = G \cap Y$.} \\
Conversely, suppose there exists some open set $G$ in $X$ such that $E = G \cap Y$. We must prove that $E$ is open relative to $Y$.

Let $p$ be an arbitrary point in $E$. Because $E = G \cap Y$, this point $p$ must reside in $G$. Since $G$ is an open set in $X$, $p$ is an interior point of $G$. This means there exists a neighborhood $V_p$ around $p$ (with some positive radius) that is entirely contained within $G$, meaning $V_p \subset G$.

If we intersect this neighborhood with $Y$, we get:
\[ V_p \cap Y \subset G \cap Y \]
But we are given that $G \cap Y = E$, which means $V_p \cap Y \subset E$. 

In practical terms, we have found a radius around $p$ such that all points within that radius that also belong to $Y$ are guaranteed to fall in $E$. Since this holds for any arbitrary point $p \in E$, the set $E$ is open relative to $Y$.
\end{proof}
\subsection*{Compact Sets}

\begin{definition}
By an \textit{open cover} of a set $E$ in a metric space $X$ we mean a collection $\{G_\alpha\}$ of open subsets of $X$ such that $E \subset \bigcup_\alpha G_\alpha$.
\end{definition}

\begin{definition}
A subset $K$ of a metric space $X$ is said to be \textit{compact} if every open cover of $K$ contains a \textit{finite subcover}.
\end{definition}

\begin{example}[Compact and Non-Compact Sets]
To understand compactness, it is helpful to look at sets that have this property and sets that fail to have it.

\textbf{1. Finite Sets (Always Compact)} \\
Any finite set is trivially compact. If a set $K = \{p_1, p_2, \dots, p_n\}$ is covered by an arbitrary open cover $\{G_\alpha\}$, we can simply pick one set from the cover for each point. Let $G_1$ cover $p_1$, $G_2$ cover $p_2$, and so on. The collection $\{G_1, G_2, \dots, G_n\}$ is a finite subcover.

\textbf{2. The Natural Numbers (Not Compact)} \\
Consider the natural numbers $\mathbb{N}$ as a subset of $\mathbb{R}$ with the standard metric. Let us construct an open cover by taking an open interval of radius $0.1$ around each natural number:
\[ G_n = (n - 0.1, n + 0.1) \quad \text{for each } n \in \mathbb{N} \]
The collection $\{G_1, G_2, G_3, \dots\}$ is clearly an open cover of $\mathbb{N}$. However, this cover has no finite subcover. Any finite subcollection can only cover a finite number of integers. Because there are infinitely many natural numbers, a finite subcollection will always leave infinitely many points of $\mathbb{N}$ uncovered. Thus, $\mathbb{N}$ is not compact.

\textbf{3. The Set $E = \{1/n \mid n \in \mathbb{N}\}$ (Not Compact)} \\
Consider the set $E = \{1, 1/2, 1/3, \dots\}$. We can cover this set with the collection of open intervals:
\[ U_n = \left(\frac{1}{n}, 2\right) \quad \text{for } n \in \mathbb{N} \]
The union of all $U_n$ covers $E$. However, if we try to take a finite subcover, it will consist of a finite number of these sets. Let $N$ be the largest index among the sets in our finite subcover. The union of this finite subcover will just be $U_N = (1/N, 2)$. This finite subcover completely fails to cover the points $\frac{1}{N+1}, \frac{1}{N+2}, \dots$ of $E$. Since we found an open cover that has no finite subcover, $E$ is not compact.

\textbf{4. The Set $K = \{1/n \mid n \in \mathbb{N}\} \cup \{0\}$ (Compact)} \\
By simply adding the limit point $0$ to our previous set, it becomes compact. Let us prove this rigorously. 

Let $\{V_\alpha\}$ be an arbitrary open cover of $K$. Since $0 \in K$, there must be some open set in our cover, let's call it $V_0$, that contains the point $0$. Because $V_0$ is an open set, there exists a neighborhood around $0$ with some radius $r > 0$ that is entirely contained within $V_0$.

By the Archimedean property, we can find a large enough integer $N$ such that $\frac{1}{N} < r$. This means that for all $n \ge N$, the points $\frac{1}{n}$ are strictly inside that neighborhood, and therefore are all safely covered by the single set $V_0$.

How many points of $K$ are left uncovered? Only finitely many: $1, 1/2, \dots, \frac{1}{N-1}$. Because there are only finitely many points remaining, we can just pick one open set from our original cover $\{V_\alpha\}$ for each of them. Let's call them $V_1, V_2, \dots, V_{N-1}$. 

The finite collection $\{V_0, V_1, V_2, \dots, V_{N-1}\}$ now covers the entire set $K$. Since we successfully found a finite subcover for an arbitrary open cover, $K$ is compact.

\textbf{5. The Set $S = \{\frac{1}{n} - \frac{1}{m}\} \cup \{0\}$} \\
Consider the set formed by the differences of unit fractions, together with $0$. The limit points of this set include sequences where $n$ is fixed and $m \to \infty$ (giving $1/n$), sequences where $m$ is fixed and $n \to \infty$ (giving $-1/m$), and sequences where both go to infinity (giving $0$). 

Because the set $S$ contains all of these limit points, $S$ is a \textit{closed} set. Furthermore, because the maximum distance from $0$ is bounded by $1$, $S$ is also a \textit{bounded} set. Just like the previous example, this set is compact (any open set covering $0$ and the other limit points will cover all but finitely many elements). This dual property—being both closed and bounded—becomes highly relevant very soon when we prove the equivalence of compact sets and closed-and-bounded sets in $\mathbb{R}^k$.
\end{example}
\begin{theorem}
Suppose $K \subset Y \subset X$. Then $K$ is compact relative to $X$ if and only if $K$ is compact relative to $Y$.
\end{theorem}

\begin{proof}
We need to prove this equivalence in both directions.

\textbf{Part 1: Assume $K$ is compact relative to $X$.} \\
Let $\{V_\alpha\}$ be an arbitrary collection of sets, open relative to $Y$, that covers $K$ (so $K \subset \bigcup_\alpha V_\alpha$).
By our previous theorem on subspace topologies, for each set $V_\alpha$ that is open relative to $Y$, there exists a corresponding set $G_\alpha$ that is open in the larger space $X$ such that $V_\alpha = Y \cap G_\alpha$.

Because the sets $V_\alpha$ cover $K$, and each $V_\alpha$ is a subset of $G_\alpha$, the collection of open sets $\{G_\alpha\}$ clearly covers $K$ in $X$.
Since we assumed $K$ is compact relative to $X$, this open cover $\{G_\alpha\}$ must have a finite subcover. Thus, there exist finitely many indices $\alpha_1, \dots, \alpha_n$ such that:
\[ K \subset G_{\alpha_1} \cup \dots \cup G_{\alpha_n} \]

Because $K$ is entirely contained within $Y$, intersecting the right side of the above inclusion with $Y$ will still fully cover $K$. Therefore:
\[ K \subset (Y \cap G_{\alpha_1}) \cup \dots \cup (Y \cap G_{\alpha_n}) \]
Substituting back our definition of $V_\alpha$, we get:
\[ K \subset V_{\alpha_1} \cup \dots \cup V_{\alpha_n} \]
We have successfully found a finite subcover from our original relative open cover $\{V_\alpha\}$. This proves that $K$ is compact relative to $Y$.

\textbf{Part 2: Assume $K$ is compact relative to $Y$.} \\
Conversely, suppose $\{G_\alpha\}$ is a collection of open subsets in $X$ that covers $K$.
For each index $\alpha$, we define $V_\alpha = Y \cap G_\alpha$. Because each $G_\alpha$ is open in $X$, each $V_\alpha$ is open relative to $Y$.

Since $K \subset Y$ and $K$ is covered by the sets $G_\alpha$, it follows that $K$ is covered by their intersections with $Y$. Thus, $\{V_\alpha\}$ is an open cover of $K$ relative to $Y$.
By our assumption that $K$ is compact relative to $Y$, there must exist a finite subcover. Therefore, for some finitely many indices $\alpha_1, \dots, \alpha_n$:
\[ K \subset V_{\alpha_1} \cup \dots \cup V_{\alpha_n} \]

Since $V_\alpha \subset G_\alpha$ for every index, it immediately follows that:
\[ K \subset G_{\alpha_1} \cup \dots \cup G_{\alpha_n} \]
We have found a finite subcover from the original open cover $\{G_\alpha\}$ in $X$. This completes the proof that $K$ is compact relative to $X$.
\end{proof}

\begin{remark}[Absolute Nature of Compactness]
By virtue of this theorem, we can regard compact sets as metric spaces in their own right, without paying attention to the embedding space. Openness and closedness are purely relative properties (they depend heavily on the parent space), but compactness is an absolute property. Proving a subset is compact within a subspace under the inherited metric is mathematically identical to proving it is compact in the whole larger metric space.
\end{remark}

\begin{example}
Consider the metric space $X = \mathbb{R}$ with the standard metric, and the subspace $Y = (0, 2]$.

By our theorem, any subset of $(0, 2]$ that is compact relative to $(0, 2]$ will also be compact in $\mathbb{R}$, and vice versa. However, this stability absolutely does not hold for closed and open sets:
\begin{itemize}
    \item The set $(0, 2]$ is \textbf{closed} in $(0, 2]$ (because every metric space is a closed subset of itself), but it is \textbf{not closed} in $\mathbb{R}$ (because it fails to contain its limit point $0$).
    \item The set $(1, 2]$ is \textbf{open} in $(0, 2]$ (it is the intersection of the standard open interval $(1, 3)$ with $Y$), but it is \textbf{not open} in $\mathbb{R}$ (because no standard open interval centered at $2$ is entirely contained within the set).
\end{itemize}
Compactness transcends this relativity, meaning a compact set remains intrinsically compact regardless of what larger space you place it in.
\end{example}
\begin{theorem}
Compact subsets of metric spaces are closed.
\end{theorem}

\begin{proof}
Let $K$ be a compact subset of a metric space $X$. To prove $K$ is closed, we will establish that its complement, $K^c$, is an open set.

Suppose $p$ is an arbitrary point in the complement $K^c$ (meaning $p \notin K$). Our objective is to find a neighborhood around $p$ that is entirely contained within $K^c$. 

For every point $q \in K$, the distance $d(p, q)$ is strictly positive because $p$ and $q$ are distinct points. Let us construct two non-overlapping neighborhoods for each $q$:
\begin{enumerate}
    \item A neighborhood $V_q$ centered at $p$ with a radius less than $\frac{1}{2} d(p, q)$.
    \item A neighborhood $W_q$ centered at $q$ with a radius less than $\frac{1}{2} d(p, q)$.
\end{enumerate}
Because the radii are less than half the distance between $p$ and $q$, these two neighborhoods are guaranteed to be disjoint.

Now, consider the collection of all these neighborhoods $\{W_q\}$ for every $q \in K$. This collection forms an open cover of the set $K$. Since $K$ is given to be compact, this arbitrary open cover must contain a finite subcover. Therefore, there are finitely many points $q_1, \dots, q_n$ in $K$ such that their corresponding neighborhoods completely cover $K$:
\[ K \subset W_{q_1} \cup \dots \cup W_{q_n} = W \]

For each of these specifically chosen points $q_i$, there is a corresponding neighborhood $V_{q_i}$ centered around our original point $p$. Let $V$ be the intersection of these finitely many neighborhoods around $p$:
\[ V = V_{q_1} \cap \dots \cap V_{q_n} \]
Since it is an intersection of a finite number of open neighborhoods, $V$ itself is an open neighborhood of $p$. 

Because each individual $V_{q_i}$ is completely disjoint from its corresponding $W_{q_i}$, their intersection $V$ is completely disjoint from the entire union $W$. Since $K$ is a subset of $W$, $V$ does not intersect $K$ at all. This means $V \subset K^c$. 

We have successfully found a neighborhood $V$ around an arbitrary point $p \in K^c$ that lies entirely outside of $K$. Therefore, $p$ is an interior point of $K^c$, making $K^c$ an open set. Consequently, $K$ is closed.
\end{proof}

\begin{theorem}
Closed subsets of compact sets are compact.
\end{theorem}

\begin{proof}
Suppose $F \subset K \subset X$, where $F$ is a closed set (relative to the metric space $X$) and $K$ is a compact set. We must prove that $F$ is compact.

Let $\{V_\alpha\}$ be an arbitrary open cover of $F$. To prove compactness, we need to extract a finite subcover from it.

Since $F$ is a closed set, its complement $F^c$ is an open set. Let us add this single open set $F^c$ to our covering collection. We will define a new collection of open sets $\Omega$:
\[ \Omega = \{V_\alpha\} \cup \{F^c\} \]
This new collection $\Omega$ is an open cover for the entire set $K$. To see why, pick any point in $K$. It either belongs to $F$ (in which case it is covered by some $V_\alpha$) or it does not belong to $F$ (in which case it is covered by $F^c$).

Since $K$ is a compact set, every open cover of $K$ must have a finite subcover. So there exists a finite subcollection $\Phi$ chosen from $\Omega$ that completely covers $K$. Because $F \subset K$, this finite subcollection $\Phi$ also completely covers $F$.

This finite subcollection $\Phi$ might happen to include the set $F^c$. If it does, we simply discard $F^c$ from $\Phi$. Because $F^c$ contains absolutely no points of $F$, removing it from our covering collection does not leave any points of $F$ exposed. 

The remaining sets in our finite subcollection belong entirely to the original cover $\{V_\alpha\}$ and still completely cover $F$. We have successfully found a finite subcover for $F$, which proves that $F$ is compact.
\end{proof}

\begin{corollary}
If $F$ is closed and $K$ is compact, then $F \cap K$ is compact.
\end{corollary}

\begin{proof}
By our above theorem, any compact subset of a metric space is closed. Since $K$ is given as a compact set, $K$ must be closed. 
We are also given that $F$ is a closed set. A fundamental property of open sets dictates that the intersection of any collection of closed sets is closed. Therefore, the intersection $F \cap K$ is a closed set.

Furthermore, it is trivially true that $F \cap K \subset K$. We have just established that $F \cap K$ is a closed subset of the compact set $K$. 
By our second theorem (which states that closed subsets of compact sets are compact), it directly follows that $F \cap K$ is compact.
\end{proof}
\begin{theorem}
If $\{K_\alpha\}$ is a collection of compact subsets of a metric space $X$ such that the intersection of every finite subcollection of $\{K_\alpha\}$ is nonempty, then the total intersection $\bigcap K_\alpha$ is nonempty.
\end{theorem}

\begin{proof}
We will prove this by contradiction. 

\textbf{Step 1: Set up the contradiction.} \\
Fix one specific compact set from the collection and call it $K_1$. 
Because each $K_\alpha$ is a compact subset of a metric space, it must be closed. Therefore, its complement $G_\alpha = K_\alpha^c$ is an open set.

Now, assume for the sake of contradiction that the total intersection $\bigcap K_\alpha$ is empty. This means there is absolutely no point that belongs to \textit{every} set $K_\alpha$. In particular, no point in our fixed set $K_1$ belongs to every $K_\alpha$.

\textbf{Step 2: Construct an open cover.} \\
If no point of $K_1$ is in every $K_\alpha$, then every point of $K_1$ is missing from at least one $K_\alpha$. If a point is missing from $K_\alpha$, it must belong to its complement, $G_\alpha$. 
This implies that the collection of open sets $\{G_\alpha\}$ completely covers $K_1$. Thus, $\{G_\alpha\}$ forms an open cover of $K_1$.

\textbf{Step 3: Use compactness to find a finite subcover.} \\
We are given that $K_1$ is a compact set. By the fundamental definition of compactness, this open cover must have a finite subcover. Thus, there exist finitely many indices $\alpha_1, \dots, \alpha_n$ such that:
\[ K_1 \subset G_{\alpha_1} \cup \dots \cup G_{\alpha_n} \]

\textbf{Step 4: Reach a contradiction.} \\
Let us translate this back to our original sets. The union of the complements $G_{\alpha_i}$ is exactly the complement of the intersection of the sets $K_{\alpha_i}$. 
If $K_1$ is entirely contained within the complement of $(K_{\alpha_1} \cap \dots \cap K_{\alpha_n})$, it means $K_1$ shares absolutely no elements with that intersection. 

Mathematically, this means:
\[ K_1 \cap K_{\alpha_1} \cap \dots \cap K_{\alpha_n} = \emptyset \]
However, this is a finite subcollection of our original sets, and our initial hypothesis explicitly stated that \textit{every} finite subcollection has a nonempty intersection. We have reached a direct contradiction. 

Therefore, our assumption must be false, and the total intersection $\bigcap K_\alpha$ cannot be empty.
\end{proof}

\begin{corollary}
If $\{K_n\}$ is a sequence of nonempty compact sets such that $K_n \supset K_{n+1}$ (for $n = 1, 2, 3, \dots$), then $\bigcap_{n=1}^\infty K_n$ is not empty.
\end{corollary}

\begin{proof}
This follows directly as a special case of the preceding theorem. 

Consider any arbitrary finite subcollection of these sets, say $K_{n_1}, K_{n_2}, \dots, K_{n_m}$. Let $N$ be the largest index among them ($N = \max(n_1, \dots, n_m)$). 

Because the sequence of sets is nested ($K_1 \supset K_2 \supset K_3 \dots \supset K_N$), the set $K_N$ with the largest index is the smallest set in the chain and is completely contained within all the earlier sets in this subcollection. 
Therefore, the intersection of this entire finite subcollection is simply equal to $K_N$. 

Since we are given that every individual set in the sequence is nonempty, $K_N$ is nonempty. 
We have just shown that the intersection of \textit{any} finite subcollection is nonempty. By the previous theorem, it immediately follows that the total intersection of all the sets, $\bigcap_{n=1}^\infty K_n$, is not empty.
\end{proof}
\begin{theorem}
If $E$ is an infinite subset of a compact set $K$, then $E$ has a limit point in $K$.
\end{theorem}

\begin{proof}
We will prove this by contradiction. 
Assume that no point of $K$ is a limit point of $E$. 

By the definition of a limit point, if a point $q \in K$ is not a limit point of $E$, there must exist a neighborhood $V_q$ around $q$ that contains no points of $E$ other than possibly $q$ itself. Thus, each neighborhood $V_q$ contains at most one point of $E$ (specifically, it contains exactly $q$ if $q \in E$, and zero points of $E$ if $q \notin E$).

Now, construct a collection of these neighborhoods $\{V_q\}$ for every single point $q \in K$. Because each point $q$ is the center of its own neighborhood, this collection clearly covers all of $K$. Therefore, $\{V_q\}$ is an open cover of $K$.

Since we are given that $K$ is a compact set, this open cover must contain a finite subcover. That means there is a finite number of neighborhoods, say $V_{q_1}, V_{q_2}, \dots, V_{q_n}$, that completely cover $K$.

Because $E$ is a subset of $K$, this finite subcover must also cover $E$. However, we established that each individual neighborhood $V_{q_i}$ contains at most one point of $E$. Therefore, the union of these finitely many neighborhoods can contain, at most, $n$ points of $E$. 

This implies that the set $E$ has only finitely many points. But our initial premise explicitly stated that $E$ is an infinite subset. We have reached a direct contradiction. 

Therefore, our assumption must be false, and $E$ must possess at least one limit point in $K$.
\end{proof}

\begin{theorem}
If $\{I_n\}$ is a sequence of intervals in $\mathbb{R}^1$, such that $I_n \supset I_{n+1}$ for $n = 1, 2, 3, \dots$, then the intersection $\bigcap_{n=1}^\infty I_n$ is not empty.
\end{theorem}

\begin{proof}
Let each interval be denoted by its endpoints: $I_n = [a_n, b_n]$. The nested condition $I_n \supset I_{n+1}$ implies that for every $n$, the next interval is contained within the previous one, which means $a_n \le a_{n+1}$ and $b_{n+1} \le b_n$.

Let $E$ be the set of all the left endpoints: $E = \{a_1, a_2, a_3, \dots\}$. 
This set $E$ is clearly nonempty. Furthermore, because every interval $I_n$ is contained within the very first interval $I_1 = [a_1, b_1]$, every left endpoint $a_n$ must be less than or equal to $b_1$. Thus, $E$ is bounded above by $b_1$.

By the Least Upper Bound Property of the real numbers, any nonempty set that is bounded above has a supremum. Let $x = \sup E$. 

We want to show that this specific point $x$ belongs to every interval $I_m$ (which would prove the total intersection is not empty). To do this, we need to show that $a_m \le x \le b_m$ for any positive integer $m$.

First, because $x$ is the supremum (the least upper bound) of the set of all left endpoints $E$, it must be greater than or equal to every element in $E$. Therefore, it is obvious that $a_m \le x$ for all $m$.

Second, let us take any two positive integers $m$ and $n$. Because the intervals are nested, we can relate their endpoints:
\[ a_n \le a_{m+n} \le b_{m+n} \le b_m \]
This inequality shows that \textit{every} right endpoint $b_m$ is an upper bound for the entire set of left endpoints $E$. 
Since $x$ is the \textit{least} upper bound of $E$, it must be less than or equal to any other upper bound. Therefore, $x \le b_m$ for each $m$.

Since we have shown that $a_m \le x \le b_m$ for every index $m = 1, 2, 3, \dots$, it follows that $x \in I_m$ for every $m$. Thus, $x \in \bigcap_{n=1}^\infty I_n$, proving the intersection is not empty.
\end{proof}
\begin{theorem}
Let $k$ be a positive integer. If $\{I_n\}$ is a sequence of $k$-cells such that $I_n \supset I_{n+1}$ ($n = 1, 2, 3, \dots$), then $\bigcap_{n=1}^\infty I_n$ is not empty.
\end{theorem}

\begin{proof}
A $k$-cell in $\mathbb{R}^k$ can be viewed as the Cartesian product of $k$ closed, one-dimensional intervals. Let us break down each $k$-cell $I_n$ into its constituent coordinate intervals. 

Let $I_n$ consist of all points $\mathbf{x} = (x_1, \dots, x_k)$ such that for each coordinate index $j$ ($1 \le j \le k$), the coordinate is bounded by:
\[ a_{n,j} \le x_j \le b_{n,j} \quad (n = 1, 2, 3, \dots) \]
We can denote the closed interval for the $j$-th coordinate of the $n$-th cell as $I_{n,j} = [a_{n,j}, b_{n,j}]$.

We are given that the $k$-cells are nested, meaning $I_n \supset I_{n+1}$ for every $n$. For a larger $k$-cell to completely contain a smaller one, every single 1-dimensional coordinate interval of the larger cell must contain the corresponding coordinate interval of the smaller cell. Therefore, for each fixed coordinate dimension $j$, we have a nested sequence of 1-dimensional intervals:
\[ I_{n,j} \supset I_{n+1,j} \quad \text{for all } n \]

Now, we can look at each of the $k$ dimensions independently. For any fixed $j$, the sequence $\{I_{n,j}\}$ is a sequence of nested closed intervals in $\mathbb{R}^1$. By the above theorem (regarding nested sequences of intervals in $\mathbb{R}^1$), we know that the intersection of such a sequence is not empty.

Consequently, for each coordinate index $j \in \{1, 2, \dots, k\}$, there must exist at least one real number, let's call it $x_j^*$, that belongs to the intersection of all the intervals in that dimension. Mathematically, this means:
\[ a_{n,j} \le x_j^* \le b_{n,j} \quad \text{for all } n = 1, 2, 3, \dots \]

Now we assemble these $k$ individual coordinates back into a single point in $\mathbb{R}^k$. Let us construct the point $\mathbf{x}^* = (x_1^*, x_2^*, \dots, x_k^*)$. 

Because every single coordinate $x_j^*$ of this point satisfies the bounding inequalities defining the $k$-cell $I_n$, the point $\mathbf{x}^*$ must belong to $I_n$. Furthermore, because these inequalities hold for \textit{all} $n \ge 1$, the point $\mathbf{x}^*$ belongs to \textit{every} $k$-cell in the sequence. 

Thus, $\mathbf{x}^* \in \bigcap_{n=1}^\infty I_n$. We have found a point that exists in the total intersection, completely proving that the intersection is not empty.
\end{proof}
\begin{theorem}
Every $k$-cell is compact.
\end{theorem}

\begin{proof}
We will prove this step-by-step using a method of contradiction and repeated subdivision.

\textbf{Step 1: The Core Idea and Initial Setup} \\
Let $I$ be a $k$-cell. By definition, $I$ consists of all points $\mathbf{x} = (x_1, \dots, x_k)$ such that $a_j \le x_j \le b_j$ for each dimension $j$ ($1 \le j \le k$). 

The maximum possible distance between any two points in this $k$-cell occurs between opposite corners (the main diagonal). We will denote this maximum distance as $\delta$. Using the standard Euclidean distance formula, we calculate this diagonal as:
\[ \delta = \left( \sum_{j=1}^k (b_j - a_j)^2 \right)^{1/2} \]
Thus, for any two points $\mathbf{x}, \mathbf{y} \in I$, the distance between them is bounded: $|\mathbf{x} - \mathbf{y}| \le \delta$.

To proceed by contradiction, assume the contrary: suppose there exists an open cover $\{G_\alpha\}$ of the $k$-cell $I$ that contains \textit{no finite subcover}. 

\textbf{Step 2: Constructing a Sequence of Nested $k$-cells} \\
We will now systematically divide our $k$-cell. For each coordinate interval $[a_j, b_j]$, find its exact midpoint $c_j = (a_j + b_j)/2$. By cutting every one of the $k$ dimensions in half, we divide our original $k$-cell $I$ into $2^k$ smaller, equal-sized $k$-cells whose union is the original cell $I$.

Since we assumed the entire cell $I$ cannot be covered by any finite subcollection of $\{G_\alpha\}$, it is a logical necessity that at least one of these $2^k$ smaller sub-cells also cannot be finitely covered. (If all $2^k$ sub-cells could be finitely covered, we could just combine those finite covers to finitely cover the whole cell $I$, which contradicts our assumption). 

Pick one of these smaller sub-cells that cannot be finitely covered, and call it $I_1$. 
We then repeat this exact subdivision process on $I_1$, slicing it into $2^k$ even smaller cells, at least one of which cannot be finitely covered. Call that one $I_2$. 

\textbf{Step 3: Properties of the Sequence} \\
By continuing this process infinitely, we obtain a sequence of $k$-cells $\{I_n\}$ that possesses three critical properties:
\begin{enumerate}[label=(\alph*)]
    \item \textbf{Nested:} Each cell is completely contained within the previous one: $I \supset I_1 \supset I_2 \supset I_3 \supset \dots$
    \item \textbf{Uncoverable:} For every integer $n$, the cell $I_n$ cannot be covered by any finite subcollection of $\{G_\alpha\}$.
    \item \textbf{Shrinking Diagonal:} If $\mathbf{x} \in I_n$ and $\mathbf{y} \in I_n$, then $|\mathbf{x} - \mathbf{y}| \le 2^{-n}\delta$. 
    \textit{(Why the $2^n$ denominator? Each time we subdivide a cell, we cut the length of every side exactly in half. After $n$ subdivisions, the length of every side is divided by $2^n$. By the Euclidean distance formula, if every side shrinks by a factor of $2^n$, the overall diagonal maximum distance also shrinks exactly by that same factor of $2^n$).}
\end{enumerate}

\textbf{Step 4: Finding the Common Point} \\
By property (a), we have a nested sequence of $k$-cells. By the above theorem (regarding nested sequences of $k$-cells), there must exist at least one point, let us call it $\mathbf{x}^*$, that lies in \textit{every} single cell $I_n$ in the sequence. 

\textbf{Step 5: Using the Openness of the Cover} \\
Because the original collection $\{G_\alpha\}$ covers the entire initial cell $I$, and $\mathbf{x}^* \in I$, this point $\mathbf{x}^*$ must belong to at least one open set in the cover. Let us call this set $G_\alpha$.

Because $G_\alpha$ is an open set, $\mathbf{x}^*$ is an interior point. This means we can draw a neighborhood around $\mathbf{x}^*$ with some strictly positive radius $r > 0$ that is entirely contained within $G_\alpha$. In other words, any point $\mathbf{y}$ satisfying $|\mathbf{y} - \mathbf{x}^*| < r$ is guaranteed to be inside $G_\alpha$.

\textbf{Step 6: Reaching the Contradiction} \\
We want to show that for a large enough $n$, the entire cell $I_n$ is swallowed completely by this neighborhood. 
The maximum size (diagonal) of the cell $I_n$ is $\delta / 2^n$. We need this maximum size to be strictly less than our radius $r$, which means we need $\delta / 2^n < r$, or equivalently, $2^n > \delta / r$.

Because the sequence $2^n$ grows infinitely large as $n$ increases, the Archimedean property of the real numbers guarantees that we can find an integer $n$ large enough such that $2^n$ surpasses the fixed number $\delta/r$. 

Let us pick such a large $n$. For this specific cell $I_n$, its maximum diagonal $\delta / 2^n$ is strictly less than $r$. 
Now, take any arbitrary point $\mathbf{y} \in I_n$. Because both $\mathbf{y}$ and $\mathbf{x}^*$ belong to $I_n$, the distance between them is at most the diagonal of the cell:
\[ |\mathbf{y} - \mathbf{x}^*| \le 2^{-n}\delta < r \]
Because the distance from $\mathbf{y}$ to $\mathbf{x}^*$ is less than $r$, every single point $\mathbf{y}$ in the cell $I_n$ falls perfectly inside our neighborhood, and therefore falls inside the single open set $G_\alpha$. 
This means $I_n \subset G_\alpha$.

We have just shown that the entire cell $I_n$ can be covered by exactly \textit{one} open set from our cover. However, property (b) explicitly stated that $I_n$ could not be covered by \textit{any} finite subcollection of $\{G_\alpha\}$. 
One set is a finite subcollection. This is a direct contradiction. 

Therefore, our initial assumption must be false. Every open cover of a $k$-cell must have a finite subcover, making every $k$-cell compact.
\end{proof}
\begin{theorem}
If a set $E$ in $\mathbb{R}^k$ has one of the following three properties, then it has the other two:
\begin{enumerate}[label=(\alph*)]
    \item $E$ is closed and bounded.
    \item $E$ is compact.
    \item Every infinite subset of $E$ has a limit point in $E$.
\end{enumerate}
\end{theorem}

\begin{proof}
We will prove this by establishing a chain of implications: (a) implies (b), (b) implies (c), and finally (c) implies (a). This will show that all three statements are logically equivalent.

\textbf{Step 1: Prove that (a) implies (b)} \\
Assume that statement (a) holds, meaning $E$ is both closed and bounded.
Because $E$ is bounded, it can be entirely enclosed within some sufficiently large $k$-cell, which we will call $I$. Thus, $E \subset I$.
By the above theorem (regarding $k$-cells), every $k$-cell is compact. Therefore, $I$ is a compact set.
We now know that $E$ is a closed subset that is contained within the compact set $I$. 
By the above theorem (regarding subsets), closed subsets of compact sets are compact. Therefore, it immediately follows that $E$ is compact. This proves (b).

\textbf{Step 2: Prove that (b) implies (c)} \\
Assume that statement (b) holds, meaning $E$ is compact. 
Let $S$ be any arbitrary infinite subset of $E$. 
By the above theorem (regarding infinite subsets of compact sets), if a set is an infinite subset of a compact set, then it must possess a limit point within that compact set. 
Since $S$ is an infinite subset of the compact set $E$, it directly follows from this previous theorem that $S$ has a limit point in $E$. This proves (c).

\textbf{Step 3: Prove that (c) implies (a)} \\
Assume that statement (c) holds, meaning every infinite subset of $E$ has a limit point in $E$. We must prove that $E$ is both bounded and closed. We will do this in two sub-parts using proof by contradiction.

\textit{Sub-part 3.1: Proving $E$ is bounded.} \\
Suppose, for the sake of contradiction, that $E$ is not bounded. 
This means we can find points in $E$ that are arbitrarily far from the origin. For each positive integer $n = 1, 2, 3, \dots$, we can select a point $\mathbf{x}_n \in E$ such that its distance from the origin is strictly greater than $n$ ($|\mathbf{x}_n| > n$).
Let $S$ be the set consisting of all these chosen points $\mathbf{x}_n$. Because the distances keep growing without bound, all these points are distinct, making $S$ an infinite subset of $E$.
However, the set $S$ clearly has no limit point anywhere in $\mathbb{R}^k$. For any proposed limit point, a bounded neighborhood around it could only ever contain at most finitely many points of $S$, because the points of $S$ eventually escape and stay outside any bounded region. 
Because $S$ has no limit point in $\mathbb{R}^k$, it certainly has no limit point in $E$. This is an infinite subset of $E$ with no limit point in $E$, which directly contradicts our assumption (c). 
Therefore, $E$ must be bounded.

\textit{Sub-part 3.2: Proving $E$ is closed.} \\
Suppose, for the sake of contradiction, that $E$ is not closed. 
By definition, this means there exists some point $\mathbf{x}_0 \in \mathbb{R}^k$ which is a limit point of $E$, but $\mathbf{x}_0$ is not an element of $E$ ($\mathbf{x}_0 \notin E$).
Because $\mathbf{x}_0$ is a limit point of $E$, every single neighborhood of $\mathbf{x}_0$ must contain infinitely many points of $E$. 
We will now construct a specific sequence of points. For each integer $n = 1, 2, 3, \dots$, consider a neighborhood around $\mathbf{x}_0$ of radius $1/n$. Because this is a neighborhood of a limit point, we can pick a point $\mathbf{x}_n \in E$ such that the distance between them is strictly less than this radius:
\[ |\mathbf{x}_n - \mathbf{x}_0| < \frac{1}{n} \]
Let $S$ be the set of these specifically chosen points $\mathbf{x}_n$. 
The set $S$ must be infinite. (If $S$ were finite, the distances $|\mathbf{x}_n - \mathbf{x}_0|$ would reach some minimum positive value and could not keep getting smaller than $1/n$ for arbitrarily large $n$).
By construction, the points of $S$ get arbitrarily close to $\mathbf{x}_0$, making $\mathbf{x}_0$ a limit point of $S$. However, we already know that $\mathbf{x}_0 \notin E$.

To complete the contradiction with (c), we must show that $S$ has \textit{no other} limit point in $E$. We will prove it has no other limit point in the entirety of $\mathbb{R}^k$.
Let $\mathbf{y}$ be any arbitrary point in $\mathbb{R}^k$ such that $\mathbf{y} \neq \mathbf{x}_0$. 
Let the distance between these two distinct points be $d = |\mathbf{x}_0 - \mathbf{y}|$. Since they are distinct, $d > 0$.
We will choose a radius exactly equal to half the distance between $\mathbf{x}_0$ and $\mathbf{y}$: let $r = d/2$. We want to examine the neighborhood around $\mathbf{y}$ with this radius $r$.
By the Reverse Triangle Inequality, the distance from our arbitrary point $\mathbf{y}$ to any sequence point $\mathbf{x}_n$ is bounded below:
\[ |\mathbf{x}_n - \mathbf{y}| \ge |\mathbf{x}_0 - \mathbf{y}| - |\mathbf{x}_n - \mathbf{x}_0| \]
We know $|\mathbf{x}_0 - \mathbf{y}| = d$. We also know that for all points in our sequence, $|\mathbf{x}_n - \mathbf{x}_0| < 1/n$. Therefore:
\[ |\mathbf{x}_n - \mathbf{y}| > d - \frac{1}{n} \]
Because $1/n$ shrinks toward zero as $n$ increases, there will be some stage in the sequence (some sufficiently large integer $N$) after which $1/n \le d/2$. 
For all terms in the sequence after this stage ($n \ge N$), we have:
\[ |\mathbf{x}_n - \mathbf{y}| > d - \frac{d}{2} = \frac{d}{2} = r \]
This strictly shows that after a certain stage $N$, every single point $\mathbf{x}_n$ in our sequence is at a distance strictly greater than $r$ from $\mathbf{y}$. 
Consequently, the neighborhood of $\mathbf{y}$ with radius $r = d/2$ contains absolutely no elements of the sequence $S$ after stage $N$. Thus, this neighborhood contains at most finitely many points of $S$ (specifically, only those from before stage $N$).
By the above theorem (regarding limit points), every neighborhood of a true limit point \textit{must} contain infinitely many points of the set. Since we found a neighborhood of $\mathbf{y}$ that contains only finitely many points of $S$, the point $\mathbf{y}$ cannot be a limit point of $S$.

Because $\mathbf{y}$ was chosen as any arbitrary point other than $\mathbf{x}_0$, we conclude that $S$ has no limit points other than $\mathbf{x}_0$. 
Since $\mathbf{x}_0 \notin E$, the infinite subset $S$ has no limit point in $E$. 
This is a direct contradiction of our assumption (c). Therefore, our premise was false, and $E$ must be closed. 
Since we have shown (c) implies $E$ is bounded and (c) implies $E$ is closed, we have shown (c) implies (a). 
The proof of equivalence is complete.
\end{proof}

\begin{remark}
The assumption that the set $E$ lies in $\mathbb{R}^k$ is crucial for the theorem above. In a general metric space, it is not always true that a closed and bounded set is compact.

\begin{itemize}
    \item \textbf{Compact implies closed and bounded, in every metric space.} Let $K$ be a compact subset of a metric space $X$. By Theorem 1.63, $K$ is closed. $K$ is also bounded: if $K$ is empty there is nothing to prove; otherwise fix a point $q \in X$. The neighborhoods $N_n(q)$, $n = 1, 2, 3, \dots$, form an open cover of $X$ (every $p \in X$ satisfies $d(p,q) < n$ for some integer $n$, by the Archimedean property), hence of $K$. By compactness, finitely many of them, say $N_{n_1}(q), \dots, N_{n_j}(q)$, cover $K$. These neighborhoods are nested, so $K \subset N_m(q)$ with $m = \max(n_1, \dots, n_j)$; that is, $d(p,q) < m$ for all $p \in K$, and $K$ is bounded.

    \textbf{The converse fails in general.} A closed and bounded set need not be compact. Besides the example $E = (0,1]$ in $X = (0,2]$ given below, consider any infinite set $X$ with the discrete metric ($d(x,y) = 1$ for $x \neq y$). $X$ is closed in itself and bounded ($d(x,y) < 2$ for all $x, y$), but it is not compact: the neighborhoods $N_{1/2}(x) = \{x\}$, $x \in X$, form an open cover of $X$, and since each contains only one point, no finite subcollection covers the infinite set $X$.

    \textbf{In $\mathbb{R}^k$ the two notions coincide.} By Theorem 1.70, a subset of $\mathbb{R}^k$ is compact if and only if it is closed and bounded.

    \item \textbf{Counterexample in a general metric space:} 
    Consider the metric space $X = (0, 2]$ equipped with the standard absolute value metric. Let $E = (0, 1]$. 
    Clearly, $E$ is bounded. Furthermore, $E$ is closed \textit{relative to} $X$, because its complement in $X$ is $(1, 2]$, which is an open set in $X$. 
    However, $E$ is not compact. We know that compactness is an absolute, intrinsic property that does not change depending on the surrounding metric space. If we view $E = (0, 1]$ as a subset of $\mathbb{R}$, it is not closed (it is missing the limit point $0$) and thus cannot be compact. Because it is not compact in $\mathbb{R}$, it cannot be compact in $X$. This demonstrates that being "closed and bounded" is not sufficient to guarantee compactness in arbitrary metric spaces.
    
    \item \textbf{The general metric space equivalent:} 
    For general metric spaces, the appropriate analog to this theorem involves completeness. A subset of a general metric space is compact if and only if it is \textit{complete} (meaning every Cauchy sequence converges) and \textit{totally bounded}.
    
    \item \textbf{Compact implies closed (Topological perspective):} 
    While "closed and bounded implies compact" fails in general metric spaces, the reverse implication—that \textit{every compact set is closed}—is true in every metric space. 
    From a broader topological perspective, this is because every metric space is a \textbf{Hausdorff space}. If you closely observe the proof that compact sets are closed, it fundamentally relies on the ability to take two distinct points and separate them using disjoint neighborhoods (choosing a radius equal to half the distance between them). This separation property is exactly what defines a Hausdorff space. 
    Therefore, the rule "compact implies closed" is not universally true for all abstract topological spaces, but it is always true in Hausdorff spaces (and consequently, all metric spaces).
\end{itemize}
\end{remark}
\begin{theorem}[Weierstrass]
Every bounded infinite subset of $\mathbb{R}^k$ has a limit point in $\mathbb{R}^k$.
\end{theorem}

\begin{proof}
Let $E$ be any bounded and infinite subset of $\mathbb{R}^k$. 

Because the set $E$ is bounded, its spatial extent is limited in every direction. This guarantees that we can construct a sufficiently large $k$-cell, which we will denote as $I$, that completely encloses the entirety of $E$. Therefore, we have $E \subset I \subset \mathbb{R}^k$.

By the above theorem (regarding the compactness of $k$-cells), we established that every $k$-cell is a compact set. Consequently, the enclosing $k$-cell $I$ is compact.

At this stage, we have framed $E$ as an infinite subset that is contained entirely within the compact set $I$. 

By the above theorem (regarding infinite subsets of compact sets), we know that if a set is an infinite subset of a compact set, it must possess a limit point that belongs to that compact set. 

Applying this directly, because $E$ is an infinite subset of the compact $k$-cell $I$, the set $E$ must have a limit point that lies within $I$. Since $I$ is just a region within the broader space $\mathbb{R}^k$, this limit point inherently exists in $\mathbb{R}^k$. 

Thus, every bounded infinite subset of $\mathbb{R}^k$ has a limit point in $\mathbb{R}^k$.
\end{proof}
\section*{Perfect Sets}

\begin{theorem}
Let $P$ be a nonempty perfect set in $\mathbb{R}^k$. Then $P$ is uncountable.
\end{theorem}

\begin{proof}
By definition, a set $P$ is \textit{perfect} if it is closed and every point of $P$ is a limit point of $P$. 

Because $P$ is nonempty and every point in $P$ is a limit point, $P$ must contain infinitely many points. We will prove that $P$ is uncountable by contradiction.

\textbf{Step 1: Set up the proof by contradiction.} \\
Suppose $P$ is countable. Then we can arrange all the elements of $P$ into a sequence and enumerate them as:
\[ P = \{\mathbf{x}_1, \mathbf{x}_2, \mathbf{x}_3, \dots\} \]
We will now construct a nested sequence of neighborhoods $\{V_n\}$ to derive a contradiction.

\textbf{Step 2: Construct the sequence of neighborhoods $\{V_n\}$.} \\
Let $V_1$ be any open neighborhood around the first point $\mathbf{x}_1$. The closure of $V_1$, denoted by $\overline{V}_1$, is closed and bounded.

Now, suppose we have constructed an open neighborhood $V_n$ such that $V_n \cap P$ is nonempty. We wish to construct the next neighborhood $V_{n+1}$.
\begin{itemize}
    \item Since $V_n \cap P$ is nonempty, pick a point $\mathbf{y} \in V_n \cap P$.
    \item Because every point of $P$ is a limit point of $P$, the open neighborhood $V_n$ must contain infinitely many points of $P$.
    \item Therefore, we can choose a point $\mathbf{y} \in V_n \cap P$ such that $\mathbf{y} \neq \mathbf{x}_n$.
    \item Since $V_n$ is an open set and $\mathbf{y} \neq \mathbf{x}_n$, we can draw a sufficiently small neighborhood $V_{n+1}$ around $\mathbf{y}$ such that its closure $\overline{V}_{n+1}$ is small enough to satisfy three conditions:
    \begin{enumerate}[label=(\roman*)]
        \item $\overline{V}_{n+1} \subset V_n$
        \item $\mathbf{x}_n \notin \overline{V}_{n+1}$
        \item $V_{n+1} \cap P \neq \emptyset \quad (\text{since } \mathbf{y} \in V_{n+1} \cap P)$
    \end{enumerate}
\end{itemize}
By condition (iii), $V_{n+1}$ satisfies our inductive hypothesis, and we can continue this process indefinitely for all $n = 1, 2, 3, \dots$.

\textbf{Step 3: Define a sequence of compact sets $K_n$.} \\
For each $n \ge 1$, define the set:
\[ K_n = \overline{V}_n \cap P \]
Let us analyze the properties of $\{K_n\}$:
\begin{enumerate}
    \item \textbf{Compactness:} The closure $\overline{V}_n$ is closed and bounded in $\mathbb{R}^k$, so it is compact. Since $P$ is a perfect set, $P$ is closed. Thus, $K_n$ is the intersection of a compact set and a closed set, making $K_n$ a closed subset of a compact set. By the above theorem (regarding closed subsets of compact sets), each $K_n$ is compact.
    \item \textbf{Nonempty:} By condition (iii) above, $V_n \cap P$ is nonempty. Since $V_n \subset \overline{V}_n$, it follows that $K_n = \overline{V}_n \cap P$ is nonempty for every $n$.
    \item \textbf{Nested:} By condition (i), $\overline{V}_{n+1} \subset V_n \subset \overline{V}_n$, which implies $K_{n+1} \subset K_n$ for all $n$.
\end{enumerate}

\textbf{Step 4: Reach a contradiction.} \\
We have built a sequence $\{K_n\}$ of nonempty compact sets such that $K_1 \supset K_2 \supset K_3 \dots$. 

By the above theorem (regarding nested sequences of nonempty compact sets), the total intersection of this sequence must be nonempty:
\[ \bigcap_{n=1}^\infty K_n \neq \emptyset \]

On the other hand, consider the points of $P$. By condition (ii) of our construction, $\mathbf{x}_n \notin \overline{V}_{n+1}$. Consequently, $\mathbf{x}_n \notin K_{n+1}$, which means $\mathbf{x}_n$ cannot belong to the total intersection $\bigcap_{i=1}^\infty K_i$.

Since every point $\mathbf{x}_n \in P$ is excluded from the intersection, **no point of $P$ lies in $\bigcap_{n=1}^\infty K_n$**. 

However, by definition, $K_n = \overline{V}_n \cap P \subset P$ for every $n$. Therefore, the total intersection $\bigcap_{n=1}^\infty K_n$ must be a subset of $P$. If no point of $P$ lies inside this intersection, then the intersection must be empty:
\[ \bigcap_{n=1}^\infty K_n = \emptyset \]

This directly contradicts our earlier result that the intersection is nonempty. 

Thus, our initial assumption that $P$ is countable must be false. We conclude that $P$ is uncountable.
\end{proof}

\begin{remark}
The assumption that the space is $\mathbb{R}^k$ (or more generally, a complete metric space) is vital here. In an arbitrary metric space, a nonempty perfect set is not guaranteed to be uncountable.
\begin{itemize}
    \item \textbf{Counterexample in $\mathbb{Q}$:} 
    Consider the set of rational numbers $\mathbb{Q}$ as a metric space in its own right, equipped with the standard metric $d(x, y) = |x - y|$. 
    Now consider the subset $P = \mathbb{Q}$ within this metric space $\mathbb{Q}$.
    \begin{itemize}
        \item $P$ is closed in $\mathbb{Q}$ (since any space is closed relative to itself).
        \item Every point of $\mathbb{Q}$ is a limit point of $\mathbb{Q}$, because for any rational number, we can construct a sequence of distinct rational numbers that converges to it.
    \end{itemize}
    Therefore, $P = \mathbb{Q}$ is a perfect set in the metric space $(\mathbb{Q}, d)$. However, $\mathbb{Q}$ is countable! 
    \textit{(The proof fails here because closed and bounded neighborhood closures in $\mathbb{Q}$ are not compact, so the nested compact sets theorem cannot be applied.)}
    \item Note that when $\mathbb{Q}$ is viewed as a subset of $\mathbb{R}$ with the standard metric, it is not closed (since its limit points include all irrational numbers), and hence it is not perfect in $\mathbb{R}$.
\end{itemize}
\end{remark}

\begin{corollary}
Every interval $[a, b]$ with $a < b$ is uncountable. In particular, the set of all real numbers $\mathbb{R}$ is uncountable.
\end{corollary}

\begin{proof}
Let $I = [a, b]$ with $a < b$ be a closed interval in $\mathbb{R}^1$.

\textbf{Step 1: Show $[a, b]$ is a perfect set.}
\begin{enumerate}
    \item The interval $[a, b]$ is closed in $\mathbb{R}^1$.
    \item Take any point $x \in [a, b]$ and any open neighborhood $(x - \epsilon, x + \epsilon)$ around $x$ for $\epsilon > 0$. The intersection $(x - \epsilon, x + \epsilon) \cap [a, b]$ always contains infinitely many distinct real numbers. Thus, every point of $[a, b]$ is a limit point of $[a, b]$.
\end{enumerate}
Since $[a, b]$ is closed and every point in $[a, b]$ is a limit point, $[a, b]$ is a nonempty perfect set in $\mathbb{R}^1$.

\textbf{Step 2: Apply the main theorem.} \\
By the above theorem (regarding perfect sets in $\mathbb{R}^k$), every nonempty perfect set in $\mathbb{R}^k$ is uncountable. 
Therefore, the interval $[a, b]$ is uncountable.

\textbf{Step 3: Show $\mathbb{R}$ is uncountable.} \\
Since the real line $\mathbb{R}$ contains the interval $[a, b]$ as a subset, and $[a, b]$ is uncountable, the set of all real numbers $\mathbb{R}$ must also be uncountable.
\end{proof}
\textbf{The Cantor Set}

The set constructed here shows that there exist nonempty, uncountable perfect sets in $\mathbb{R}^1$ that contain no open segments (intervals).

\textbf{Step 1: Construction of the Cantor Set}
\begin{enumerate}
    \item Let $E_0$ be the closed unit interval $[0, 1]$.
    \item Remove the open middle third interval $\left(\frac{1}{3}, \frac{2}{3}\right)$, leaving the set $E_1$, which is the union of two closed intervals:
    \[ E_1 = \left[0, \frac{1}{3}\right] \cup \left[\frac{2}{3}, 1\right] \]
    \item Remove the open middle third of each of those remaining two intervals. This gives $E_2$, which is the union of $2^2 = 4$ closed intervals:
    \[ E_2 = \left[0, \frac{1}{9}\right] \cup \left[\frac{2}{9}, \frac{1}{3}\right] \cup \left[\frac{2}{3}, \frac{7}{9}\right] \cup \left[\frac{8}{9}, 1\right] \]
    \item Continue this process inductively. At the $n$-th step, we obtain a set $E_n$ with the following properties:
    \begin{enumerate}[label=(\alph*)]
        \item $E_0 \supset E_1 \supset E_2 \supset E_3 \supset \dots$
        \item $E_n$ is the union of $2^n$ disjoint closed intervals, each having length $3^{-n}$.
    \end{enumerate}
\end{enumerate}

The \textbf{Cantor set} $P$ is defined as the infinite intersection of all these sets:
\[ P = \bigcap_{n=1}^\infty E_n \]
\\
\textbf{Step 2: Non-Emptiness and Compactness}
\begin{enumerate}
    \item \textbf{Non-Empty:} Each $E_n$ is a non-empty closed set. Because $E_n \supset E_{n+1}$ for all $n$, $\{E_n\}$ is a nested sequence of non-empty compact sets. By the above theorem (regarding nested compact sets), their intersection $P = \bigcap_{n=1}^\infty E_n$ is non-empty.
    \item \textbf{Compact:} Each set $E_n$ is closed and bounded in $\mathbb{R}^1$. Since $P$ is the intersection of closed sets, $P$ is closed. Furthermore, $P \subset E_0 = [0, 1]$, so $P$ is bounded. Because $P$ is a closed and bounded subset of $\mathbb{R}^1$, $P$ is compact.
\end{enumerate}

---

\textbf{Step 3: The Cantor Set Contains No Open Interval}

By construction, no segment (open interval) of the form
\[ \left(\frac{3k+1}{3^m}, \frac{3k+2}{3^m}\right) \]
(where $k$ and $m$ are positive integers) can contain any point of $P$, because all such middle-third segments were removed during step $m$.

Now, take any arbitrary open segment $(\alpha, \beta)$ with $\beta > \alpha$. 
Since the lengths $3^{-m}$ shrink to zero as $m \to \infty$, we can choose an integer $m$ large enough so that:
\[ 3^{-m} < \frac{\beta - \alpha}{6} \]
Because the removed middle-third segments of length $3^{-m}$ are spaced at distances no larger than $6 \cdot 3^{-m}$ apart, the interval $(\alpha, \beta)$ is long enough that it must contain at least one complete removed segment of the form $\left(\frac{3k+1}{3^m}, \frac{3k+2}{3^m}\right)$. 

Since that removed segment contains no points of $P$, the larger segment $(\alpha, \beta)$ cannot be entirely contained in $P$. Thus, \textbf{$P$ contains no open interval}.

\textbf{Step 4: Every Point of $P$ is a Limit Point}

To show that $P$ is perfect, it remains to prove that $P$ contains no isolated points.

\begin{enumerate}
    \item Take any arbitrary point $x \in P$, and let $S = (x - \epsilon, x + \epsilon)$ be any open segment containing $x$ for $\epsilon > 0$.
    \item For each stage $n$, $x \in P \subset E_n$. Let $I_n$ be the specific closed interval of $E_n$ that contains $x$.
    \item The length of $I_n$ is $3^{-n}$. Choose $n$ large enough so that $3^{-n} < \epsilon$. This guarantees that $I_n \subset S$.
    \item The interval $I_n$ has two endpoints. Let $x_n$ be an endpoint of $I_n$ such that $x_n \neq x$.
    \item By the construction of the Cantor set, the endpoints of all sub-intervals at every stage $n$ are never removed, so $x_n \in P$.
    \item Thus, in any neighborhood $S$ of $x$, we have found a point $x_n \in P$ with $x_n \neq x$.
\end{enumerate}

Therefore, \textbf{every point of $P$ is a limit point of $P$}.

\textbf{Step 5: Conclusion - $P$ is Perfect and Uncountable}

Since $P$ is closed and every point of $P$ is a limit point of $P$, \textbf{$P$ is a perfect set}. 
By Theorem 1.73 (which states that every non-empty perfect set in $\mathbb{R}^k$ is uncountable), it immediately follows that \textbf{the Cantor set $P$ is uncountable}.

\textbf{Remark: Total Length and Lebesgue Outer Measure}

Although $P$ is uncountable, its total "length" is zero:
\begin{itemize}
    \item At stage 1, we remove 1 interval of length $\frac{1}{3}$.
    \item At stage 2, we remove 2 intervals of length $\frac{1}{9}$.
    \item At stage $n$, we remove $2^{n-1}$ intervals of length $\frac{1}{3^n}$.
\end{itemize}

Summing the lengths of all removed intervals yields a geometric series:
\[ \sum_{n=1}^\infty 2^{n-1} \cdot \frac{1}{3^n} = \frac{1}{3} \sum_{n=0}^\infty \left(\frac{2}{3}\right)^n = \frac{1}{3} \left( \frac{1}{1 - \frac{2}{3}} \right) = \frac{1}{3} \cdot 3 = 1 \]

Since the original interval $[0, 1]$ has a total length of $1$, and we removed open intervals whose total lengths sum to $1$, the remaining set $P$ has total length zero. In Measure Theory, this is expressed rigorously by stating that \textbf{the Cantor set $P$ has Lebesgue outer measure zero} ($\mu(P) = 0$).
\subsection*{Connected Sets}

\textbf{Definition (Separated and Connected Sets).} \\
Two subsets $A$ and $B$ of a metric space $X$ are said to be \textbf{separated} if both $A \cap \overline{B}$ and $\overline{A} \cap B$ are empty, i.e., if no point of $A$ lies in the closure of $B$ and no point of $B$ lies in the closure of $A$.

A set $E \subset X$ is said to be \textbf{connected} if $E$ is not the union of two nonempty separated sets.


\textbf{Remark (Separated versus Disjoint).} \\
Separated sets are naturally disjoint, but disjoint sets are not necessarily separated:
\begin{itemize}
    \item \textbf{Disjoint but NOT separated:} Consider the closed interval $[0, 1]$ and the open segment $(1, 2)$. Although they do not share any points, they are not separated because $1 \in [0, 1]$ and $1$ is a limit point of $(1, 2)$ (meaning $1 \in \overline{(1, 2)}$). Thus, $[0, 1] \cap \overline{(1, 2)} \neq \emptyset$.
    \item \textbf{Separated sets:} Consider the open segments $(0, 1)$ and $(1, 2)$. The closure of $(0, 1)$ is $[0, 1]$ and the closure of $(1, 2)$ is $[1, 2]$. Neither set contains any point from the closure of the other:
    \[ (0, 1) \cap [1, 2] = \emptyset \quad \text{and} \quad [0, 1] \cap (1, 2) = \emptyset \]
\end{itemize}

The connected subsets of the real line $\mathbb{R}^1$ have a particularly simple characterization, which we will establish in the following theorem.
\begin{theorem}
A subset $E$ of the real line $\mathbb{R}^1$ is connected if and only if it has the following property: 
If $x \in E$, $y \in E$, and $x < z < y$, then $z \in E$.
\end{theorem}

\begin{proof}
We prove both directions of the logical equivalence step-by-step.

\textbf{Step 1: Forward Direction ($\implies$)} \\
We assume $E$ is connected. We must show that for any $x, y \in E$ with $x < z < y$, the intermediate point $z$ must also belong to $E$.

Suppose, for the sake of contradiction, that there exist points $x \in E$ and $y \in E$ with $x < y$, and some point $z \in (x, y)$ such that $z \notin E$.

We use this point $z$ to decompose $E$ into two sets:
\[ A_z = E \cap (-\infty, z) \quad \text{and} \quad B_z = E \cap (z, \infty) \]

Let us examine the properties of $A_z$ and $B_z$:
\begin{enumerate}[label=(\roman*)]
    \item \textbf{Non-empty:} Since $x \in E$ and $x < z$, we have $x \in A_z$, so $A_z \neq \emptyset$. Similarly, since $y \in E$ and $y > z$, we have $y \in B_z$, so $B_z \neq \emptyset$.
    \item \textbf{Union is $E$:} Since $z \notin E$, every point in $E$ is either strictly less than $z$ or strictly greater than $z$. Thus, $E = A_z \cup B_z$.
    \item \textbf{Separated:} Since $A_z \subset (-\infty, z)$, its closure satisfies $\overline{A}_z \subset (-\infty, z]$. Therefore, $\overline{A}_z \cap B_z \subset (-\infty, z] \cap (z, \infty) = \emptyset$. \\
    Similarly, since $B_z \subset (z, \infty)$, its closure satisfies $\overline{B}_z \subset [z, \infty)$, which implies $A_z \cap \overline{B}_z \subset (-\infty, z) \cap [z, \infty) = \emptyset$.
\end{enumerate}

Thus, $E$ is the union of two non-empty separated sets $A_z$ and $B_z$, which means $E$ is not connected. This directly contradicts our initial assumption that $E$ is connected.

Therefore, if $E$ is connected, then for all $x, y \in E$ and $x < z < y$, we must have $z \in E$.

\vspace{0.5em}

\textbf{Step 2: Backward Direction ($\impliedby$)} \\
Now we assume $E$ possesses the property that whenever $x, y \in E$ and $x < z < y$, then $z \in E$. We must prove that $E$ is connected.

Suppose, for the sake of contradiction, that $E$ is not connected. Then $E$ can be written as the union of two non-empty separated sets:
\[ E = A \cup B, \quad \text{where } A \neq \emptyset, \, B \neq \emptyset, \, A \cap \overline{B} = \emptyset, \, \text{and } \overline{A} \cap B = \emptyset \]

Pick points $x \in A$ and $y \in B$. Without loss of generality, assume $x < y$. Define the set:
\[ A \cap [x, y] \]
This set is non-empty (since $x \in A \cap [x, y]$) and bounded above by $y$. Therefore, by the Least Upper Bound Property of $\mathbb{R}$, its supremum exists. Let:
\[ z = \sup(A \cap [x, y]) \]

Since $x \le z \le y$, $z$ is a limit point or an element of $A \cap [x, y]$. By the above theorem (regarding properties of closures), $z \in \overline{A}$.

Because $A$ and $B$ are separated, $\overline{A} \cap B = \emptyset$. Since $z \in \overline{A}$, it follows that $z \notin B$. In particular, since $y \in B$, we know $z \neq y$, so $x \le z < y$.

Now we analyze two possible cases for $z$:

\begin{itemize}
    \item \textbf{Case 1: $z \notin A$.} \\
    Since $z \notin A$ and $z \notin B$, $z$ does not belong to $A \cup B = E$. 
    Thus, $x \in E$, $y \in E$, and $x < z < y$, but $z \notin E$. This contradicts our hypothesis.

    \item \textbf{Case 2: $z \in A$.} \\
    Since $A \cap \overline{B} = \emptyset$ and $z \in A$, we must have $z \notin \overline{B}$. 
    Because $z \notin \overline{B}$, $z$ is not in $B$ and is not a limit point of $B$. Therefore, there exists a neighborhood around $z$ that contains no points of $B$. Specifically, we can find a point $z_1$ such that $z < z_1 < y$ and the interval $(z, z_1]$ contains no points of $B$:
    \[ (z, z_1] \cap B = \emptyset \]
    Furthermore, since $z$ is the supremum of $A \cap [x, y]$ and $z_1 > z$, there can be no points of $A$ in $(z, z_1]$ either, so $(z, z_1] \cap A = \emptyset$.
    
    Combining these two facts, $z_1 \notin A$ and $z_1 \notin B$, which implies $z_1 \notin A \cup B = E$. 
    Thus, $x \in E$, $y \in E$, and $x < z_1 < y$, but $z_1 \notin E$. This again contradicts our hypothesis.
\end{itemize}

Since both cases yield a contradiction, our assumption that $E$ is not connected must be false. Hence, $E$ is connected.
\end{proof}

\begin{remark}

\begin{enumerate}
    \item \textbf{Connected Subsets of the Real Line are Intervals:} \\
    The property "$x \in E, y \in E, x < z < y \implies z \in E$" is precisely the defining characteristic of an interval in $\mathbb{R}^1$. Therefore, the connected subsets of $\mathbb{R}^1$ are exactly intervals of every kind: open intervals $(a, b)$, closed intervals $[a, b]$, half-open intervals $[a, b)$ or $(a, b]$, unbounded rays $(-\infty, a)$, $[a, \infty)$, single points, and the entire line $\mathbb{R}^1$.

    \item \textbf{Connectedness vs. Path-Connectedness (Topological Intuition):} \\
    In general metric and topological spaces, connectedness mathematically means that a set \textit{cannot be split into separated pieces}. 
    
    In higher topology, you will encounter a related concept called \textbf{path-connectedness}, where a space is path-connected if any two points can be joined by a continuous path/curve. While every path-connected set is connected, the converse is not true in general topological spaces—there exist sets that are connected but not path-connected (such as the Topologist's Sine Curve). 
    
    This distinction highlights that topological connectedness is fundamentally about the \textit{absence of separation or gaps}, rather than the existence of geometric paths.
\end{enumerate}
\end{remark}

\clearpage
\section{Continuity}

Throughout this section, $X$, $Y$ and $Z$ denote metric spaces. When we need to distinguish the distances in different spaces, we write $d_X$, $d_Y$ and $d_Z$. When a space is the real line, the complex plane or a euclidean space $\mathbb{R}^k$, the distance is, of course, replaced by absolute values or by norms of differences: $d(x,y) = |x-y|$.

\subsection{Limits of Functions}

Before defining limits of functions, we recall the notion of a convergent sequence in a metric space, since we shall use sequences freely in what follows.

\begin{definition}\label{def:seq}
A sequence $\{p_n\}$ in a metric space $X$ is said to \textit{converge} if there is a point $p \in X$ with the following property: for every $\varepsilon > 0$ there is a positive integer $N$ such that $n \geq N$ implies $d(p_n, p) < \varepsilon$. In this case we say that $\{p_n\}$ converges to $p$, or that $p$ is the limit of $\{p_n\}$, and we write $p_n \to p$, or
\[
\lim_{n \to \infty} p_n = p.
\]
In the language of neighborhoods: $p_n \to p$ if and only if every neighborhood of $p$ contains all but finitely many of the terms $p_n$.
\end{definition}

\begin{definition}\label{def:limit}
Let $X$ and $Y$ be metric spaces; suppose $E \subset X$, $f$ maps $E$ into $Y$, and $p$ is a limit point of $E$. We write $f(x) \to q$ as $x \to p$, or
\[
\lim_{x \to p} f(x) = q
\]
if there is a point $q \in Y$ with the following property: For every $\varepsilon > 0$ there exists a $\delta > 0$ such that
\[
d_Y(f(x), q) < \varepsilon
\]
for all points $x \in E$ for which
\[
0 < d_X(x, p) < \delta.
\]
The symbols $d_X$ and $d_Y$ refer to the distances in $X$ and $Y$, respectively. If $X$ and/or $Y$ are replaced by the real line, the complex plane, or by some euclidean space $\mathbb{R}^k$, the distances $d_X$, $d_Y$ are of course replaced by absolute values, or by norms of differences.
\end{definition}

\begin{remark}\label{rem:limitpoint}
We pay close attention to \textit{where} a limit is defined.

\begin{enumerate}
    \item \textbf{Limits are taken only at limit points of the domain.} The point $p$ in the definition is a limit point of $E$. It is a point of $X$, but it need \textit{not} be a point of $E$. For instance, let $E = (0,1] \subset \mathbb{R}$ and $f(x) = x$. The point $0$ is a limit point of $E$ which does not belong to $E$, and $\lim_{x\to 0} f(x) = 0$ (given $\varepsilon > 0$ take $\delta = \varepsilon$). So $f$ has a limit at a point where $f$ is not even defined.

    \item \textbf{The value at $p$ plays no role.} Even if $p \in E$, the condition $0 < d_X(x,p)$ excludes the point $x = p$ itself. Thus the value $f(p)$ (if it exists) is never used, and we may very well have $f(p) \neq \lim_{x\to p} f(x)$. For instance, if $f(x) = 0$ for $x \neq 0$ and $f(0) = 1$ on $\mathbb{R}$, then $\lim_{x\to 0} f(x) = 0 \neq f(0)$.

    \item \textbf{Limits are not defined at isolated points of the domain.} Suppose $p$ is an isolated point of $E$, say $E = (0,1] \cup \{2\}$ and $p = 2$. Then there is an $r > 0$ such that $N_r(p)$ contains no point of $E$ other than $p$. If we tried to use the definition with $\delta \leq r$, there would be \textit{no} point $x \in E$ with $0 < d_X(x,p) < \delta$, so the condition ``$d_Y(f(x),q) < \varepsilon$ for all such $x$'' would be satisfied vacuously, by \textit{every} $q \in Y$. Every point of $Y$ would be ``a limit'', and the notion would be meaningless. This is exactly why the definition demands that $p$ be a limit point of $E$: then every punctured neighborhood of $p$ contains points of $E$ (in fact infinitely many, by Theorem 1.43), and the condition has real content. We shall see in Corollary \ref{cor:unique} that, at a limit point, the limit (if it exists) is unique.
\end{enumerate}
\end{remark}

\begin{theorem}\label{thm:seqlimit}
Let $X$, $Y$, $E$, $f$, and $p$ be as in Definition \ref{def:limit}. Then
\[
\lim_{x \to p} f(x) = q
\]
if and only if
\[
\lim_{n \to \infty} f(p_n) = q
\]
for every sequence $\{p_n\}$ in $E$ such that
\[
p_n \neq p, \qquad \lim_{n\to\infty} p_n = p.
\]
\end{theorem}

\begin{proof}
\textbf{Part 1: Assume $\lim_{x\to p} f(x) = q$.}\\
Let $\{p_n\}$ be any sequence in $E$ with $p_n \neq p$ and $p_n \to p$. We must show $f(p_n) \to q$.

Let $\varepsilon > 0$ be given. By the definition of the limit of $f$, there exists $\delta > 0$ such that $d_Y(f(x), q) < \varepsilon$ whenever $x \in E$ and $0 < d_X(x,p) < \delta$.

Now we use the convergence of the sequence, with this $\delta$ playing the role of ``$\varepsilon$'' in Definition \ref{def:seq}: there is an $N$ such that $d_X(p_n, p) < \delta$ for all $n \geq N$. Since $p_n \neq p$, we also have $d_X(p_n,p) > 0$. Thus for $n \geq N$, the point $x = p_n$ satisfies $x \in E$ and $0 < d_X(x,p) < \delta$, and hence
\[
d_Y(f(p_n), q) < \varepsilon \qquad (n \geq N).
\]
Since $\varepsilon > 0$ was arbitrary, $f(p_n) \to q$.

\textbf{Part 2: Assume $f(p_n) \to q$ for every such sequence.}\\
We argue by contradiction. Suppose $\lim_{x\to p} f(x) = q$ is false. Let us write carefully what this means. The definition says ``for every $\varepsilon > 0$ there is a $\delta > 0$ which works''. Its negation is: \textit{there is some $\varepsilon > 0$ such that no $\delta > 0$ works}. That is, there exists $\varepsilon > 0$ such that for every $\delta > 0$ there is a point $x \in E$ (depending on $\delta$) with
\[
0 < d_X(x, p) < \delta \quad \text{but} \quad d_Y(f(x), q) \geq \varepsilon.
\]

\textbf{Step 2.1: Constructing a sequence.}\\
Fix such an $\varepsilon$. Apply the above statement with $\delta = 1/n$ for $n = 1, 2, 3, \dots$. For each $n$ we obtain a point $p_n \in E$ such that
\[
0 < d_X(p_n, p) < \frac{1}{n} \quad \text{and} \quad d_Y(f(p_n), q) \geq \varepsilon.
\]
Let us understand why such points are available at all. Since $p$ is a limit point of $E$, every neighborhood $N_{1/n}(p)$ contains points of $E$ different from $p$; in fact, by Theorem 1.43, it contains infinitely many points of $E$. So the punctured neighborhoods of $p$ are never empty of points of $E$, and the negation of the limit statement tells us that, among these points, there is always at least one ``bad'' point where $f$ is at distance at least $\varepsilon$ from $q$. We pick one such bad point in each $N_{1/n}(p)$.

\textbf{Step 2.2: Properties of the sequence.}\\
By construction $p_n \in E$ and $p_n \neq p$ (because $d_X(p_n,p) > 0$). Moreover $p_n \to p$: given any $\eta > 0$, choose $N$ with $1/N < \eta$; then $d_X(p_n,p) < 1/n \leq 1/N < \eta$ for all $n \geq N$.

We also observe that this sequence has \textit{infinitely many distinct terms}. Indeed, suppose the set $\{p_n : n \in \mathbb{N}\}$ were finite, say equal to $\{y_1, \dots, y_m\}$. Each $y_i \neq p$, so
\[
r = \min_{1 \leq i \leq m} d_X(y_i, p) > 0,
\]
being the minimum of finitely many positive numbers. Choose $n$ with $1/n < r$. Then $d_X(p_n,p) < 1/n < r$, but $p_n$ is one of the $y_i$, for which $d_X(y_i,p) \geq r$. This is a contradiction. (This is exactly the argument of Theorem 1.43: a limit point cannot be approached using only finitely many points.) So the sequence genuinely runs through infinitely many different points of $E$ approaching $p$.

\textbf{Step 2.3: The contradiction.}\\
We have produced a sequence $\{p_n\}$ in $E$ with $p_n \neq p$ and $p_n \to p$. By hypothesis, $f(p_n) \to q$, so $d_Y(f(p_n), q) < \varepsilon$ for all large $n$. But by construction $d_Y(f(p_n),q) \geq \varepsilon$ for every $n$. This contradiction shows that $\lim_{x\to p} f(x) = q$.
\end{proof}

\begin{corollary}\label{cor:unique}
If $f$ has a limit at $p$, this limit is unique.
\end{corollary}

\begin{proof}
Suppose $\lim_{x\to p} f(x) = q$ and $\lim_{x\to p} f(x) = q'$. We show that $q = q'$.

\textbf{Step 1: A sequence approaching $p$.}\\
Since $p$ is a limit point of $E$, for each $n$ the neighborhood $N_{1/n}(p)$ contains a point $p_n \in E$ with $p_n \neq p$. Then $d_X(p_n,p) < 1/n$, so $p_n \to p$, and $p_n \neq p$.

\textbf{Step 2: Apply the theorem.}\\
By Theorem \ref{thm:seqlimit}, $f(p_n) \to q$ and also $f(p_n) \to q'$.

\textbf{Step 3: A sequence has at most one limit.}\\
Let $\varepsilon > 0$. There are $N_1$ and $N_2$ such that $d_Y(f(p_n), q) < \varepsilon$ for $n \geq N_1$ and $d_Y(f(p_n), q') < \varepsilon$ for $n \geq N_2$. Take $n = \max(N_1, N_2)$. By the triangle inequality,
\[
d_Y(q, q') \leq d_Y(q, f(p_n)) + d_Y(f(p_n), q') < 2\varepsilon.
\]
Since $\varepsilon > 0$ is arbitrary, the non-negative number $d_Y(q,q')$ is smaller than every positive number, so $d_Y(q,q') = 0$, that is, $q = q'$.
\end{proof}

\begin{definition}\label{def:algebra}
Suppose we have two complex functions $f$ and $g$, both defined on a set $E$. By $f + g$ we mean the function which assigns to each point $x$ of $E$ the number $f(x) + g(x)$. Similarly we define the difference $f - g$, the product $fg$, and the quotient $f/g$ of the two functions:
\[
(f+g)(x) = f(x) + g(x), \quad (f-g)(x) = f(x) - g(x), \quad (fg)(x) = f(x)g(x), \quad \Big(\frac{f}{g}\Big)(x) = \frac{f(x)}{g(x)},
\]
with the understanding that the quotient is defined only at those points $x$ of $E$ at which $g(x) \neq 0$. If $f$ assigns to each point $x$ of $E$ the same number $c$, then $f$ is said to be a \textit{constant function}, or simply a \textit{constant}, and we write $f = c$. If $\lambda$ is a complex number, $\lambda f$ is the function $(\lambda f)(x) = \lambda f(x)$. If $f$ and $g$ are real functions, and if $f(x) \geq g(x)$ for every $x \in E$, we shall sometimes write $f \geq g$, for brevity.

Similarly, if $f$ and $g$ map $E$ into $\mathbb{R}^k$, we define $f + g$ and $f \cdot g$ by
\[
(f+g)(x) = f(x) + g(x), \qquad (f \cdot g)(x) = f(x) \cdot g(x);
\]
and if $\lambda$ is a real number, $(\lambda f)(x) = \lambda f(x)$. Here $f(x) \cdot g(x)$ is the inner (dot) product of the two vectors $f(x)$ and $g(x)$ in $\mathbb{R}^k$, so $f\cdot g$ is a \textit{real} function on $E$.
\end{definition}

\begin{theorem}\label{thm:limitalgebra}
Suppose $E \subset X$, a metric space, $p$ is a limit point of $E$, $f$ and $g$ are complex functions on $E$, and
\[
\lim_{x \to p} f(x) = A, \qquad \lim_{x \to p} g(x) = B.
\]
Then
\begin{enumerate}[label=(\alph*)]
    \item $\displaystyle \lim_{x \to p} (f+g)(x) = A + B$;
    \item $\displaystyle \lim_{x \to p} (fg)(x) = AB$;
    \item $\displaystyle \lim_{x \to p} \Big(\frac{f}{g}\Big)(x) = \frac{A}{B}$, if $B \neq 0$.
\end{enumerate}
\end{theorem}

\begin{proof}
The idea in every part is the same: we want to make the \textit{combined} quantity close to its expected limit, and we do this by making $f(x)$ close to $A$ and $g(x)$ close to $B$ by suitable amounts. Each of the two hypotheses supplies its own $\delta$; to have both closeness conditions at the same time we take the \textit{smaller} of the two $\delta$'s.

\textbf{Part (a): The sum.}\\
Let $\varepsilon > 0$. Using the hypothesis on $f$ with $\varepsilon/2$ in place of $\varepsilon$, there is $\delta_1 > 0$ such that $|f(x) - A| < \varepsilon/2$ for $x \in E$, $0 < d_X(x,p) < \delta_1$. Similarly there is $\delta_2 > 0$ such that $|g(x) - B| < \varepsilon/2$ for $x \in E$, $0 < d_X(x,p) < \delta_2$. Put $\delta = \min(\delta_1, \delta_2) > 0$. If $x \in E$ and $0 < d_X(x,p) < \delta$, then both inequalities hold, and by the triangle inequality
\[
|(f+g)(x) - (A+B)| = |(f(x) - A) + (g(x) - B)| \leq |f(x) - A| + |g(x) - B| < \frac{\varepsilon}{2} + \frac{\varepsilon}{2} = \varepsilon.
\]
Hence $(f+g)(x) \to A+B$.

\textbf{Part (b): The product.}\\
We begin with the algebraic identity
\[
f(x)g(x) - AB = \big(f(x) - A\big) g(x) + A \big(g(x) - B\big),
\]
which is verified by expanding the right side. It shows that $fg - AB$ is small if $f - A$ and $g - B$ are small, \textit{provided} the factor $g(x)$ does not become large. So we first control $g$.

\textit{Step 1: $g$ stays bounded near $p$.} Using the hypothesis on $g$ with $\varepsilon = 1$, there is $\delta_0 > 0$ such that $|g(x) - B| < 1$ for $x \in E$, $0 < d_X(x,p) < \delta_0$. For such $x$,
\[
|g(x)| \leq |g(x) - B| + |B| < |B| + 1.
\]

\textit{Step 2: Choosing the $\delta$.} Let $\varepsilon > 0$. Choose $\delta_1 > 0$ such that
\[
|f(x) - A| < \frac{\varepsilon}{2(|B|+1)} \quad (x \in E,\ 0 < d_X(x,p) < \delta_1),
\]
and $\delta_2 > 0$ such that
\[
|g(x) - B| < \frac{\varepsilon}{2(|A|+1)} \quad (x \in E,\ 0 < d_X(x,p) < \delta_2).
\]
(We divide by $|A| + 1$ and $|B|+1$ rather than by $|A|$ and $|B|$, simply to avoid dividing by zero.) Put $\delta = \min(\delta_0, \delta_1, \delta_2)$.

\textit{Step 3: Estimate.} For $x \in E$ with $0 < d_X(x,p) < \delta$, all three estimates hold and
\[
|f(x)g(x) - AB| \leq |f(x) - A|\,|g(x)| + |A|\,|g(x) - B| < \frac{\varepsilon}{2(|B|+1)}(|B|+1) + |A| \frac{\varepsilon}{2(|A|+1)} < \frac{\varepsilon}{2} + \frac{\varepsilon}{2} = \varepsilon.
\]
Hence $(fg)(x) \to AB$.

\textbf{Part (c): The quotient.}\\
Assume $B \neq 0$. We first show that $1/g(x) \to 1/B$, and then write $f/g = f \cdot (1/g)$ and use part (b).

\textit{Step 1: $g$ stays away from zero near $p$.} Applying the hypothesis on $g$ with $\varepsilon = |B|/2 > 0$, there is $\delta_0 > 0$ such that $|g(x) - B| < |B|/2$ for $x \in E$, $0 < d_X(x,p) < \delta_0$. For such $x$,
\[
|g(x)| \geq |B| - |B - g(x)| > |B| - \frac{|B|}{2} = \frac{|B|}{2} > 0.
\]
In particular $g(x) \neq 0$, so $f/g$ is defined at all points $x \in E$ with $0 < d_X(x,p) < \delta_0$. Consequently $p$ is still a limit point of the set $E' = \{x \in E : g(x) \neq 0\}$ on which $f/g$ is defined: every neighborhood $N_r(p)$ contains the smaller neighborhood $N_{\min(r,\delta_0)}(p)$, which contains a point $x \in E$, $x \neq p$, and such a point lies in $E'$.

\textit{Step 2: $1/g \to 1/B$.} For $x \in E$, $0 < d_X(x,p) < \delta_0$, we have
\[
\Big| \frac{1}{g(x)} - \frac{1}{B} \Big| = \frac{|B - g(x)|}{|g(x)|\,|B|} < \frac{2}{|B|^2} |g(x) - B|.
\]
Given $\varepsilon > 0$, choose $\delta_1 > 0$ with $|g(x) - B| < \varepsilon |B|^2 / 2$ for $x \in E$, $0 < d_X(x,p) < \delta_1$, and put $\delta = \min(\delta_0, \delta_1)$. Then $|1/g(x) - 1/B| < \varepsilon$ for $x \in E$, $0 < d_X(x,p) < \delta$. Thus $1/g(x) \to 1/B$.

\textit{Step 3: Conclusion.} Regard $f$ and $1/g$ as functions on $E'$. They still satisfy $f(x) \to A$ and $1/g(x) \to 1/B$ as $x \to p$ (the defining inequalities now have to be checked only at the points of the smaller set $E'$, so the same $\delta$'s work). By part (b),
\[
\frac{f(x)}{g(x)} = f(x) \cdot \frac{1}{g(x)} \longrightarrow A \cdot \frac{1}{B} = \frac{A}{B}. \qedhere
\]
\end{proof}

\begin{remark}\label{rem:Rk}
If $f$ and $g$ map $E$ into $\mathbb{R}^k$, with $\lim_{x\to p} f(x) = A$ and $\lim_{x\to p} g(x) = B$ ($A, B \in \mathbb{R}^k$), then (a) remains true, and (b) becomes
\[
\text{(b$'$)} \qquad \lim_{x \to p} (f \cdot g)(x) = A \cdot B,
\]
where the product is the dot product. Also $\lim_{x \to p} (\lambda f)(x) = \lambda A$ for every real $\lambda$. We prove these.

\textbf{Proof of (a) in $\mathbb{R}^k$.} The proof of part (a) above used only the triangle inequality, which holds for the norm in $\mathbb{R}^k$. So the same proof, word for word, gives $(f+g)(x) \to A + B$.

\textbf{Proof of $\lambda f \to \lambda A$.} $|\lambda f(x) - \lambda A| = |\lambda|\,|f(x) - A|$. Given $\varepsilon > 0$, choose $\delta$ such that $|f(x) - A| < \varepsilon/(|\lambda|+1)$ for $x \in E$, $0 < d_X(x,p) < \delta$; then $|\lambda f(x) - \lambda A| < \varepsilon$.

\textbf{Proof of (b$'$).} Write $f = (f_1, \dots, f_k)$, $g = (g_1,\dots,g_k)$, $A = (A_1, \dots, A_k)$ and $B = (B_1, \dots, B_k)$, where $f_j$, $g_j$ are real functions on $E$.

\textit{Step 1: Each component converges.} For each $j$,
\[
|f_j(x) - A_j| = \big( (f_j(x) - A_j)^2 \big)^{1/2} \leq \Big( \sum_{i=1}^k (f_i(x) - A_i)^2 \Big)^{1/2} = |f(x) - A|,
\]
because a sum of non-negative squares is at least as large as any single one of them. So whenever $|f(x) - A| < \varepsilon$, also $|f_j(x) - A_j| < \varepsilon$, and the same $\delta$ shows $f_j(x) \to A_j$. Similarly $g_j(x) \to B_j$.

\textit{Step 2: Products and sums of the components.} The functions $f_j$, $g_j$ are real, hence complex, functions. By part (b) of the theorem, $f_j(x) g_j(x) \to A_j B_j$ for each $j$. Now
\[
(f \cdot g)(x) = \sum_{j=1}^k f_j(x) g_j(x),
\]
which is a sum of $k$ functions. Applying part (a) of the theorem repeatedly ($k-1$ times; formally, by induction on the number of summands), we get
\[
(f\cdot g)(x) \to \sum_{j=1}^k A_j B_j = A \cdot B.
\]
This proves (b$'$).
\end{remark}

\subsection{Continuous Functions}

\begin{definition}\label{def:cont}
Suppose $X$ and $Y$ are metric spaces, $E \subset X$, $p \in E$, and $f$ maps $E$ into $Y$. Then $f$ is said to be \textit{continuous at $p$} if for every $\varepsilon > 0$ there exists a $\delta > 0$ such that
\[
d_Y(f(x), f(p)) < \varepsilon
\]
for all points $x \in E$ for which $d_X(x, p) < \delta$.

If $f$ is continuous at every point of $E$, then $f$ is said to be \textit{continuous on $E$}.
\end{definition}

\begin{remark}\label{rem:contcompare}
\textbf{1. Comparison with the definition of limit.} Place Definition \ref{def:limit} and Definition \ref{def:cont} side by side.
\begin{itemize}
    \item For a \textit{limit} at $p$: $p$ must be a limit point of $E$; it may or may not belong to $E$; the value $f(p)$ never enters; the points $x$ considered satisfy $0 < d_X(x,p) < \delta$ (the point $p$ is excluded); and $q$ is some point of $Y$.
    \item For \textit{continuity} at $p$: $p$ must be a point of $E$ (we need the value $f(p)$); $p$ may be a limit point of $E$ or an isolated point of $E$; the points $x$ considered satisfy $d_X(x,p) < \delta$ (the point $p$ is allowed, but for $x = p$ the inequality $d_Y(f(p), f(p)) = 0 < \varepsilon$ is automatic); and the role of $q$ is played by $f(p)$ itself.
\end{itemize}
So continuity is defined at \textit{every point of the domain}, including isolated points, but it is \textit{not} defined at limit points of $E$ which lie outside $E$. For instance, for $f(x) = x$ on $E = (0,1]$, the limit at $0$ exists, but it makes no sense to ask whether $f$ is continuous at $0$, since $f(0)$ is not defined.

\textbf{2. Isolated points.} If $p$ is an isolated point of $E$, then \textit{every} function $f$ which has $E$ as its domain of definition is continuous at $p$. Indeed, since $p$ is not a limit point of $E$, there is an $r > 0$ such that the neighborhood $N_r(p)$ contains no point of $E$ other than $p$. Now let $\varepsilon > 0$ be arbitrary, and take $\delta = r$. The only point $x \in E$ with $d_X(x,p) < \delta$ is $x = p$ itself, and for it $d_Y(f(x), f(p)) = 0 < \varepsilon$. So the defining condition holds, no matter what $\varepsilon$ was and no matter what $f$ is. For instance, on $E = (0,1] \cup \{2\}$ the function which is $0$ on $(0,1]$ and $100$ at $2$ is continuous at $2$; and every function defined on $\mathbb{N}$ (as a subset of $\mathbb{R}$) is continuous at every point of $\mathbb{N}$.

\textbf{3. The neighborhood formulation.} The definition can be read as follows:
\begin{quote}
\textit{$f$ is continuous at $p$ if and only if for every neighborhood $U$ of $f(p)$ in $Y$ there is a neighborhood $V$ of $p$ in $X$ such that every point $x \in V \cap E$ is mapped into $U$, that is, $f(V \cap E) \subset U$.}
\end{quote}
In words: whatever small target $U$ around $f(p)$ is prescribed, we can find a small window $V$ around $p$ which $f$ sends entirely into the target.

To see that this is the same as the $\varepsilon$-$\delta$ definition, recall (Definition 1.33) that a neighborhood is a ball: $U = N_\varepsilon(f(p))$ and $V = N_\delta(p)$. Now
\[
d_Y(f(x), f(p)) < \varepsilon \iff f(x) \in N_\varepsilon(f(p)) = U, \qquad d_X(x,p) < \delta \iff x \in N_\delta(p) = V.
\]
So ``for every $\varepsilon > 0$ there is $\delta > 0$ such that $x \in E$, $d_X(x,p) < \delta$ implies $d_Y(f(x), f(p)) < \varepsilon$'' is, sentence for sentence, ``for every neighborhood $U$ of $f(p)$ there is a neighborhood $V$ of $p$ such that $x \in V \cap E$ implies $f(x) \in U$''. The two definitions are the same statement written in two languages.

In fact, we may even let $U$ be \textit{any open set} containing $f(p)$. If $G$ is open in $Y$ and $f(p) \in G$, then $f(p)$ is an interior point of $G$, so some $N_\varepsilon(f(p)) \subset G$; a neighborhood $V$ which works for $N_\varepsilon(f(p))$ then also works for $G$. Conversely every neighborhood is an open set (Theorem 1.42). So nothing changes if ``neighborhood $U$'' is replaced by ``open set $U$ containing $f(p)$''.

In exactly the same way, Definition \ref{def:limit} can be written in the neighborhood language:
\begin{quote}
\textit{$\lim_{x\to p} f(x) = q$ if and only if for every neighborhood $U$ of $q$ there is a neighborhood $V$ of $p$ such that $f(x) \in U$ for all $x \in V \cap E$ with $x \neq p$.}
\end{quote}
This formulation does not mention distances explicitly; it will be the key to defining limits at $\pm\infty$ in Section \ref{sec:infinite}.
\end{remark}

\begin{remark}\label{rem:local}
Two simple consequences of the definition will be used repeatedly.
\begin{enumerate}[label=(\alph*)]
    \item \textbf{Restriction.} If $f$ is continuous at $p$ as a function on $E$, and $A \subset E$ with $p \in A$, then the restriction of $f$ to $A$ is continuous at $p$. Indeed, the same $\delta$ works: the condition now has to be checked only for the points $x$ of the smaller set $A$.
    \item \textbf{Continuity is local.} Suppose $f$ and $g$ are defined on $E$, $p \in E$, and $f(x) = g(x)$ for all $x \in E \cap N_r(p)$ for some $r > 0$. If $g$ is continuous at $p$, then so is $f$. Indeed, given $\varepsilon > 0$, let $\delta$ work for $g$ and put $\delta' = \min(\delta, r)$; for $x \in E$ with $d_X(x,p) < \delta'$ we have $f(x) = g(x)$ and $f(p) = g(p)$, so $d_Y(f(x), f(p)) = d_Y(g(x), g(p)) < \varepsilon$. The same remark (with $x \neq p$) applies to limits: the limit of $f$ at $p$ depends only on the values of $f$ near $p$.
\end{enumerate}
\end{remark}

\begin{theorem}\label{thm:contlimit}
In the situation given in Definition \ref{def:cont}, assume also that $p$ is a limit point of $E$. Then $f$ is continuous at $p$ if and only if
\[
\lim_{x \to p} f(x) = f(p).
\]
\end{theorem}

\begin{proof}
Since $p$ is a limit point of $E$, the limit of $f$ at $p$ makes sense, and since $p \in E$, the value $f(p)$ makes sense. We compare the two definitions with $q = f(p)$.

\textbf{Part 1: Assume $f$ is continuous at $p$.}\\
Let $\varepsilon > 0$. By continuity there is $\delta > 0$ such that $d_Y(f(x), f(p)) < \varepsilon$ for all $x \in E$ with $d_X(x,p) < \delta$. In particular, this holds for all $x \in E$ with $0 < d_X(x,p) < \delta$ (we are just looking at fewer points). This is precisely the statement $\lim_{x\to p} f(x) = f(p)$.

\textbf{Part 2: Assume $\lim_{x\to p} f(x) = f(p)$.}\\
Let $\varepsilon > 0$. There is $\delta > 0$ such that $d_Y(f(x), f(p)) < \varepsilon$ for all $x \in E$ with $0 < d_X(x,p) < \delta$. The points $x \in E$ with $d_X(x,p) < \delta$ are these points together with the single point $x = p$. For $x = p$ we have $d_Y(f(p), f(p)) = 0 < \varepsilon$. Hence $d_Y(f(x), f(p)) < \varepsilon$ for all $x \in E$ with $d_X(x,p) < \delta$, which is continuity at $p$.
\end{proof}

Thus, at a point $p \in E$: if $p$ is isolated, $f$ is automatically continuous at $p$ (Remark \ref{rem:contcompare}); if $p$ is a limit point of $E$, $f$ is continuous at $p$ exactly when the limit of $f$ at $p$ exists and equals the value $f(p)$.

\begin{theorem}\label{thm:composition}
Suppose $X$, $Y$, $Z$ are metric spaces, $E \subset X$, $f$ maps $E$ into $Y$, $g$ maps the range of $f$, $f(E)$, into $Z$, and $h$ is the mapping of $E$ into $Z$ defined by
\[
h(x) = g(f(x)) \qquad (x \in E).
\]
If $f$ is continuous at a point $p \in E$ and if $g$ is continuous at the point $f(p)$, then $h$ is continuous at $p$.

This function $h$ is called the \textit{composition} or the \textit{composite} of $f$ and $g$. The notation $h = g \circ f$ is frequently used.
\end{theorem}

\textbf{The idea of the proof.} We have to produce, for a given $\varepsilon > 0$, a $\delta > 0$ such that $d_Z(h(x), h(p)) < \varepsilon$ whenever $x \in E$ and $d_X(x,p) < \delta$. The trick is to work \textit{backwards}, from $Z$ to $Y$ to $X$:
\begin{itemize}
    \item We want $g(f(x))$ to be within $\varepsilon$ of $g(f(p))$. The continuity of $g$ at $f(p)$ says this happens as soon as the point $y = f(x)$ is close enough to $f(p)$, say within a distance $\eta$.
    \item So now we want $f(x)$ to be within $\eta$ of $f(p)$. The continuity of $f$ at $p$, applied with $\eta$ \textit{in the role of epsilon}, says this happens as soon as $x$ is within some distance $\delta$ of $p$.
\end{itemize}
Thus the ``$\delta$'' of $g$ becomes the ``$\varepsilon$'' of $f$; this is the whole proof.

\begin{proof}
Let $\varepsilon > 0$ be given.

\textbf{Step 1: Use the continuity of $g$ at $f(p)$.}\\
Since $g$ is continuous at $f(p)$, there exists $\eta > 0$ such that
\[
d_Z(g(y), g(f(p))) < \varepsilon \quad \text{if } d_Y(y, f(p)) < \eta \text{ and } y \in f(E).
\]

\textbf{Step 2: Use the continuity of $f$ at $p$, with $\eta$ in place of $\varepsilon$.}\\
Since $f$ is continuous at $p$ and $\eta > 0$, there exists $\delta > 0$ such that
\[
d_Y(f(x), f(p)) < \eta \quad \text{if } d_X(x,p) < \delta \text{ and } x \in E.
\]

\textbf{Step 3: Combine.}\\
Let $x \in E$ with $d_X(x,p) < \delta$. By Step 2, the point $y = f(x)$ belongs to $f(E)$ and satisfies $d_Y(y, f(p)) < \eta$. By Step 1, applied to this $y$,
\[
d_Z(h(x), h(p)) = d_Z(g(f(x)), g(f(p))) < \varepsilon.
\]
Since $\varepsilon > 0$ was arbitrary, $h$ is continuous at $p$.
\end{proof}

\begin{theorem}\label{thm:openpre}
A mapping $f$ of a metric space $X$ into a metric space $Y$ is continuous on $X$ if and only if $f^{-1}(V)$ is open in $X$ for every open set $V$ in $Y$.
\end{theorem}

\begin{proof}
\textbf{Part 1: Assume $f$ is continuous on $X$; let $V$ be open in $Y$.}\\
We have to show that every point of $f^{-1}(V)$ is an interior point of $f^{-1}(V)$.

\textit{Step 1: Pick a point.} Let $p \in f^{-1}(V)$, that is, $p \in X$ and $f(p) \in V$.

\textit{Step 2: Use that $V$ is open.} Since $V$ is open, $f(p)$ is an interior point of $V$: there is $\varepsilon > 0$ such that $y \in V$ whenever $d_Y(f(p), y) < \varepsilon$. In other words, $N_\varepsilon(f(p)) \subset V$.

\textit{Step 3: Use continuity at $p$ with this $\varepsilon$.} Since $f$ is continuous at $p$, there is $\delta > 0$ such that $d_Y(f(x), f(p)) < \varepsilon$ whenever $d_X(x,p) < \delta$.

\textit{Step 4: Conclude.} Let $x \in N_\delta(p)$. By Step 3, $f(x) \in N_\varepsilon(f(p))$, and by Step 2, $f(x) \in V$; that is, $x \in f^{-1}(V)$. Hence $N_\delta(p) \subset f^{-1}(V)$, so $p$ is an interior point of $f^{-1}(V)$. Since $p$ was an arbitrary point of $f^{-1}(V)$, the set $f^{-1}(V)$ is open.

\textbf{Part 2: Assume $f^{-1}(V)$ is open in $X$ for every open set $V$ in $Y$.}\\
Fix $p \in X$ and $\varepsilon > 0$. We find a $\delta$.

\textit{Step 1: Choose the open set.} Let $V$ be the set of all $y \in Y$ such that $d_Y(y, f(p)) < \varepsilon$, that is, $V = N_\varepsilon(f(p))$. By Theorem 1.42, $V$ is open.

\textit{Step 2: Pull it back.} By hypothesis, $f^{-1}(V)$ is open in $X$. It contains $p$, because $f(p) \in V$. Hence $p$ is an interior point of $f^{-1}(V)$: there is $\delta > 0$ such that $x \in f^{-1}(V)$ whenever $d_X(x,p) < \delta$.

\textit{Step 3: Conclude.} If $d_X(x,p) < \delta$, then $x \in f^{-1}(V)$, so $f(x) \in V$, which means $d_Y(f(x), f(p)) < \varepsilon$. Thus $f$ is continuous at $p$. Since $p$ was arbitrary, $f$ is continuous on $X$.
\end{proof}

\begin{corollary}\label{cor:closedpre}
A mapping $f$ of a metric space $X$ into a metric space $Y$ is continuous if and only if $f^{-1}(C)$ is closed in $X$ for every closed set $C$ in $Y$.
\end{corollary}

\begin{proof}
\textbf{Step 1: Preimages respect complements.} For every $C \subset Y$,
\[
f^{-1}(C^c) = \big[f^{-1}(C)\big]^c.
\]
Indeed, for $x \in X$: $x \in f^{-1}(C^c) \iff f(x) \in C^c \iff f(x) \notin C \iff x \notin f^{-1}(C) \iff x \in [f^{-1}(C)]^c$.

\textbf{Step 2: Assume $f$ is continuous; let $C$ be closed in $Y$.} Then $C^c$ is open (Corollary 1.45.1), so $f^{-1}(C^c)$ is open by Theorem \ref{thm:openpre}. By Step 1, $[f^{-1}(C)]^c$ is open, hence $f^{-1}(C)$ is closed (Theorem 1.45).

\textbf{Step 3: Assume $f^{-1}(C)$ is closed for every closed $C$; let $V$ be open in $Y$.} Then $V^c$ is closed (Theorem 1.45), so $f^{-1}(V^c) = [f^{-1}(V)]^c$ is closed, hence $f^{-1}(V)$ is open (Corollary 1.45.1). By Theorem \ref{thm:openpre}, $f$ is continuous.
\end{proof}

\begin{remark}\label{rem:topology}
Theorem \ref{thm:openpre} is exactly the definition of a continuous function in \textit{topology}: in a general topological space one is only told which sets are open; there is no distance, so there are no $\varepsilon$'s and $\delta$'s, and continuity is \textit{defined} by the requirement that inverse images of open sets be open. In a metric space, the open sets are determined by the metric (Definition 1.33), and the theorem says that the $\varepsilon$-$\delta$ definition and the topological definition are \textit{equivalent}: both single out exactly the same class of functions. This is what ``equivalent statement'' means here. So in metric spaces we may use whichever is more convenient: $\varepsilon$-$\delta$ for computations at a point, and open (or closed) sets for global arguments, as we shall do repeatedly with compactness and connectedness.

Note that the theorem characterizes continuity \textit{on all of $X$}, not continuity at a single point. Also, the domain here is the whole metric space $X$; see Remark \ref{rem:subset} for functions defined on a subset.
\end{remark}

\begin{theorem}\label{thm:contalgebra}
Let $f$ and $g$ be complex continuous functions on a metric space $X$. Then $f + g$, $fg$, and $f/g$ are continuous on $X$.

In the last case, we must of course assume that $g(x) \neq 0$, for all $x \in X$.
\end{theorem}

\begin{proof}
Fix a point $p \in X$. There are two possibilities.

\textbf{Case 1: $p$ is an isolated point of $X$.} Then every function on $X$ is continuous at $p$ (Remark \ref{rem:contcompare}), in particular $f+g$, $fg$ and $f/g$. There is nothing to prove.

\textbf{Case 2: $p$ is a limit point of $X$.} Since $f$ and $g$ are continuous at $p$, Theorem \ref{thm:contlimit} gives
\[
\lim_{x\to p} f(x) = f(p), \qquad \lim_{x\to p} g(x) = g(p).
\]
By Theorem \ref{thm:limitalgebra}(a),
\[
\lim_{x \to p} (f+g)(x) = f(p) + g(p) = (f+g)(p),
\]
and so, by Theorem \ref{thm:contlimit} again (now in the reverse direction), $f+g$ is continuous at $p$. In the same way, Theorem \ref{thm:limitalgebra}(b) gives $\lim_{x\to p} (fg)(x) = f(p)g(p) = (fg)(p)$, so $fg$ is continuous at $p$. Finally, since $g(p) \neq 0$, Theorem \ref{thm:limitalgebra}(c) applies with $B = g(p) \neq 0$ and gives $\lim_{x\to p} (f/g)(x) = f(p)/g(p) = (f/g)(p)$, so $f/g$ is continuous at $p$.

Since $p$ was arbitrary, the three functions are continuous on $X$. Note also that constant functions are continuous ($d(c,c) = 0 < \varepsilon$ for any $\delta$), so $\lambda f$ (the product of the constant $\lambda$ and $f$) and $f - g = f + (-1)g$ are continuous as well.
\end{proof}

\begin{theorem}\label{thm:components}
\begin{enumerate}[label=(\alph*)]
    \item Let $f_1, \dots, f_k$ be real functions on a metric space $X$, and let $\mathbf{f}$ be the mapping of $X$ into $\mathbb{R}^k$ defined by
    \[
    \mathbf{f}(x) = \big(f_1(x), \dots, f_k(x)\big) \qquad (x \in X);
    \]
    then $\mathbf{f}$ is continuous if and only if each of the functions $f_1, \dots, f_k$ is continuous.
    \item If $\mathbf{f}$ and $\mathbf{g}$ are continuous mappings of $X$ into $\mathbb{R}^k$, then $\mathbf{f} + \mathbf{g}$ and $\mathbf{f} \cdot \mathbf{g}$ are continuous on $X$.
\end{enumerate}
The functions $f_1, \dots, f_k$ are called the \textit{components} of $\mathbf{f}$. Note that $\mathbf{f} + \mathbf{g}$ is a mapping into $\mathbb{R}^k$, whereas $\mathbf{f} \cdot \mathbf{g}$ is a real function on $X$.
\end{theorem}

\begin{proof}
\textbf{Part (a).} Everything rests on the following inequality and equality, valid for all $x, y \in X$ and each $j = 1, \dots, k$:
\[
|f_j(x) - f_j(y)| \;\leq\; |\mathbf{f}(x) - \mathbf{f}(y)| \;=\; \Big\{ \sum_{i=1}^k |f_i(x) - f_i(y)|^2 \Big\}^{1/2}. \tag{$*$}
\]
The equality is nothing but the definition of the euclidean norm of the vector $\mathbf{f}(x) - \mathbf{f}(y) = \big(f_1(x) - f_1(y), \dots, f_k(x) - f_k(y)\big)$. The inequality holds because the sum of the $k$ non-negative squares is at least as large as its $j$-th term $|f_j(x)-f_j(y)|^2$; taking square roots gives the inequality.

The two halves of $(*)$ give the two directions of (a).

\textit{($\Rightarrow$) $\mathbf{f}$ continuous implies each $f_j$ continuous: use the inequality.} The inequality says that \textit{a component cannot move more than the whole vector}. Fix $p \in X$ and $\varepsilon > 0$. Since $\mathbf{f}$ is continuous at $p$, there is $\delta > 0$ with $|\mathbf{f}(x) - \mathbf{f}(p)| < \varepsilon$ whenever $d_X(x,p) < \delta$. For the same $x$, by the inequality in $(*)$ with $y = p$,
\[
|f_j(x) - f_j(p)| \leq |\mathbf{f}(x) - \mathbf{f}(p)| < \varepsilon.
\]
So the same $\delta$ works for every component, and each $f_j$ is continuous at $p$.

\textit{($\Leftarrow$) each $f_j$ continuous implies $\mathbf{f}$ continuous: use the equality.} The equality says that \textit{the vector is controlled by its components}: if all $k$ components are small, the vector is small. Fix $p \in X$ and $\varepsilon > 0$. We want the sum of $k$ squares in $(*)$ to be less than $\varepsilon^2$; it suffices that each square be less than $\varepsilon^2/k$, i.e., each component difference be less than $\varepsilon/\sqrt{k}$. So, for each $i = 1, \dots, k$, by continuity of $f_i$ at $p$ choose $\delta_i > 0$ with
\[
|f_i(x) - f_i(p)| < \frac{\varepsilon}{\sqrt{k}} \qquad \text{whenever } d_X(x,p) < \delta_i.
\]
Let $\delta = \min(\delta_1, \dots, \delta_k) > 0$ (a minimum of finitely many positive numbers). If $d_X(x,p) < \delta$, then all $k$ estimates hold simultaneously, and by the equality in $(*)$,
\[
|\mathbf{f}(x) - \mathbf{f}(p)| = \Big\{\sum_{i=1}^k |f_i(x) - f_i(p)|^2\Big\}^{1/2} < \Big\{ k \cdot \frac{\varepsilon^2}{k} \Big\}^{1/2} = \varepsilon.
\]
So $\mathbf{f}$ is continuous at $p$. Since $p$ was arbitrary, (a) is proved in both directions.

\textbf{Part (b).} Let $\mathbf{f} = (f_1, \dots, f_k)$ and $\mathbf{g} = (g_1, \dots, g_k)$. By part (a), each $f_i$ and each $g_i$ is continuous.

\textit{The sum.} The components of $\mathbf{f} + \mathbf{g}$ are $f_1 + g_1, \dots, f_k + g_k$, each continuous by Theorem \ref{thm:contalgebra}. By part (a) (direction $\Leftarrow$), $\mathbf{f} + \mathbf{g}$ is continuous.

\textit{The dot product.} $(\mathbf{f}\cdot\mathbf{g})(x) = \sum_{i=1}^k f_i(x) g_i(x)$. Each product $f_i g_i$ is continuous by Theorem \ref{thm:contalgebra}, and a sum of finitely many continuous functions is continuous (Theorem \ref{thm:contalgebra} applied $k-1$ times). Hence $\mathbf{f}\cdot\mathbf{g}$ is continuous.
\end{proof}

\begin{example}\label{ex:polynomials}
We now build up a large stock of continuous functions, step by step, starting from the simplest ones.

\textbf{Step 1: The coordinate functions.} If $x_1, \dots, x_k$ are the coordinates of the point $\mathbf{x} \in \mathbb{R}^k$, the functions $\phi_i$ defined by
\[
\phi_i(\mathbf{x}) = x_i \qquad (\mathbf{x} \in \mathbb{R}^k)
\]
are called the \textit{coordinate functions}. Each $\phi_i$ is continuous on $\mathbb{R}^k$. Indeed, for $\mathbf{x}, \mathbf{y} \in \mathbb{R}^k$,
\[
|\phi_i(\mathbf{x}) - \phi_i(\mathbf{y})| = |x_i - y_i| \leq \Big\{\sum_{j=1}^k (x_j - y_j)^2\Big\}^{1/2} = |\mathbf{x} - \mathbf{y}|,
\]
since one square is at most the sum of all the squares. Given $\varepsilon > 0$, take $\delta = \varepsilon$: if $|\mathbf{x} - \mathbf{y}| < \delta$, then $|\phi_i(\mathbf{x}) - \phi_i(\mathbf{y})| < \varepsilon$. (Equivalently: the identity map $\mathbf{x} \mapsto \mathbf{x}$ of $\mathbb{R}^k$ is clearly continuous, and its components are exactly $\phi_1, \dots, \phi_k$, so Theorem \ref{thm:components}(a) applies.)

\textbf{Step 2: Constants.} Every constant function $c$ is continuous, since $|c - c| = 0 < \varepsilon$ for every $\varepsilon$, whatever $\delta$ we choose.

\textbf{Step 3: Monomials.} A \textit{monomial} is a function of the form
\[
\mathbf{x} \mapsto x_1^{n_1} x_2^{n_2} \cdots x_k^{n_k} \qquad (\mathbf{x} \in \mathbb{R}^k), \tag{M}
\]
where $n_1, \dots, n_k$ are non-negative integers. It is a product of $n_1 + \cdots + n_k$ coordinate functions ($n_1$ copies of $\phi_1$, $n_2$ copies of $\phi_2$, and so on); when all $n_i = 0$ it is the constant $1$. By Theorem \ref{thm:contalgebra}, the product of two continuous functions is continuous. Hence, by induction on the number of factors (if a product of $m$ coordinate functions is continuous, then multiplying it by one more continuous coordinate function gives a continuous product of $m+1$ factors), every monomial is continuous on $\mathbb{R}^k$.

\textbf{Step 4: Constant multiples of monomials.} If $c$ is a complex number, then $c \, x_1^{n_1} \cdots x_k^{n_k}$ is the product of the continuous constant function $c$ and a continuous monomial, so it is continuous (Theorem \ref{thm:contalgebra}).

\textbf{Step 5: Polynomials.} A \textit{polynomial} on $\mathbb{R}^k$ is a function of the form
\[
P(\mathbf{x}) = \sum c_{n_1 \cdots n_k} \, x_1^{n_1} \cdots x_k^{n_k} \qquad (\mathbf{x} \in \mathbb{R}^k), \tag{P}
\]
where the coefficients $c_{n_1 \cdots n_k}$ are complex numbers, $n_1, \dots, n_k$ are non-negative integers, and the sum in (P) has \textit{finitely many} terms. By Theorem \ref{thm:contalgebra}, the sum of two continuous functions is continuous; by induction on the number of terms, a finite sum of continuous functions is continuous. By Step 4, every polynomial is continuous on $\mathbb{R}^k$.

\textbf{Step 6: The case $k = 1$.} When $k = 1$, a monomial is just $x^n$, and (P) becomes
\[
P(x) = c_0 + c_1 x + c_2 x^2 + \cdots + c_n x^n \qquad (x \in \mathbb{R}),
\]
which is the familiar polynomial in one variable. So every polynomial in the usual sense is continuous on $\mathbb{R}$. The same argument, starting from the continuous identity function $z \mapsto z$ on the complex plane, shows that every complex polynomial $c_0 + c_1 z + \cdots + c_n z^n$ is continuous on $\mathbb{C}$.

\textbf{Step 7: Rational functions.} A \textit{rational function} in $x_1, \dots, x_k$ is a quotient $P/Q$ of two polynomials of the form (P). Let $X = \{\mathbf{x} \in \mathbb{R}^k : Q(\mathbf{x}) \neq 0\}$, regarded as a metric space with the distance of $\mathbb{R}^k$. The restrictions of $P$ and $Q$ to $X$ are continuous (Remark \ref{rem:local}(a)), and $Q$ does not vanish on $X$. By Theorem \ref{thm:contalgebra}, $P/Q$ is continuous on $X$. Thus \textit{every rational function is continuous wherever its denominator is different from zero}. For example, $1/x$ is continuous on $\mathbb{R} \setminus \{0\}$, and $\dfrac{x^2 + 1}{x - 1}$ is continuous on $\mathbb{R} \setminus \{1\}$.

\textbf{Step 8: The norm (modulus) function.} From the triangle inequality we obtain
\[
\big|\, |\mathbf{x}| - |\mathbf{y}| \,\big| \leq |\mathbf{x} - \mathbf{y}| \qquad (\mathbf{x}, \mathbf{y} \in \mathbb{R}^k).
\]
Indeed, $|\mathbf{x}| = |(\mathbf{x} - \mathbf{y}) + \mathbf{y}| \leq |\mathbf{x} - \mathbf{y}| + |\mathbf{y}|$ gives $|\mathbf{x}| - |\mathbf{y}| \leq |\mathbf{x} - \mathbf{y}|$; interchanging $\mathbf{x}$ and $\mathbf{y}$ gives $|\mathbf{y}| - |\mathbf{x}| \leq |\mathbf{x} - \mathbf{y}|$; the two together are the stated inequality. Hence, taking $\delta = \varepsilon$, the mapping $\mathbf{x} \mapsto |\mathbf{x}|$ is a continuous real function on $\mathbb{R}^k$. In particular the absolute value $x \mapsto |x|$ is continuous on $\mathbb{R}$, and the modulus $z \mapsto |z|$ is continuous on $\mathbb{C}$.

\textbf{Step 9: The modulus of a continuous function.} If now $\mathbf{f}$ is a continuous mapping from a metric space $X$ into $\mathbb{R}^k$, and if $\phi$ is defined on $X$ by setting $\phi(p) = |\mathbf{f}(p)|$, then $\phi$ is the composition of $\mathbf{f}$ (from $X$ to $\mathbb{R}^k$) followed by the norm function $\mathbf{y} \mapsto |\mathbf{y}|$ (from $\mathbb{R}^k$ to $\mathbb{R}$):
\[
\phi = |\cdot| \circ \mathbf{f}.
\]
Both are continuous, so by Theorem \ref{thm:composition}, $\phi$ is a continuous real function on $X$. For instance, $|x^2 - 1|$ and $|x - x_0|$ are continuous on $\mathbb{R}$, and $|f|$ is continuous whenever $f$ is a continuous real or complex function.
\end{example}

\begin{remark}\label{rem:subset}
We defined the notion of continuity for functions defined on a subset $E$ of a metric space $X$. Should we talk about continuous functions on a subset of a metric space, or on a full metric space? The answer is that talking only about functions on an entire metric space causes no problem at all.

Indeed, look again at Definition \ref{def:cont}. To check continuity at $p \in E$ we only look at points $x \in E$ and at the distances $d_X(x,p)$ between points of $E$. The complement of $E$ in $X$ plays \textit{no role} in the definition. Now $E$, with the distance $d_X$ restricted to pairs of points of $E$, is itself a metric space. Hence $f$ is continuous on $E$ (as a subset of $X$) if and only if $f$ is continuous as a mapping of the metric space $E$ into $Y$. So if we discard the complement of $E$, we lose nothing: the definitions and all operations with continuous functions remain exactly the same. This allows us to state theorems for continuous mappings of one metric space into another, rather than for mappings of subsets; this simplifies statements and proofs. We have already made use of this principle in Theorems \ref{thm:openpre} to \ref{thm:components} and in Example \ref{ex:polynomials}, and will continue to do so.

Two comments are in order.
\begin{enumerate}
    \item The situation is different for \textit{limits}: the limit point $p$ in Definition \ref{def:limit} may lie in $X \setminus E$, so for limits the complement of $E$ does matter.
    \item When Theorem \ref{thm:openpre} is applied to $f : E \to Y$, ``$f^{-1}(V)$ is open'' means open in the metric space $E$, i.e., open \textit{relative to} $E$; by Theorem 1.56 this means $f^{-1}(V) = E \cap G$ for some open set $G$ in $X$. For instance, $f(x) = 1/x$ is continuous on $E = (0,1]$, and $f^{-1}\big((1/2,2)\big) = \{x \in (0,1] : 1/2 < 1/x < 2\} = (1/2, 1]$, which is open relative to $(0,1]$ but not open in $\mathbb{R}$ (compare Example 1.53).
\end{enumerate}
\end{remark}

\subsection{Continuity and Compactness}

\begin{definition}\label{def:bounded}
A mapping $\mathbf{f}$ of a set $E$ into $\mathbb{R}^k$ is said to be \textit{bounded} if there is a real number $M$ such that $|\mathbf{f}(x)| \leq M$ for all $x \in E$.
\end{definition}

\begin{theorem}\label{thm:compactimage}
Suppose $f$ is a continuous mapping of a compact metric space $X$ into a metric space $Y$. Then $f(X)$ is compact.
\end{theorem}

\textbf{The idea of the proof.} We are given an open cover of $f(X)$ in $Y$ and we must extract a finite subcover. We cannot use compactness in $Y$ (we know nothing about $Y$); compactness is available only in $X$. So we \textit{pull the cover back} to $X$ by taking inverse images (which are open, by continuity), use the compactness of $X$ to choose finitely many of them, and then \textit{push forward} again to $Y$.

\begin{proof}
\textbf{Step 1: An open cover of $f(X)$.}\\
Let $\{V_\alpha\}$ be an open cover of $f(X)$; that is, each $V_\alpha$ is an open subset of $Y$ and $f(X) \subset \bigcup_\alpha V_\alpha$.

\textbf{Step 2: Pull back to $X$.}\\
Since $f$ is continuous, Theorem \ref{thm:openpre} shows that each of the sets $f^{-1}(V_\alpha)$ is open in $X$. These sets cover $X$: if $x \in X$, then $f(x) \in f(X) \subset \bigcup_\alpha V_\alpha$, so $f(x) \in V_\alpha$ for some $\alpha$, i.e., $x \in f^{-1}(V_\alpha)$. Thus $\{f^{-1}(V_\alpha)\}$ is an open cover of $X$.

\textbf{Step 3: Use compactness of $X$.}\\
Since $X$ is compact, there are finitely many indices, say $\alpha_1, \dots, \alpha_n$, such that
\[
X \subset f^{-1}(V_{\alpha_1}) \cup \cdots \cup f^{-1}(V_{\alpha_n}).
\]

\textbf{Step 4: Push forward to $Y$.}\\
Let $y \in f(X)$, say $y = f(x)$ with $x \in X$. By Step 3, $x \in f^{-1}(V_{\alpha_i})$ for some $i$, which means $y = f(x) \in V_{\alpha_i}$. Hence
\[
f(X) \subset V_{\alpha_1} \cup \cdots \cup V_{\alpha_n}.
\]
(In the language of sets, we have used the inclusion $f(f^{-1}(E)) \subset E$, valid for every $E \subset Y$; see Remark \ref{rem:imagepre}.) We have found a finite subcover of the arbitrary open cover $\{V_\alpha\}$, so $f(X)$ is compact.
\end{proof}

\begin{remark}\label{rem:imagepre}
Let $f$ be a mapping of a set $A$ into a set $B$. For $E \subset B$ and $F \subset A$ we always have
\[
\text{(i)}\ \ f(f^{-1}(E)) \subset E, \qquad \text{(ii)}\ \ F \subset f^{-1}(f(F)).
\]
\textit{Proof of (i):} if $y \in f(f^{-1}(E))$, then $y = f(x)$ for some $x \in f^{-1}(E)$, i.e., $f(x) \in E$; so $y \in E$. \textit{Proof of (ii):} if $x \in F$, then $f(x) \in f(F)$, so $x \in f^{-1}(f(F))$.

Equality may fail in both. Why do we lose equality?

\textbf{(i) and ontoness.} Going from $E$ to $f^{-1}(E)$ and back, we recover only those points of $E$ which are values of $f$:
\[
f(f^{-1}(E)) = E \cap f(A).
\]
Points of $E$ outside the range of $f$ are lost. Thus \textit{$f(f^{-1}(E)) = E$ for every $E \subset B$ if and only if $f$ is onto}. Indeed, if $f$ is onto, then $f(A) = B$ and $E \cap f(A) = E$. If $f$ is not onto, pick $y \in B \setminus f(A)$ and $E = \{y\}$; then $f^{-1}(E) = \emptyset$ and $f(f^{-1}(E)) = \emptyset \neq E$.

\textbf{(ii) and one-to-oneness.} Going from $F$ to $f(F)$ and back, we pick up every point which has the same value as some point of $F$. Thus \textit{$f^{-1}(f(F)) = F$ for every $F \subset A$ if and only if $f$ is 1-1}. Indeed, if $f$ is 1-1 and $x \in f^{-1}(f(F))$, then $f(x) = f(x')$ for some $x' \in F$, and so $x = x' \in F$. If $f$ is not 1-1, pick $x_1 \neq x_2$ with $f(x_1) = f(x_2)$ and $F = \{x_1\}$; then $f^{-1}(f(F))$ contains both $x_1$ and $x_2$, so it is not equal to $F$.

\textbf{Example.} Let $f(x) = x^2$ on $\mathbb{R}$; it is neither 1-1 nor onto $\mathbb{R}$. Then
\[
f(f^{-1}([-1,1])) = f([-1,1]) = [0,1] \neq [-1,1], \qquad f^{-1}(f([0,1])) = f^{-1}([0,1]) = [-1,1] \neq [0,1].
\]
In the first, the negative numbers are lost because they are not values of $f$; in the second, the points of $[-1,0)$ are gained because they have the same values as points of $(0,1]$.

\textbf{(iii) Bijections.} If $f$ is 1-1 and onto, then images also respect complements:
\[
f(A \setminus F) = B \setminus f(F) \qquad (F \subset A).
\]
Indeed, since $f$ is onto, $f(F) \cup f(A\setminus F) = f(A) = B$; since $f$ is 1-1, no point of $A \setminus F$ has the same value as a point of $F$, so $f(F) \cap f(A\setminus F) = \emptyset$. This will be used in Theorem \ref{thm:inverse}.
\end{remark}

\begin{theorem}\label{thm:closedbounded}
If $\mathbf{f}$ is a continuous mapping of a compact metric space $X$ into $\mathbb{R}^k$, then $\mathbf{f}(X)$ is closed and bounded. Thus, $\mathbf{f}$ is bounded.
\end{theorem}

\begin{proof}
\textbf{Step 1: $\mathbf{f}(X)$ is compact.} By Theorem \ref{thm:compactimage}, $\mathbf{f}(X)$ is a compact subset of $\mathbb{R}^k$.

\textbf{Step 2: Compact subsets are closed and bounded.} By Theorem 1.70, a compact subset of $\mathbb{R}^k$ is closed and bounded. Hence $\mathbf{f}(X)$ is closed and bounded. We emphasize that here \textit{we have used the fact that a compact set is bounded} (and closed). Recall from Section 1 that this direction holds in every metric space (Theorem 1.63 and Remark 1.71), while the converse ``closed and bounded implies compact'' may fail in a general metric space; in $\mathbb{R}^k$ the two are equivalent (Theorem 1.70).

\textbf{Step 3: $\mathbf{f}$ is bounded.} Since $\mathbf{f}(X)$ is bounded, there are a point $\mathbf{q} \in \mathbb{R}^k$ and a number $R$ with $|\mathbf{y} - \mathbf{q}| < R$ for all $\mathbf{y} \in \mathbf{f}(X)$ (Definition 1.33(i)). Then for every $x \in X$,
\[
|\mathbf{f}(x)| \leq |\mathbf{f}(x) - \mathbf{q}| + |\mathbf{q}| < R + |\mathbf{q}|,
\]
so $\mathbf{f}$ is bounded in the sense of Definition \ref{def:bounded}, with $M = R + |\mathbf{q}|$.
\end{proof}

\begin{theorem}\label{thm:maxmin}
Suppose $f$ is a continuous real function on a (nonempty) compact metric space $X$, and
\[
M = \sup_{p \in X} f(p), \qquad m = \inf_{p \in X} f(p).
\]
Then there exist points $p, q \in X$ such that $f(p) = M$ and $f(q) = m$.

The notation means that $M$ is the least upper bound of the set of all numbers $f(p)$, where $p$ ranges over $X$, and that $m$ is the greatest lower bound of this set of numbers. The conclusion may also be stated as follows: There exist points $p$ and $q$ in $X$ such that $f(q) \leq f(x) \leq f(p)$ for all $x \in X$; that is, $f$ attains its maximum (at $p$) and its minimum (at $q$).
\end{theorem}

\begin{proof}
\textbf{Step 1: $f(X)$ is a nonempty, closed and bounded set of real numbers.} This is Theorem \ref{thm:closedbounded} with $k = 1$. Being bounded, $f(X)$ is bounded above and below, so by the least upper bound property $M = \sup f(X)$ and $m = \inf f(X)$ exist as real numbers.

\textbf{Step 2: $M \in f(X)$.} By Theorem 1.50, $M = \sup f(X)$ belongs to the closure $\overline{f(X)}$. Since $f(X)$ is closed, $\overline{f(X)} = f(X)$ (Theorem 1.49(b)). Hence $M \in f(X)$, that is, $M = f(p)$ for some $p \in X$.

\textbf{Step 3: $m \in f(X)$.} The same argument as in Theorem 1.50, with the inequalities reversed, shows that $m \in \overline{f(X)}$: if $m \notin f(X)$, then for every $h > 0$ the number $m + h$ is not a lower bound of $f(X)$ (because $m$ is the \textit{greatest} lower bound), so there is $y \in f(X)$ with $m < y < m + h$; thus every neighborhood of $m$ contains a point of $f(X)$ other than $m$, and $m$ is a limit point of $f(X)$. In either case $m \in \overline{f(X)} = f(X)$, so $m = f(q)$ for some $q \in X$.
\end{proof}

\begin{remark}\label{rem:supnotlimit}
Do not be misled by the proof into thinking that the maximum value $M$ must be a limit point of $f(X)$. The proof of Theorem 1.50 shows that $\sup E$ is \textit{either} a point of $E$ \textit{or} a limit point of $E$, and it may well be an isolated point of $E$. For example, let $X = [0,1] \cup \{2\}$, which is closed and bounded in $\mathbb{R}$, hence compact (Theorem 1.70), and let $f(x) = x$. Then $f(X) = [0,1] \cup \{2\}$ and the maximum $M = 2 = f(2)$ is an isolated point of $f(X)$; it is not approached by other values of $f$. Similarly, if $f = 0$ on $[0,1]$ and $f(2) = 5$, the maximum $5$ is attained only at the isolated point $2$, and is not a limit point of $f(X) = \{0, 5\}$.
\end{remark}

\begin{theorem}\label{thm:inverse}
Suppose $f$ is a continuous 1-1 mapping of a compact metric space $X$ onto a metric space $Y$. Then the inverse mapping $f^{-1}$ defined on $Y$ by
\[
f^{-1}(f(x)) = x \qquad (x \in X)
\]
is a continuous mapping of $Y$ onto $X$.
\end{theorem}

\begin{proof}
\textbf{Step 1: What has to be shown.}\\
Applying Theorem \ref{thm:openpre} to the mapping $f^{-1} : Y \to X$ in place of $f$, it suffices to prove that $(f^{-1})^{-1}(V)$ is open in $Y$ for every open set $V$ in $X$. Now the inverse image of $V$ under the map $f^{-1}$ is
\[
(f^{-1})^{-1}(V) = \{y \in Y : f^{-1}(y) \in V\} = f(V),
\]
because $f$ is a bijection: $f^{-1}(y) \in V$ means exactly that $y = f(x)$ for some $x \in V$. So it suffices to prove that \textit{$f(V)$ is open in $Y$ for every open set $V$ in $X$}.

\textbf{Step 2: Pass to complements.}\\
Fix an open set $V$ in $X$. Its complement $V^c = X \setminus V$ is closed in $X$ (Theorem 1.45).

\textbf{Step 3: $V^c$ is compact.}\\
$V^c$ is a closed subset of the compact space $X$, hence compact (Theorem 1.64).

\textbf{Step 4: $f(V^c)$ is compact, hence closed.}\\
The restriction of $f$ to $V^c$ is continuous (Remark \ref{rem:local}(a)), and $V^c$ is a compact metric space in its own right (Theorem 1.60). By Theorem \ref{thm:compactimage}, $f(V^c)$ is a compact subset of $Y$, and therefore $f(V^c)$ is closed in $Y$ (Theorem 1.63). (If $V^c$ is empty, $f(V^c)$ is empty, which is closed.)

\textbf{Step 5: Back to $f(V)$.}\\
Since $f$ is 1-1 and onto, $f(V)$ is the complement of $f(V^c)$ in $Y$ (Remark \ref{rem:imagepre}(iii)). Since $f(V^c)$ is closed, $f(V)$ is open (Corollary 1.45.1). By Step 1, $f^{-1}$ is continuous.
\end{proof}

\begin{remark}\label{rem:homeo}
Thus, if the domain is compact, a continuous bijection automatically has a continuous inverse: it is a homeomorphism in the sense of Remark 1.12, and we get the continuity of the inverse for free. The compactness of the domain is essential for this luxury. In Example \ref{ex:circle} we shall see a continuous 1-1 mapping of a non-compact space onto a compact one whose inverse is not continuous.
\end{remark}

\begin{definition}\label{def:unif}
Let $f$ be a mapping of a metric space $X$ into a metric space $Y$. We say that $f$ is \textit{uniformly continuous} on $X$ if for every $\varepsilon > 0$ there exists $\delta > 0$ such that
\[
d_Y(f(p), f(q)) < \varepsilon
\]
for all $p$ and $q$ in $X$ for which $d_X(p, q) < \delta$.
\end{definition}

\begin{remark}\label{rem:unifvscont}
Let us consider the differences between the concepts of continuity and of uniform continuity.
\begin{enumerate}
    \item \textbf{A global property.} Uniform continuity is a property of a function on a \textit{set}, whereas continuity can be defined at a single point. To ask whether a given function is uniformly continuous at a certain point is meaningless. Uniform continuity describes the behaviour of the function over the whole domain at once.

    \item \textbf{Who chooses $\delta$.} If $f$ is continuous on $X$, then for each $\varepsilon > 0$ and for each point $p$ of $X$ we can find a $\delta > 0$ having the property of Definition \ref{def:cont}. This $\delta$ depends on $\varepsilon$ \textit{and on $p$}. If $f$ is uniformly continuous on $X$, then for each $\varepsilon > 0$ it is possible to find \textit{one} number $\delta > 0$ which will do for \textit{all} points $p$ of $X$. In terms of quantifiers:
    Continuity on $X$:
    \[
    \forall p \in X\ \ \forall \varepsilon > 0\ \ \exists \delta > 0\ \ \forall q \in X: \quad d_X(p,q) < \delta \Rightarrow d_Y(f(p), f(q)) < \varepsilon.
    \]
    Uniform continuity on $X$:
    \[
    \forall \varepsilon > 0\ \ \exists \delta > 0\ \ \forall p \in X\ \ \forall q \in X: \quad d_X(p,q) < \delta \Rightarrow d_Y(f(p), f(q)) < \varepsilon.
    \]
    The only difference is the position of ``$\forall p$''; in the second, $\delta$ is chosen before $p$, so it cannot depend on $p$.

    \item Evidently, every uniformly continuous function is continuous: given $p$ and $\varepsilon$, the uniform $\delta$ works at $p$.

    \item \textbf{An illustration.} On $\mathbb{R}$, $f(x) = x$ is uniformly continuous (take $\delta = \varepsilon$). But $f(x) = x^2$ is continuous and not uniformly continuous on $\mathbb{R}$: since $|x^2 - p^2| = |x-p|\,|x+p|$, the $\delta$ that works at $p$ has to shrink as $|p|$ grows. Precisely, for $\varepsilon = 1$ and any $\delta > 0$, let $p = 1/\delta$ and $q = p + \delta/2$; then $|q - p| < \delta$ but $|q^2 - p^2| = \frac{\delta}{2}(2p + \frac{\delta}{2}) > \frac{\delta}{2} \cdot \frac{2}{\delta} = 1$. So no single $\delta$ serves the whole line.
\end{enumerate}
\end{remark}

\begin{theorem}\label{thm:unifcompact}
Let $f$ be a continuous mapping of a compact metric space $X$ into a metric space $Y$. Then $f$ is uniformly continuous on $X$.
\end{theorem}

\textbf{The idea of the proof.} Continuity gives, for each point $p$, its own $\delta_p$. There may be infinitely many points, and the numbers $\delta_p$ may have infimum $0$, so we cannot simply take the smallest. Compactness rescues us: finitely many neighborhoods cover $X$, and the minimum of \textit{finitely many} positive numbers is positive. The radii are halved so that if $p$ lies in one of these finitely many neighborhoods and $q$ is close to $p$, then \textit{both} $p$ and $q$ lie within $\delta_{p_m}$ of the same centre $p_m$; and $\varepsilon/2$ is used because we travel from $f(p)$ to $f(q)$ via $f(p_m)$, in two steps.

\begin{proof}
Let $\varepsilon > 0$ be given.

\textbf{Step 1: A $\delta_p$ for each point.}\\
Since $f$ is continuous, we can associate to each point $p \in X$ a positive number $\delta_p$ such that
\[
q \in X, \ d_X(p,q) < \delta_p \quad \text{implies} \quad d_Y(f(p), f(q)) < \frac{\varepsilon}{2}. \tag{1}
\]

\textbf{Step 2: An open cover by half-size neighborhoods.}\\
Let $J(p)$ be the set of all $q \in X$ for which $d_X(p,q) < \frac{1}{2}\delta_p$; that is, $J(p) = N_{\delta_p/2}(p)$. Each $J(p)$ is open (Theorem 1.42), and since $p \in J(p)$, the collection of all sets $J(p)$, $p \in X$, is an open cover of $X$.

\textbf{Step 3: A finite subcover and the choice of $\delta$.}\\
Since $X$ is compact, there is a finite set of points $p_1, \dots, p_n$ in $X$ such that
\[
X \subset J(p_1) \cup \cdots \cup J(p_n).
\]
We put
\[
\delta = \min\Big( \frac{\delta_{p_1}}{2}, \dots, \frac{\delta_{p_n}}{2} \Big).
\]
Then $\delta > 0$, being the minimum of finitely many positive numbers. (This is where compactness is used: an infimum over infinitely many points could be $0$.)

\textbf{Step 4: The $\delta$ works for all pairs.}\\
Now let $p$ and $q$ be points of $X$ such that $d_X(p,q) < \delta$. Since the $J(p_i)$ cover $X$, there is an integer $m$, $1 \leq m \leq n$, such that $p \in J(p_m)$; hence
\[
d_X(p, p_m) < \frac{1}{2}\delta_{p_m} < \delta_{p_m},
\]
and we also have
\[
d_X(q, p_m) \leq d_X(q, p) + d_X(p, p_m) < \delta + \frac{1}{2}\delta_{p_m} \leq \frac{1}{2}\delta_{p_m} + \frac{1}{2}\delta_{p_m} = \delta_{p_m}.
\]
So both $p$ and $q$ lie within $\delta_{p_m}$ of $p_m$. Applying (1) at the point $p_m$ twice,
\[
d_Y(f(p), f(q)) \leq d_Y(f(p), f(p_m)) + d_Y(f(p_m), f(q)) < \frac{\varepsilon}{2} + \frac{\varepsilon}{2} = \varepsilon.
\]
Since $\delta$ depends only on $\varepsilon$ (not on $p$ or $q$), $f$ is uniformly continuous on $X$.
\end{proof}

The next theorem shows that compactness is essential in the hypotheses of Theorems \ref{thm:closedbounded}, \ref{thm:maxmin} and \ref{thm:unifcompact}.

\begin{theorem}\label{thm:noncompact}
Let $E$ be a noncompact set in $\mathbb{R}^1$. Then
\begin{enumerate}[label=(\alph*)]
    \item there exists a continuous function on $E$ which is not bounded;
    \item there exists a continuous and bounded function on $E$ which has no maximum.
\end{enumerate}
If, in addition, $E$ is bounded, then
\begin{enumerate}[label=(\alph*)]
    \setcounter{enumi}{2}
    \item there exists a continuous function on $E$ which is not uniformly continuous.
\end{enumerate}
\end{theorem}

\begin{proof}
\textbf{Case 1: $E$ is bounded.}

\textit{Step 1.1: A ``hole'' in $E$.} By Theorem 1.70, a closed and bounded subset of $\mathbb{R}^1$ is compact. Since $E$ is bounded but not compact, $E$ is not closed. Hence there exists a limit point $x_0$ of $E$ which is not a point of $E$.

\textit{Step 1.2: Proof of (a).} Consider
\[
f(x) = \frac{1}{x - x_0} \qquad (x \in E).
\]
This is a rational function whose denominator does not vanish on $E$ (since $x_0 \notin E$), so $f$ is continuous on $E$ (Example \ref{ex:polynomials}, Step 7, and Remark \ref{rem:local}(a)). It is unbounded: given $M > 0$, since $x_0$ is a limit point of $E$, there is $x \in E$ with $0 < |x - x_0| < 1/M$, and then $|f(x)| > M$.

\textit{Step 1.3: Proof of (c).} We show that $f$ is not uniformly continuous on $E$; we show that $\varepsilon = 1$ admits no $\delta$. Let $\delta > 0$ be arbitrary. Choose $x \in E$ with $|x - x_0| < \delta/2$ (possible since $x_0$ is a limit point of $E$). Now $|f(t)|$ becomes as large as we please when $t \in E$ is taken close to $x_0$; so choose $t \in E$ with
\[
|t - x_0| < \min\Big( \frac{\delta}{2}, \frac{1}{|f(x)| + 1} \Big).
\]
Then $|t - x| \leq |t - x_0| + |x_0 - x| < \delta$, while $|f(t)| = 1/|t - x_0| > |f(x)| + 1$, so that
\[
|f(t) - f(x)| \geq |f(t)| - |f(x)| > 1.
\]
Thus for every $\delta > 0$ there are points $x, t \in E$ with $|t - x| < \delta$ but $|f(t) - f(x)| > 1$. So $f$ is not uniformly continuous on $E$.

\textit{Step 1.4: Proof of (b).} The function
\[
g(x) = \frac{1}{1 + (x - x_0)^2} \qquad (x \in E)
\]
is continuous on $E$ (a rational function whose denominator is never zero), and it is bounded, since $0 < g(x) \leq 1$. Moreover $g(x) < 1$ for every $x \in E$, because $x \neq x_0$. On the other hand, $\sup_{x\in E} g(x) = 1$: given $0 < \eta < 1$, choose $x \in E$ with $(x - x_0)^2 < \eta$ (again possible since $x_0$ is a limit point of $E$); then, since $(1-\eta)(1+\eta) = 1 - \eta^2 < 1$,
\[
g(x) = \frac{1}{1 + (x-x_0)^2} > \frac{1}{1 + \eta} > 1 - \eta.
\]
So the least upper bound $1$ is not attained: $g$ has no maximum on $E$.

\textbf{Case 2: $E$ is unbounded.}

\textit{Proof of (a).} Let $f(x) = x$ on $E$. It is continuous. Since $E$ is unbounded, for every $M$ there is $x \in E$ with $|x| > M$ (otherwise $E \subset [-M, M]$ would be bounded). So $f$ is not bounded.

\textit{Proof of (b).} Let
\[
h(x) = \frac{x^2}{1 + x^2} = 1 - \frac{1}{1 + x^2} \qquad (x \in E).
\]
$h$ is continuous and $0 \leq h(x) < 1$ for all $x$. Given $0 < \eta < 1$, choose $x \in E$ with $x^2 > 1/\eta$; then $h(x) > 1 - \eta$. Hence $\sup_{x \in E} h(x) = 1$, and it is not attained, since $h(x) < 1$ for all $x \in E$. So $h$ is a bounded continuous function with no maximum.

\textbf{(c) fails without boundedness.} Assertion (c) would be false if boundedness were omitted from the hypotheses. For, let $E = \mathbb{N}$, which is closed and unbounded, hence not compact. Every function defined on $\mathbb{N}$ is uniformly continuous on $\mathbb{N}$; see Example \ref{ex:N} below.
\end{proof}

\begin{remark}\label{rem:otherfns}
The particular functions used in the proof are not sacred; it is instructive to see which other functions work and which do not.
\begin{enumerate}
    \item \textbf{Bounded case, (b).} Instead of $g$ we may use
    \[
    \tilde g(x) = \frac{1}{1 + |x - x_0|} \qquad (x \in E).
    \]
    It is continuous: $x \mapsto |x - x_0|$ is continuous (Example \ref{ex:polynomials}, Step 9), and $\tilde g$ is its composition with $s \mapsto 1/(1+s)$, which is continuous for $s \geq 0$ (Theorem \ref{thm:composition}). We have $0 < \tilde g(x) < 1$ on $E$ (because $x \neq x_0$), and $\tilde g(x) > 1/(1+\eta) > 1 - \eta$ as soon as $|x - x_0| < \eta$. Hence $\sup_E \tilde g = 1$ is not attained.

    \item \textbf{Unbounded case, (b).} Instead of $h$ we may use
    \[
    \tilde h(x) = \frac{|x|}{1 + |x|} = 1 - \frac{1}{1 + |x|} \qquad (x \in E).
    \]
    It is continuous, $0 \leq \tilde h < 1$, and $\tilde h(x) > 1 - \eta$ as soon as $|x| > 1/\eta$. Since $E$ is unbounded, $|x|$ takes arbitrarily large values on $E$, whether $E$ is unbounded above or unbounded below. Hence $\sup_E \tilde h = 1$ is not attained.

    \item \textbf{What goes wrong with $x/(1+|x|)$.} Consider
    \[
    k(x) = \frac{x}{1 + |x|} \qquad (x \in \mathbb{R}).
    \]
    It is continuous and $-1 < k(x) < 1$. It is strictly increasing: for $x \geq 0$, $k(x) = 1 - \frac{1}{1+x}$ increases from $0$ towards $1$; for $x \leq 0$, $k(x) = \frac{x}{1-x} = -1 + \frac{1}{1-x}$ increases from values near $-1$ up to $0$; and $k < 0$ on $(-\infty,0)$ while $k \geq 0$ on $[0,\infty)$. If $E$ is unbounded \textit{above}, then $\sup_E k = 1$ is not attained, and $k$ works. But suppose $E$ is unbounded only \textit{below}, for instance $E = (-\infty, \alpha]$, which is closed but unbounded, hence not compact. Since $k$ is increasing, it attains its maximum on $E$ at the right end point: $\max_E k = k(\alpha)$; for $E = (-\infty, 0]$ the maximum $0$ is attained at $0$. So $k$ does \textit{not} establish (b) for such $E$. The modulus in $\tilde h$ (or the square in $h$) cures this: it makes the function approach its supremum $1$ at \textit{both} ends of the line.
\end{enumerate}
\end{remark}

\begin{example}\label{ex:01}
Let $E = (0,1)$. It is bounded but not closed: the points $0$ and $1$ are limit points of $E$ not in $E$. So Theorem \ref{thm:noncompact} applies, and either of these points may be used as $x_0$.

\textbf{With $x_0 = 0$:} $f(x) = 1/x$ is continuous and unbounded on $(0,1)$; $g(x) = 1/(1+x^2)$ (or $\tilde g(x) = 1/(1+x)$) is continuous and bounded with supremum $1$ not attained. To see directly that $1/x$ is not uniformly continuous, take $x_n = 1/n$ and $y_n = 1/(n+1)$ ($n \geq 2$). Then
\[
|x_n - y_n| = \frac{1}{n(n+1)} \to 0, \qquad |f(x_n) - f(y_n)| = |n - (n+1)| = 1.
\]
So for $\varepsilon = 1$ (or any $\varepsilon \leq 1$) and any $\delta > 0$, choosing $n$ with $\frac{1}{n(n+1)} < \delta$ gives two points closer than $\delta$ whose values differ by $1$; no $\delta$ works.

\textbf{With $x_0 = 1$:} $f(x) = 1/(x-1)$ is continuous and unbounded on $(0,1)$ (it tends to $-\infty$ near $1$); $g(x) = 1/(1 + (x-1)^2)$ is bounded with supremum $1$ not attained; and with $x_n = 1 - \frac{1}{n}$, $y_n = 1 - \frac{1}{n+1}$ we get $|x_n - y_n| = \frac{1}{n(n+1)} \to 0$ while $|f(x_n) - f(y_n)| = |-n + (n+1)| = 1$, so $1/(x-1)$ is not uniformly continuous on $(0,1)$.
\end{example}

\begin{example}\label{ex:Q}
In the bounded case of Theorem \ref{thm:noncompact}, boundedness was used \textit{only} to guarantee that $E$ is not closed, so that there is a limit point $x_0$ of $E$ outside $E$. The function $1/(x - x_0)$ then misbehaves near this ``hole''. Analysing this, we can produce such a function even on an unbounded set, as long as it has a hole.

Let $E = \mathbb{Q}$, which is unbounded. The irrational number $\pi$ is a limit point of $\mathbb{Q}$ (every interval around $\pi$ contains rationals, since $\mathbb{Q}$ is dense in $\mathbb{R}$; Example 1.51), and $\pi \notin \mathbb{Q}$. Let
\[
f(x) = \frac{1}{x - \pi} \qquad (x \in \mathbb{Q}).
\]
\textbf{$f$ is continuous on $\mathbb{Q}$:} it is the restriction to $\mathbb{Q}$ of the rational function $1/(x - \pi)$, which is continuous on $\mathbb{R} \setminus \{\pi\} \supset \mathbb{Q}$ (Example \ref{ex:polynomials}, Remark \ref{rem:local}(a)).

\textbf{$f$ is not uniformly continuous on $\mathbb{Q}$:} For each $n$, choose a rational $q_n$ with $\pi < q_n < \pi + \frac{1}{n}$, and then a rational $p_n$ with $\pi < p_n < \pi + \frac{q_n - \pi}{2}$ (both possible by the density of $\mathbb{Q}$). Then $0 < p_n - \pi < \frac{1}{2}(q_n - \pi)$, so
\[
f(p_n) = \frac{1}{p_n - \pi} > \frac{2}{q_n - \pi} = 2 f(q_n),
\]
and hence
\[
|f(p_n) - f(q_n)| > f(q_n) = \frac{1}{q_n - \pi} > n \geq 1, \qquad\text{while}\qquad |p_n - q_n| < q_n - \pi < \frac{1}{n}.
\]
So for $\varepsilon = 1$ no $\delta$ works: choose $n$ with $1/n < \delta$. Thus $f$ is continuous but not uniformly continuous on $\mathbb{Q}$, though $\mathbb{Q}$ is unbounded. (It is also unbounded.) The lesson: in Theorem \ref{thm:noncompact}(c) what really matters is the presence of a limit point outside $E$, not the boundedness of $E$ as such.
\end{example}

\begin{example}\label{ex:N}
\textbf{Every function on $\mathbb{N}$ is uniformly continuous.} Let $f$ be any function from $\mathbb{N}$ (with the usual distance $|m - n|$) into a metric space $Y$. Every point of $\mathbb{N}$ is isolated, so $f$ is automatically continuous at every point (Remark \ref{rem:contcompare}). But more is true: $f$ is uniformly continuous. Given $\varepsilon > 0$, take $\delta = \frac{1}{2}$ (any $\delta \leq 1$ will do). If $m, n \in \mathbb{N}$ and $|m - n| < \frac{1}{2}$, then $m = n$, since distinct integers are at distance at least $1$; hence $d_Y(f(m), f(n)) = 0 < \varepsilon$. The point is that the points of $\mathbb{N}$ are \textit{uniformly} isolated: every point has an isolating neighborhood of the \textit{same} radius $\frac{1}{2}$, and this one radius serves as a $\delta$ for all points at once.

\textbf{Discrete metric spaces.} The same argument applies to the \textit{discrete metric} on any set $X$: $d(x,y) = 1$ if $x \neq y$ and $d(x,x) = 0$. (This is a metric: (a) and (b) of Definition 1.26 are clear, and in (c), if $x \neq y$ then $r$ differs from at least one of $x$, $y$, so $d(x,r) + d(r,y) \geq 1 = d(x,y)$.) Every function $f$ from a discrete metric space $X$ into any metric space $Y$ is uniformly continuous: take $\delta = \frac{1}{2}$; then $d(x,y) < \frac12$ forces $x = y$ and $d_Y(f(x), f(y)) = 0 < \varepsilon$.
\end{example}

\begin{example}\label{ex:1overn}
\textbf{All points isolated does not imply uniform continuity.} Every function on a domain consisting only of isolated points is continuous, but it need not be uniformly continuous. Let
\[
E = \Big\{ \frac{1}{n} : n \in \mathbb{N} \Big\} \subset \mathbb{R}, \qquad f(x) = \frac{1}{x} \quad (x \in E), \ \text{ i.e., } f\Big(\frac{1}{n}\Big) = n.
\]
Each point $\frac{1}{n}$ is isolated in $E$: its nearest neighbour in $E$ is $\frac{1}{n+1}$, at distance $\frac{1}{n(n+1)}$, so the neighborhood of $\frac1n$ of radius $\frac{1}{n(n+1)}$ contains no other point of $E$. Hence $f$ is continuous on $E$.

But $f$ is not uniformly continuous on $E$. Let $0 < \varepsilon < 1$ (indeed $\varepsilon = 1$ will also do), and let $\delta > 0$ be arbitrary. Choose $n$ with $\frac{1}{n(n+1)} < \delta$, and take $x = \frac{1}{n}$, $y = \frac{1}{n+1}$. Then
\[
|x - y| = \frac{1}{n(n+1)} < \delta, \qquad |f(x) - f(y)| = |n - (n+1)| = 1 \geq \varepsilon.
\]
So uniform continuity is lost for every $\varepsilon \leq 1$. (Taking $x = \frac1n$, $y = \frac{1}{2n}$ instead, we have $|x - y| = \frac{1}{2n} \to 0$ while $|f(x) - f(y)| = n \to \infty$; so in fact no $\varepsilon > 0$ at all admits a $\delta$.)

Compare with $\mathbb{N}$: in $E$ the isolating radii $\frac{1}{n(n+1)}$ shrink to $0$, so no single $\delta$ can serve all points. The points crowd towards $0$, which is a limit point of $E$ not belonging to $E$; once again it is a ``hole'' in the domain that destroys uniform continuity, exactly as in Theorem \ref{thm:noncompact}(c) and Example \ref{ex:Q}.
\end{example}

\begin{example}\label{ex:circle}
\textbf{A continuous bijection whose inverse is not continuous.} As promised in Remark \ref{rem:homeo}, we show that compactness cannot be dropped in Theorem \ref{thm:inverse}. In this example we assume the familiar properties of the functions $\cos$ and $\sin$: they are continuous on $\mathbb{R}$ (see Remark \ref{rem:sincos}), $\cos^2 t + \sin^2 t = 1$, $\cos$ is even and $\sin$ is odd, both have period $2\pi$, $\sin t = 0$ exactly when $t$ is an integer multiple of $\pi$, $\cos$ decreases strictly from $1$ to $-1$ on $[0,\pi]$ and increases strictly from $-1$ to $1$ on $[\pi, 2\pi]$, $\sin \geq 0$ on $[0,\pi]$ and $\sin \leq 0$ on $[\pi, 2\pi]$, and
\[
\cos t_1 = \cos t_2 \iff t_2 = t_1 + 2n\pi \ \text{ or } \ t_2 = -t_1 + 2n\pi \ \text{ for some integer } n.
\]
(The last fact follows from $\cos t_1 - \cos t_2 = -2 \sin\frac{t_1 + t_2}{2} \sin \frac{t_1 - t_2}{2}$ and the zeros of $\sin$.)

Let $X = [0, 2\pi) \subset \mathbb{R}$, let $Y = \{(x,y) \in \mathbb{R}^2 : x^2 + y^2 = 1\}$ be the unit circle, and define
\[
\mathbf{f}(t) = (\cos t, \sin t) \qquad (0 \leq t < 2\pi).
\]
Since $\cos^2 t + \sin^2 t = 1$, $\mathbf{f}$ maps $X$ into $Y$.

\textbf{Step 1: $\mathbf{f}$ is continuous.} Its components $\cos t$ and $\sin t$ are continuous, so $\mathbf{f}$ is continuous by Theorem \ref{thm:components}(a) (restricted to $X$, Remark \ref{rem:local}(a)).

\textbf{Step 2: $\mathbf{f}$ is 1-1.} Let $t_1, t_2 \in [0,2\pi)$ with $\mathbf{f}(t_1) = \mathbf{f}(t_2)$, i.e., $\cos t_1 = \cos t_2$ and $\sin t_1 = \sin t_2$. From $\cos t_1 = \cos t_2$, for some integer $n$, either $t_2 = t_1 + 2n\pi$ or $t_2 = -t_1 + 2n\pi$.
\begin{itemize}
    \item \textit{Case $t_2 = t_1 + 2n\pi$.} Since $t_1, t_2 \in [0, 2\pi)$, we have $|t_2 - t_1| < 2\pi$, i.e., $|2n\pi| < 2\pi$, so $n = 0$ and $t_1 = t_2$.
    \item \textit{Case $t_2 = -t_1 + 2n\pi$.} Now we use the second coordinate. By periodicity and oddness, $\sin t_2 = \sin(-t_1) = -\sin t_1$. Together with $\sin t_2 = \sin t_1$, this gives $\sin t_1 = 0$, so $t_1 \in \{0, \pi\}$ (the multiples of $\pi$ in $[0, 2\pi)$). If $t_1 = 0$, then $t_2 = 2n\pi \in [0, 2\pi)$ forces $t_2 = 0 = t_1$. If $t_1 = \pi$, then $t_2 = (2n-1)\pi \in [0,2\pi)$ forces $t_2 = \pi = t_1$.
\end{itemize}
In all cases $t_1 = t_2$. So $\mathbf{f}$ is 1-1.

\textbf{Step 3: $\mathbf{f}$ is onto $Y$.} Let $(x,y) \in Y$, so $x^2 + y^2 = 1$. Then $x^2 \leq 1$, i.e., $-1 \leq x \leq 1$, and $y = \sqrt{1 - x^2}$ or $y = -\sqrt{1-x^2}$. Thus $(x, \sqrt{1-x^2})$ is a point of the upper half of the circle and $(x, -\sqrt{1-x^2})$ a point of the lower half.
\begin{itemize}
    \item \textit{Upper half, $y \geq 0$.} On $[0,\pi]$, $\cos$ is continuous and strictly decreasing from $\cos 0 = 1$ to $\cos\pi = -1$; hence it takes every value in $[-1,1]$ exactly once (by the intermediate value theorem, Theorem \ref{thm:ivt} below). So there is a unique $t \in [0,\pi]$ with $\cos t = x$. Since $\sin t \geq 0$ on $[0,\pi]$, $\sin t = \sqrt{1 - \cos^2 t} = \sqrt{1 - x^2} = y$. Thus $\mathbf{f}(t) = (x,y)$; indeed $\mathbf{f}$ is a bijection of $[0,\pi]$ onto the closed upper half circle.
    \item \textit{Lower half, $y < 0$.} Then $x^2 = 1 - y^2 < 1$, so $-1 < x < 1$. On $[\pi, 2\pi]$, $\cos$ increases strictly from $-1$ to $1$, so there is a unique $t \in [\pi, 2\pi]$ with $\cos t = x$, and $t \neq \pi, 2\pi$ because $x \neq \pm 1$; thus $t \in (\pi, 2\pi) \subset X$. Since $\sin t \leq 0$ there, $\sin t = -\sqrt{1 - x^2} = y$. Thus $\mathbf{f}(t) = (x,y)$; $\mathbf{f}$ is a bijection of $(\pi, 2\pi)$ onto the open lower half circle.
\end{itemize}
So $\mathbf{f}$ is onto $Y$. Hence $\mathbf{f}$ is a continuous 1-1 mapping of $X$ onto $Y$.

\textbf{Step 4: The inverse is not continuous.} Let $\mathbf{g} = \mathbf{f}^{-1} : Y \to [0, 2\pi)$. Consider the points $\mathbf{z}_n = \mathbf{f}(2\pi - \frac{1}{n}) \in Y$, $n = 1, 2, \dots$. By continuity of $\cos$ and $\sin$ at $2\pi$ (and Theorem \ref{thm:seqlimit}),
\[
\mathbf{z}_n = \Big(\cos\big(2\pi - \tfrac{1}{n}\big), \sin\big(2\pi - \tfrac{1}{n}\big)\Big) \to (\cos 2\pi, \sin 2\pi) = (1, 0) = \mathbf{f}(0),
\]
and $\mathbf{z}_n \neq (1,0)$ because $\mathbf{f}$ is 1-1. But
\[
\mathbf{g}(\mathbf{z}_n) = 2\pi - \frac{1}{n} \to 2\pi \neq 0 = \mathbf{g}\big((1,0)\big).
\]
By Theorems \ref{thm:seqlimit} and \ref{thm:contlimit}, $\mathbf{g}$ is not continuous at $(1,0)$. Geometrically, points of the circle just below $(1,0)$ are sent by $\mathbf{g}$ near $2\pi$, while $(1,0)$ itself is sent to $0$: the inverse tears the circle apart at $(1,0)$.

One can also see this without computation. $Y$ is closed in $\mathbb{R}^2$ (it is the inverse image of the closed set $\{1\}$ under the continuous function $x^2 + y^2$; Corollary \ref{cor:closedpre}) and bounded, hence compact (Theorem 1.70). If $\mathbf{g}$ were continuous, then $\mathbf{g}(Y) = [0,2\pi)$ would be compact (Theorem \ref{thm:compactimage}); but $[0, 2\pi)$ is not closed ($2\pi$ is a limit point not in it), so it is not compact.

\textbf{Conclusion.} Everything in Theorem \ref{thm:inverse} holds here except the compactness of the domain $X = [0, 2\pi)$, and the conclusion fails. So compactness is essential, in general, to obtain the continuity of the inverse.
\end{example}

\begin{remark}\label{rem:sincos}
For those who prefer not to assume the continuity of $\cos$ and $\sin$: one may \textit{define}
\[
E(z) = \sum_{n=0}^\infty \frac{z^n}{n!} \quad (z \in \mathbb{C}), \qquad \cos x = \operatorname{Re} E(ix), \quad \sin x = \operatorname{Im} E(ix) \quad (x \in \mathbb{R}),
\]
the series converging absolutely for every $z$. Using the product of absolutely convergent series one shows $E(z + w) = E(z)E(w)$. Then for $|h| \leq 1$,
\[
|E(z + h) - E(z)| = |E(z)|\,|E(h) - 1| \leq |E(z)| \sum_{n=1}^\infty \frac{|h|^n}{n!} \leq |E(z)|\, |h| \sum_{n=1}^\infty \frac{|h|^{n-1}}{(n-1)!} \leq e\,|E(z)|\,|h|,
\]
so $E$ is continuous at every $z$ (take $\delta = \min(1, \varepsilon/(e|E(z)| + 1))$). Since $x \mapsto ix$ is continuous, and $|\operatorname{Re} w| \leq |w|$, $|\operatorname{Im} w| \leq |w|$, we get $|\cos x - \cos y| \leq |E(ix) - E(iy)|$ and $|\sin x - \sin y| \leq |E(ix) - E(iy)|$, so $\cos$ and $\sin$ are continuous on $\mathbb{R}$. The remaining properties of $\cos$ and $\sin$ used above can also be derived from these definitions; we do not pursue this here.
\end{remark}

\subsection{Continuity and Connectedness}

\begin{theorem}\label{thm:connimage}
If $f$ is a continuous mapping of a metric space $X$ into a metric space $Y$, and if $E$ is a connected subset of $X$, then $f(E)$ is connected.
\end{theorem}

\textbf{The idea of the proof.} We use nothing but the definition of connectedness (Section 1): a set is connected if it is not the union of two nonempty separated sets. Suppose $f(E)$ were \textit{not} connected, so that $f(E) = A \cup B$ with $A$, $B$ nonempty and separated. Pull this splitting back to $E$: let $G$ be the part of $E$ mapped into $A$ and $H$ the part mapped into $B$. Continuity (through Corollary \ref{cor:closedpre}) will show that $G$ and $H$ are separated as well. Then $E$ would be split into two nonempty separated sets, contradicting the connectedness of $E$.

\begin{proof}
\textbf{Step 1: Assume the contrary.}\\
Assume, on the contrary, that $f(E) = A \cup B$, where $A$ and $B$ are nonempty separated subsets of $Y$; that is, $\overline{A} \cap B = \emptyset$ and $A \cap \overline{B} = \emptyset$.

\textbf{Step 2: Pull back the splitting.}\\
Put
\[
G = E \cap f^{-1}(A), \qquad H = E \cap f^{-1}(B).
\]
Then $E = G \cup H$: every $x \in E$ has $f(x) \in f(E) = A \cup B$, so $x \in G$ or $x \in H$. Neither $G$ nor $H$ is empty: since $A \neq \emptyset$ and $A \subset f(E)$, there is $x \in E$ with $f(x) \in A$, i.e., $x \in G$; similarly $H \neq \emptyset$. Also, $f(H) = B$: clearly $f(H) \subset B$, and every $b \in B \subset f(E)$ is $f(x)$ for some $x \in E$, and this $x$ lies in $H$.

\textbf{Step 3: $\overline{G} \cap H = \emptyset$.}\\
Since $A \subset \overline{A}$, we have $G \subset f^{-1}(\overline{A})$. The set $\overline{A}$ is closed (Theorem 1.49(a)), and $f$ is continuous, so $f^{-1}(\overline{A})$ is closed in $X$ (Corollary \ref{cor:closedpre}). Since $\overline{G}$ is the smallest closed set containing $G$ (Theorem 1.49(c)), we get
\[
\overline{G} \subset f^{-1}(\overline{A}), \qquad\text{hence}\qquad f(\overline{G}) \subset \overline{A}.
\]
Now if there were a point $x \in \overline{G} \cap H$, then $f(x) \in \overline{A}$ (because $x \in \overline{G}$) and $f(x) \in B$ (because $x \in H$), i.e., $f(x) \in \overline{A} \cap B = \emptyset$, which is impossible. Hence $\overline{G} \cap H = \emptyset$.

\textbf{Step 4: $G \cap \overline{H} = \emptyset$.}\\
The same argument with the roles of $A$, $G$ and $B$, $H$ interchanged (using $A \cap \overline{B} = \emptyset$) shows that $G \cap \overline{H}$ is empty.

\textbf{Step 5: Contradiction.}\\
Thus $G$ and $H$ are nonempty separated sets with $E = G \cup H$. This is impossible, since $E$ is connected. Hence $f(E)$ is connected.
\end{proof}

\begin{theorem}[Intermediate Value Theorem]\label{thm:ivt}
Let $f$ be a continuous real function on the interval $[a,b]$. If $f(a) < f(b)$ and if $c$ is a number such that $f(a) < c < f(b)$, then there exists a point $x \in (a,b)$ such that $f(x) = c$.

A similar result holds, of course, if $f(a) > f(b)$. Roughly speaking, a continuous real function assumes all intermediate values on an interval.
\end{theorem}

\begin{proof}
The proof uses, twice, the characterization of the connected subsets of $\mathbb{R}^1$ (Theorem 1.75): a set $E \subset \mathbb{R}^1$ is connected if and only if $x, y \in E$ and $x < z < y$ imply $z \in E$; that is, connected subsets of $\mathbb{R}^1$ are exactly the intervals (Remark 1.76).

\textbf{Step 1: $[a,b]$ is connected.} If $x, y \in [a,b]$ and $x < z < y$, then $a \leq x < z < y \leq b$, so $z \in [a,b]$. By Theorem 1.75, $[a,b]$ is connected.

\textbf{Step 2: $f([a,b])$ is connected.} By Theorem \ref{thm:connimage}, $f([a,b])$ is a connected subset of $\mathbb{R}^1$. Hence, again by Theorem 1.75, $f([a,b])$ has the interval property: whenever two numbers belong to it, every number between them belongs to it.

\textbf{Step 3: $c$ is a value of $f$.} The numbers $f(a)$ and $f(b)$ belong to $f([a,b])$, and $f(a) < c < f(b)$. By Step 2, $c \in f([a,b])$; that is, $c$ has a preimage: $c = f(x)$ for some $x \in [a,b]$.

\textbf{Step 4: $x$ is an interior point of the interval.} Since $f(a) \neq c$ and $f(b) \neq c$, we have $x \neq a$ and $x \neq b$. Thus $x \in (a,b)$.

If $f(a) > f(b)$ and $f(b) < c < f(a)$, the same argument applies with the roles of $f(a)$ and $f(b)$ interchanged in Step 3.
\end{proof}

\begin{remark}\label{rem:ivtconverse}
At first glance, it might seem that the intermediate value theorem has a converse: that a function which assumes all intermediate values must be continuous. This is not so.

\textbf{1. A swap.} On $[0,1]$ let $f(x) = x$, except that we swap the values at two points:
\[
f\big(\tfrac14\big) = \tfrac34, \qquad f\big(\tfrac34\big) = \tfrac14, \qquad f(x) = x \ \text{ otherwise}.
\]
Then $f$ is a bijection of $[0,1]$ onto $[0,1]$, so $f([0,1]) = [0,1]$ and every number $c$ between $f(0) = 0$ and $f(1) = 1$ is a value of $f$ at some point of $(0,1)$: the conclusion of the intermediate value theorem on $[0,1]$ holds. But $f$ is not continuous at $\frac14$: the points $x_n = \frac14 + \frac{1}{n}$ ($n \geq 5$) tend to $\frac14$ and are different from $\frac14$ and $\frac34$, so $f(x_n) = x_n \to \frac14 \neq \frac34 = f(\frac14)$ (Theorems \ref{thm:seqlimit} and \ref{thm:contlimit}). Similarly at $\frac34$. So taking all intermediate values on an interval does not force continuity.

One may object that this $f$ fails the intermediate value property on \textit{subintervals}: on $[0, \frac14]$ we have $f(0) = 0$, $f(\frac14) = \frac34$, but the value $c = \frac12$ is not taken on $(0, \frac14)$, where $f(x) = x < \frac14$.

\textbf{2. The stronger property still does not imply continuity.} Even if we require that \textit{for any two points $x_1 < x_2$ and any number $c$ between $f(x_1)$ and $f(x_2)$, there is a point $x \in (x_1, x_2)$ with $f(x) = c$}, the function need not be continuous. The function $f(x) = \sin(1/x)$ ($x \neq 0$), $f(0) = 0$, has this property but is discontinuous at $0$; see Example \ref{ex:rudin}(d).

\textbf{3. Derivatives.} In the future we shall learn that the derivative of a differentiable function always has the intermediate value property, although it need not be continuous. This is \textit{Darboux's theorem}.
\end{remark}

\subsection{Discontinuities}

If $x$ is a point in the domain of definition of the function $f$ at which $f$ is not continuous, we say that $f$ is \textit{discontinuous} at $x$, or that $f$ has a \textit{discontinuity} at $x$. If $f$ is defined on an interval or on a segment, it is customary to divide discontinuities into two types. Before giving this classification, we define the right-hand and the left-hand limits of $f$ at $x$, which we denote by $f(x+)$ and $f(x-)$, respectively.

\begin{definition}\label{def:onesided}
Let $f$ be defined on $(a,b)$. Consider any point $x$ such that $a \leq x < b$. We write
\[
f(x+) = q
\]
if $f(t_n) \to q$ as $n \to \infty$, for all sequences $\{t_n\}$ in $(x,b)$ such that $t_n \to x$. To obtain the definition of $f(x-)$, for $a < x \leq b$, we restrict ourselves to sequences $\{t_n\}$ in $(a,x)$.

$f(x+)$ is called the \textit{right-hand limit} and $f(x-)$ the \textit{left-hand limit} of $f$ at $x$. We also write $f(x+) = \lim_{t \to x+} f(t)$ and $f(x-) = \lim_{t \to x-} f(t)$.
\end{definition}

The definition uses sequences. We now show that, as for ordinary limits, there is an equivalent $\varepsilon$-$\delta$ formulation: we simply restrict the $\delta$-neighborhood $(x - \delta, x + \delta)$ of $x$ to its right half $(x, x+\delta)$, or to its left half $(x - \delta, x)$.

\begin{theorem}\label{thm:onesidedeps}
Let $f$ be defined on $(a,b)$ and $a \leq x < b$. Then $f(x+) = q$ if and only if for every $\varepsilon > 0$ there is $\delta > 0$ such that
\[
|f(t) - q| < \varepsilon \qquad \text{whenever } x < t < x + \delta \text{ and } t < b.
\]
Similarly, for $a < x \leq b$, $f(x-) = q$ if and only if for every $\varepsilon > 0$ there is $\delta > 0$ such that $|f(t) - q| < \varepsilon$ whenever $x - \delta < t < x$ and $t > a$.
\end{theorem}

\begin{proof}
Let $g$ be the restriction of $f$ to the set $E = (x,b)$. The point $x$ is a limit point of $E$: for every $r > 0$, the point $x + \frac{1}{2}\min(r, b - x)$ lies in $E \cap N_r(x)$ and differs from $x$.

\textbf{Step 1: $f(x+)$ is an ordinary limit of $g$.} The sequences $\{t_n\}$ in $E$ with $t_n \neq x$ and $t_n \to x$ are exactly the sequences in $(x,b)$ with $t_n \to x$ (no term of a sequence in $(x,b)$ equals $x$). So by Definition \ref{def:onesided} and Theorem \ref{thm:seqlimit} (applied to $g$ on $E$),
\[
f(x+) = q \iff \lim_{t \to x} g(t) = q.
\]

\textbf{Step 2: Unwrap the $\varepsilon$-$\delta$ definition.} By Definition \ref{def:limit}, $\lim_{t\to x} g(t) = q$ means: for every $\varepsilon > 0$ there is $\delta > 0$ such that $|g(t) - q| < \varepsilon$ for all $t \in E$ with $0 < |t - x| < \delta$. For $t \in E = (x,b)$, the condition $0 < |t - x| < \delta$ says exactly $x < t < x + \delta$, and $g(t) = f(t)$. This is the stated condition.

The left-hand case is identical, with $E = (a,x)$, where $0 < |t - x| < \delta$ means $x - \delta < t < x$.
\end{proof}

\begin{theorem}\label{thm:contonesided}
Let $f$ be defined on $(a,b)$ and $a < x < b$. Then:
\begin{enumerate}[label=(\alph*)]
    \item $\lim_{t \to x} f(t)$ exists if and only if $f(x+)$ and $f(x-)$ both exist and are equal; in that case $f(x+) = f(x-) = \lim_{t \to x} f(t)$.
    \item $f$ is continuous at $x$ if and only if
    \[
    f(x+) = f(x-) = f(x).
    \]
\end{enumerate}
\end{theorem}

\begin{proof}
We prove (b) carefully; (a) is proved in exactly the same way, with $0 < |t - x| < \delta$ in place of $|t - x| < \delta$ and with $q = \lim_{t\to x} f(t)$ in place of $f(x)$.

\textbf{($\Rightarrow$)} Suppose $f$ is continuous at $x$. Let $\varepsilon > 0$. There is $\delta > 0$ such that $|f(t) - f(x)| < \varepsilon$ whenever $t \in (a,b)$ and $|t - x| < \delta$. In particular this holds for $x < t < x + \delta$ (the right half) and for $x - \delta < t < x$ (the left half). By Theorem \ref{thm:onesidedeps}, $f(x+) = f(x)$ and $f(x-) = f(x)$. In words: a $\delta$-window around $x$ which works for continuity works, in particular, on each of its two halves.

\textbf{($\Leftarrow$)} Suppose $f(x+) = f(x-) = f(x)$. Let $\varepsilon > 0$. By Theorem \ref{thm:onesidedeps}, there is $\delta_1 > 0$ such that $|f(t) - f(x)| < \varepsilon$ for $x < t < x + \delta_1$, and there is $\delta_2 > 0$ such that $|f(t) - f(x)| < \varepsilon$ for $x - \delta_2 < t < x$. Let $\delta = \min(\delta_1, \delta_2)$ (we may also take $\delta \leq \min(x - a, b - x)$). If $|t - x| < \delta$, there are three possibilities: $t > x$, and then $x < t < x + \delta_1$; $t < x$, and then $x - \delta_2 < t < x$; or $t = x$, and then $|f(t) - f(x)| = 0$. In each case $|f(t) - f(x)| < \varepsilon$. Thus $f$ is continuous at $x$. In words: if each half-window works, the smaller of the two half-windows, taken on both sides, gives a full window.
\end{proof}

\begin{definition}\label{def:kinds}
Let $f$ be defined on $(a,b)$. If $f$ is discontinuous at a point $x$, and if $f(x+)$ and $f(x-)$ exist, then $f$ is said to have a \textit{discontinuity of the first kind}, or a \textit{simple discontinuity}, at $x$. Otherwise the discontinuity is said to be of the \textit{second kind}.

There are two ways in which a function can have a simple discontinuity: either $f(x+) \neq f(x-)$ (in which case the value $f(x)$ is immaterial), or $f(x+) = f(x-) \neq f(x)$.
\end{definition}

In the examples below we use the density of the rationals and of the irrationals in $\mathbb{R}$ (Example 1.51): every open interval, however small, contains rational numbers as well as irrational numbers. Consequently, for every $x$ and every $n$ we may choose a rational $t_n$ and an irrational $s_n$ in $(x, x + \frac{1}{n})$, and similarly in $(x - \frac1n, x)$; such sequences tend to $x$ from the right (respectively, from the left).

\begin{example}\label{ex:rudin}
\textbf{(a) The Dirichlet function.} Define
\[
f(x) = \begin{cases} 1 & (x \text{ rational}), \\ 0 & (x \text{ irrational}). \end{cases}
\]
Then $f$ has a discontinuity of the second kind at every point $x$, since neither $f(x+)$ nor $f(x-)$ exists.

\textit{Explanation.} Fix $x \in \mathbb{R}$. Choose rationals $t_n \in (x, x + \frac1n)$ and irrationals $s_n \in (x, x + \frac1n)$. Both sequences lie to the right of $x$ and tend to $x$. But $f(t_n) = 1 \to 1$ and $f(s_n) = 0 \to 0$. If $f(x+) = q$ existed, Definition \ref{def:onesided} would force $f(t_n) \to q$ and $f(s_n) \to q$, so $q = 1$ and $q = 0$, which is absurd. So $f(x+)$ does not exist. In the $\varepsilon$-$\delta$ language: no $q$ works with $\varepsilon = \frac12$, because every right half-window $(x, x+\delta)$ contains points where $f = 1$ and points where $f = 0$, and these two values cannot both lie within $\frac12$ of the same $q$. The same argument with points in $(x - \frac1n, x)$ shows that $f(x-)$ does not exist. Hence $f$ is discontinuous at $x$ (Theorem \ref{thm:contonesided}), and the discontinuity is of the second kind.

\textbf{(b)} Define
\[
f(x) = \begin{cases} x & (x \text{ rational}), \\ 0 & (x \text{ irrational}). \end{cases}
\]
Then $f$ is continuous at $x = 0$ and has a discontinuity of the second kind at every other point.

\textit{Continuity at $0$.} We have $f(0) = 0$, and for every $t$, $|f(t) - f(0)| = |f(t)|$ is either $|t|$ or $0$; so $|f(t) - f(0)| \leq |t|$. Given $\varepsilon > 0$, take $\delta = \varepsilon$: if $|t| < \delta$ then $|f(t) - f(0)| < \varepsilon$.

\textit{Discontinuity at $x \neq 0$.} Choose rationals $t_n$ and irrationals $s_n$ in $(x, x + \frac{1}{n})$. Then $f(t_n) = t_n \to x$ while $f(s_n) = 0 \to 0$. Since $x \neq 0$, the two limits differ, so $f(x+)$ does not exist. Similarly, with points in $(x - \frac1n, x)$, $f(x-)$ does not exist. Thus the discontinuity at $x$ is of the second kind.

\textbf{(c)} Define
\[
f(x) = \begin{cases}
x + 2 & (-3 < x < -2), \\
-x - 2 & (-2 \leq x < 0), \\
x + 2 & (0 \leq x < 1).
\end{cases}
\]
Then $f$ has a simple discontinuity at $x = 0$ and is continuous at every other point of $(-3,1)$.

\textit{Points other than $-2$ and $0$.} If $-3 < x < -2$, then on the neighborhood $N_r(x)$ with $r = \min(x + 3, -2 - x) > 0$, $f$ coincides with the polynomial $t + 2$, which is continuous; by Remark \ref{rem:local}(b), $f$ is continuous at $x$. In the same way, $f$ coincides with $-t - 2$ near each $x \in (-2, 0)$ and with $t + 2$ near each $x \in (0,1)$, so it is continuous there.

\textit{The point $-2$.} For $-3 < t < -2$, $f(t) = t + 2$ and $|f(t) - 0| = |t + 2| < \varepsilon$ if $-2 - \varepsilon < t < -2$; so $f(-2-) = 0$ (Theorem \ref{thm:onesidedeps} with $\delta = \varepsilon$). For $-2 < t < 0$, $f(t) = -t - 2$ and $|f(t) - 0| = |t + 2| < \varepsilon$ if $-2 < t < -2 + \varepsilon$; so $f(-2+) = 0$. And $f(-2) = -(-2) - 2 = 0$. By Theorem \ref{thm:contonesided}, $f$ is continuous at $-2$: the two pieces fit together there.

\textit{The point $0$.} For $-2 < t < 0$, $f(t) = -t - 2$ and $|f(t) - (-2)| = |t| < \varepsilon$ if $-\varepsilon < t < 0$; so $f(0-) = -2$. For $0 < t < 1$, $f(t) = t + 2$ and $|f(t) - 2| = |t| < \varepsilon$ if $0 < t < \varepsilon$; so $f(0+) = 2$. Also $f(0) = 2$. Both one-sided limits exist but $f(0+) = 2 \neq -2 = f(0-)$, so $f$ is discontinuous at $0$ (Theorem \ref{thm:contonesided}), and the discontinuity is simple (a jump of size $4$). Note that $f(0+) = f(0)$: $f$ is ``continuous from the right'' at $0$.

\textbf{(d)} Define
\[
f(x) = \begin{cases} \sin \dfrac{1}{x} & (x \neq 0), \\[2mm] 0 & (x = 0). \end{cases}
\]
Since neither $f(0+)$ nor $f(0-)$ exists, $f$ has a discontinuity of the second kind at $x = 0$. We have not yet shown that $\sin x$ is a continuous function; if we assume this result for the moment (see Remark \ref{rem:sincos}), $f$ is continuous at every point $x \neq 0$.

\textit{Continuity at $x \neq 0$.} On the neighborhood $N_{|x|}(x)$, which does not contain $0$, $f(t) = \sin(1/t)$. The function $t \mapsto 1/t$ is continuous on $\mathbb{R}\setminus\{0\}$ (Example \ref{ex:polynomials}) and $\sin$ is continuous, so their composition is continuous (Theorem \ref{thm:composition}). By Remark \ref{rem:local}(b), $f$ is continuous at $x$.

\textit{No right-hand limit at $0$.} Let
\[
t_n = \frac{1}{\frac{\pi}{2} + 2n\pi}, \qquad s_n = \frac{1}{\frac{\pi}{2} + \pi + 2n\pi} \qquad (n = 1, 2, \dots).
\]
Both are positive and tend to $0$. But
\[
f(t_n) = \sin\Big(\frac{\pi}{2} + 2n\pi\Big) = 1, \qquad f(s_n) = \sin\Big(\frac{3\pi}{2} + 2n\pi\Big) = -1.
\]
So $f(t_n) \to 1$ and $f(s_n) \to -1$; by Definition \ref{def:onesided}, $f(0+)$ does not exist.

\textit{No left-hand limit at $0$.} The sequences $-t_n$ and $-s_n$ are negative and tend to $0$, and since $\sin$ is odd, $f(-t_n) = -1$ and $f(-s_n) = 1$. So $f(0-)$ does not exist.

Thus $f$ is discontinuous at $0$, with a discontinuity of the second kind. Near $0$ the graph oscillates between $-1$ and $1$ infinitely often.

\textit{The intermediate value property (Remark \ref{rem:ivtconverse}).} Let $x_1 < x_2$ and let $c$ lie strictly between $f(x_1)$ and $f(x_2)$; note $c \in [-1,1]$, since all values of $f$ lie in $[-1,1]$. If $0 \notin [x_1, x_2]$, then $f$ is continuous on $[x_1, x_2]$ and Theorem \ref{thm:ivt} gives $x \in (x_1,x_2)$ with $f(x) = c$. If $0 \in [x_1, x_2]$, then $(x_1, x_2)$ contains an interval $(0, \eta)$ or $(-\eta, 0)$ with $\eta > 0$. Say it contains $(0, \eta)$. For large $n$ both $s_n$ and $t_n$ lie in $(0, \eta)$; on the closed interval $[s_n, t_n]$, $f$ is continuous with $f(s_n) = -1$ and $f(t_n) = 1$, so by Theorem \ref{thm:ivt} it takes every value in $[-1,1]$, in particular $c$, at some point of $(s_n, t_n) \subset (x_1,x_2)$. (If only $(-\eta,0)$ is available, use $[-t_n, -s_n]$.) So $f$ has the intermediate value property on every interval, yet it is not continuous at $0$.
\end{example}

\begin{example}\label{ex:possibilities}
\textbf{The possible behaviours at a point.} We now list examples showing that every combination of existence and non-existence of $f(x-)$, $f(x+)$, and agreement with $f(x)$, can occur. In each case the point in question is $x = 0$, the functions are defined on $(-1,1)$ (except in (vi)), and the one-sided limits are checked with Theorem \ref{thm:onesidedeps} or with sequences as in Example \ref{ex:rudin}.
\begin{enumerate}[label=(\roman*)]
    \item \textbf{Left limit exists, right limit does not.}
    \[
    f(x) = \begin{cases} 0 & (-1 < x \leq 0), \\ \sin \frac{1}{x} & (0 < x < 1). \end{cases}
    \]
    For $-1 < t < 0$, $f(t) = 0$, so $|f(t) - 0| = 0 < \varepsilon$ for any $\delta$: $f(0-) = 0 = f(0)$. For $t > 0$, the sequences $t_n$, $s_n$ of Example \ref{ex:rudin}(d) give $f(t_n) = 1$ and $f(s_n) = -1$, so $f(0+)$ does not exist. A discontinuity of the second kind, although $f$ is continuous from the left at $0$.

    \item \textbf{Right limit exists, left limit does not.} The mirror image:
    \[
    f(x) = \begin{cases} \sin \frac{1}{x} & (-1 < x < 0), \\ 0 & (0 \leq x < 1). \end{cases}
    \]
    Now $f(0+) = 0 = f(0)$, while $f(0-)$ does not exist (use $-t_n$ and $-s_n$).

    \item \textbf{Both one-sided limits exist but are different (a jump).}
    \[
    f(x) = \begin{cases} -1 & (-1 < x < 0), \\ c & (x = 0), \\ 1 & (0 < x < 1). \end{cases}
    \]
    Here $f(t) = -1$ for all $t$ to the left of $0$ and $f(t) = 1$ for all $t$ to the right, so $f(0-) = -1$ and $f(0+) = 1$. Whatever the value $c = f(0)$, $f$ is discontinuous at $0$ (Theorem \ref{thm:contonesided}); the discontinuity is simple. If $c = -1$, $f$ is continuous from the left ($f(0-) = f(0)$); if $c = 1$, continuous from the right; if, say, $c = 0$, neither.

    \item \textbf{Both one-sided limits exist and are equal, but differ from the value (removable discontinuity).}
    \[
    f(x) = \begin{cases} 0 & (x \neq 0), \\ 1 & (x = 0). \end{cases}
    \]
    $f(0-) = f(0+) = 0$, so $\lim_{t\to 0} f(t) = 0$ (Theorem \ref{thm:contonesided}(a)), but $f(0) = 1$. A simple discontinuity. Redefining $f(0) = 0$ removes it.

    \item \textbf{Both one-sided limits exist and are equal to the value.} By Theorem \ref{thm:contonesided}(b) this is exactly continuity at the point; e.g., $f(x) = x$, or the function of Example \ref{ex:rudin}(c) at $-2$.

    \item \textbf{Both one-sided limits exist and are equal, but the function has no value at the point.} Let $f(x) = x \sin\frac{1}{x}$ on $(-1,0) \cup (0,1)$; $f$ is not defined at $0$. Since $|f(t) - 0| \leq |t|$, taking $\delta = \varepsilon$ shows $f(0+) = f(0-) = 0$, so $\lim_{t\to 0} f(t) = 0$. As $0$ is not in the domain, it makes no sense to ask whether $f$ is continuous at $0$ (Remark \ref{rem:contcompare}); but setting $f(0) = 0$ produces a function continuous at $0$. Contrast this with $\sin\frac1x$, which cannot be made continuous at $0$ by any choice of $f(0)$, since its one-sided limits do not exist.

    \item \textbf{Neither one-sided limit exists.} Example \ref{ex:rudin}(d) at $0$, or the Dirichlet function of Example \ref{ex:rudin}(a) at any point.
\end{enumerate}
\end{example}

\subsection{Monotonic Functions}

\begin{definition}\label{def:monotone}
Let $f$ be real on $(a,b)$. Then $f$ is said to be \textit{monotonically increasing} on $(a,b)$ if $a < x < y < b$ implies $f(x) \leq f(y)$. If the last inequality is reversed, we obtain the definition of a \textit{monotonically decreasing} function: $a < x < y < b$ implies $f(x) \geq f(y)$. The class of \textit{monotonic} functions consists of both the increasing and the decreasing functions. If $a < x < y < b$ implies $f(x) < f(y)$ (respectively $f(x) > f(y)$), $f$ is said to be \textit{strictly} increasing (respectively strictly decreasing).
\end{definition}

\begin{example}\label{ex:monotone}
\begin{enumerate}
    \item \textbf{Increasing, with a constant part.} Let
    \[
    f(x) = \begin{cases} x & (x < 0), \\ 0 & (0 \leq x \leq 1), \\ x - 1 & (x > 1). \end{cases}
    \]
    $f$ is monotonically increasing on $\mathbb{R}$: on each of the three pieces it is increasing, and the pieces fit in order, since $f < 0$ on $(-\infty,0)$, $f = 0$ on $[0,1]$ and $f > 0$ on $(1,\infty)$. So if $x < y$ lie in different pieces, again $f(x) \leq f(y)$. It is not strictly increasing, because it is constant on $[0,1]$. Monotonicity allows the function to stay constant for a while.
    \item \textbf{A step function.} The greatest integer function $f(x) = [x]$ (the largest integer $\leq x$) is monotonically increasing on $\mathbb{R}$: if $x < y$ then $[x] \leq x < y$, so $[x]$ is an integer $\leq y$, and hence $[x] \leq [y]$. It is constant on each interval $[n, n+1)$ and jumps by $1$ at each integer.
    \item \textbf{Constant functions.} A constant function $f = c$ satisfies both $f(x) \leq f(y)$ and $f(x) \geq f(y)$ for all $x < y$. So it is \textit{both} monotonically increasing and monotonically decreasing. (Conversely, a function which is both is constant: $f(x) \leq f(y) \leq f(x)$.)
    \item $f(x) = 1/x$ is strictly decreasing on $(0, \infty)$, and $f(x) = x^3$ is strictly increasing on $\mathbb{R}$.
\end{enumerate}
\end{example}

\begin{theorem}\label{thm:monotonelimits}
Let $f$ be monotonically increasing on $(a,b)$. Then $f(x+)$ and $f(x-)$ exist at every point $x$ of $(a,b)$. More precisely,
\[
\sup_{a < t < x} f(t) = f(x-) \leq f(x) \leq f(x+) = \inf_{x < t < b} f(t). \tag{M}
\]
Furthermore, if $a < x < y < b$, then
\[
f(x+) \leq f(y-).
\]
Analogous results evidently hold for monotonically decreasing functions.
\end{theorem}

\begin{proof}
Fix $x \in (a,b)$.

\textbf{Step 1: The supremum $A$ of the values to the left of $x$.}\\
By hypothesis, $f(t) \leq f(x)$ whenever $a < t < x$. So the set of numbers $\{f(t) : a < t < x\}$ is nonempty and bounded above by $f(x)$, and therefore (least upper bound property) it has a least upper bound, which we denote by $A$. Since $f(x)$ is an upper bound and $A$ is the least one, $A \leq f(x)$.

\textbf{Step 2: $f(x-) = A$.}\\
We verify the $\varepsilon$-$\delta$ condition of Theorem \ref{thm:onesidedeps}. Let $\varepsilon > 0$ be given.
\begin{itemize}
    \item Since $A$ is the \textit{least} upper bound, $A - \varepsilon$ is \textit{not} an upper bound of $\{f(t) : a < t < x\}$. Hence there is some $t_0 \in (a,x)$ with $f(t_0) > A - \varepsilon$.
    \item Write $t_0 = x - \delta$, i.e., $\delta = x - t_0 > 0$. Then $a < x - \delta < x$ and
    \[
    A - \varepsilon < f(x - \delta) \leq A.
    \]
    \item Now let $x - \delta < t < x$. Since $f$ is increasing, $f(x - \delta) \leq f(t)$; since $A$ is an upper bound, $f(t) \leq A$. Hence
    \[
    A - \varepsilon < f(x - \delta) \leq f(t) \leq A, \qquad\text{so}\qquad |f(t) - A| < \varepsilon \quad (x - \delta < t < x).
    \]
\end{itemize}
This is precisely the $\varepsilon$-$\delta$ definition of the left-hand limit established in Theorem \ref{thm:onesidedeps}. Hence $f(x-)$ exists and $f(x-) = A = \sup_{a<t<x} f(t)$. Together with Step 1, $f(x-) \leq f(x)$.

The picture: once one value $f(x - \delta)$ is within $\varepsilon$ of the supremum, monotonicity traps all later values $f(t)$, $x - \delta < t < x$, between $f(x-\delta)$ and $A$.

\textbf{Step 3: The right-hand side, $f(x+) = B = \inf_{x<t<b} f(t)$.}\\
The set $\{f(t) : x < t < b\}$ is nonempty and bounded below by $f(x)$, so it has a greatest lower bound $B$, and $f(x) \leq B$. Let $\varepsilon > 0$. Since $B$ is the greatest lower bound, $B + \varepsilon$ is not a lower bound, so there is $t_0 \in (x,b)$ with $f(t_0) < B + \varepsilon$. Put $\delta = t_0 - x > 0$. For $x < t < x + \delta = t_0$, monotonicity and the definition of $B$ give
\[
B \leq f(t) \leq f(x + \delta) < B + \varepsilon, \qquad\text{so}\qquad |f(t) - B| < \varepsilon.
\]
By Theorem \ref{thm:onesidedeps}, $f(x+) = B$, and $f(x) \leq f(x+)$. This completes the proof of (M).

\textbf{Step 4: Shortening the interval does not change the infimum or the supremum.}\\
Now let $a < x < y < b$. We claim
\[
f(x+) = \inf_{x < t < b} f(t) = \inf_{x < t < y} f(t), \qquad f(y-) = \sup_{a < t < y} f(t) = \sup_{x < t < y} f(t).
\]
In each claim the first equality is (M) (applied at $x$ and at $y$). For the second equality in the first claim, let $S_1 = \{f(t) : x < t < b\}$ and $S_2 = \{f(t) : x < t < y\}$.
\begin{itemize}
    \item $S_2 \subset S_1$, and the infimum of a smaller set can only be larger: $\inf S_1 \leq \inf S_2$.
    \item Conversely, $\inf S_2$ is a lower bound for $S_1$. Indeed, let $t \in (x,b)$. If $t < y$, then $f(t) \in S_2$, so $f(t) \geq \inf S_2$. If $t \geq y$, pick any $s \in (x,y)$; then $s < t$, so $f(t) \geq f(s) \geq \inf S_2$. Hence $\inf S_1 \geq \inf S_2$.
\end{itemize}
So $\inf S_1 = \inf S_2$. Intuitively, for an increasing function, the values far to the right of $x$ are larger anyway, and the infimum over $(x,b)$ is decided by the values close to $x$; cutting the interval down to $(x,y)$ loses nothing. The second claim is proved in the same way: the values of $f$ on $(a, x]$ are not larger than any value on $(x,y)$, so they do not affect the supremum.

\textbf{Step 5: $f(x+) \leq f(y-)$.}\\
By Step 4, $f(x+)$ and $f(y-)$ are the infimum and the supremum of $f$ over the \textit{same} interval $(x,y)$. Pick any $s \in (x,y)$. Then
\[
f(x+) = \inf_{x<t<y} f(t) \leq f(s) \leq \sup_{x<t<y} f(t) = f(y-).
\]

\textbf{Decreasing functions.} If $f$ is monotonically decreasing, then $-f$ is monotonically increasing, and $(-f)(x\pm) = -f(x\pm)$. Applying the above to $-f$ we obtain
\[
\inf_{a<t<x} f(t) = f(x-) \geq f(x) \geq f(x+) = \sup_{x<t<b} f(t), \qquad\text{and}\qquad f(x+) \geq f(y-) \ \text{ if } a < x < y < b. \qedhere
\]
\end{proof}

\begin{corollary}\label{cor:nosecondkind}
Monotonic functions have no discontinuities of the second kind.
\end{corollary}

\begin{proof}
Let $f$ be monotonic on $(a,b)$ and suppose $f$ is discontinuous at $x \in (a,b)$. By Theorem \ref{thm:monotonelimits} (applied to $f$ or to $-f$), both $f(x+)$ and $f(x-)$ exist. By Definition \ref{def:kinds}, the discontinuity is of the first kind. So a monotonic function can only have simple discontinuities (jumps).

In fact more is true for increasing $f$: since $f(x-) \leq f(x) \leq f(x+)$, and $f$ is continuous at $x$ exactly when $f(x-) = f(x) = f(x+)$ (Theorem \ref{thm:contonesided}), $f$ is discontinuous at $x$ if and only if $f(x-) < f(x+)$.
\end{proof}

\begin{theorem}\label{thm:countable}
Let $f$ be monotonic on $(a,b)$. Then the set of points of $(a,b)$ at which $f$ is discontinuous is at most countable.
\end{theorem}

\begin{proof}
\textbf{Step 1: Reduction.} Suppose, for the sake of definiteness, that $f$ is increasing. (If $f$ is decreasing, apply the argument to the increasing function $-f$, which has the same points of discontinuity.) Let $E$ be the set of points at which $f$ is discontinuous.

\textbf{Step 2: An open interval for each discontinuity.} For every $x \in E$, $f(x-) < f(x+)$ (Corollary \ref{cor:nosecondkind}). So the open interval
\[
I_x = \big(f(x-), f(x+)\big)
\]
is nonempty. Since the rationals are dense in $\mathbb{R}$, we may choose a rational number $r(x) \in I_x$:
\[
f(x-) < r(x) < f(x+).
\]

\textbf{Step 3: Different points get different rationals.} Let $x_1, x_2 \in E$ with $x_1 \neq x_2$, say $x_1 < x_2$. By Theorem \ref{thm:monotonelimits}, $f(x_1+) \leq f(x_2-)$. Hence
\[
r(x_1) < f(x_1+) \leq f(x_2-) < r(x_2),
\]
so $r(x_1) \neq r(x_2)$. Geometrically, the intervals $I_{x_1}$ and $I_{x_2}$ are disjoint, so the rationals chosen in them are distinct.

\textbf{Step 4: Conclusion.} We have thus established a 1-1 correspondence $x \mapsto r(x)$ between the set $E$ and a subset $r(E)$ of the set of rational numbers. The rational numbers form a countable set (Corollary 1.23.1). A subset of a countable set is either finite or, if infinite, countable (Theorem 1.17). Hence $r(E)$, and therefore $E$, is at most countable.
\end{proof}

\begin{example}\label{ex:densejumps}
\textbf{A monotonic function discontinuous exactly on a countable dense set.} The discontinuities of a monotonic function need not be isolated. In fact, given any countable subset $E$ of $(a,b)$, which may even be dense, we can construct a function $f$, monotonic on $(a,b)$, discontinuous at every point of $E$, and at no other point of $(a,b)$. The most striking case is $E = \mathbb{Q} \cap (a,b)$, which is countable (an infinite subset of $\mathbb{Q}$) and dense in $(a,b)$.

\textbf{The construction.} Arrange the points of $E$ in a sequence $\{x_n\}$ of distinct terms, so that $E = \{x_1, x_2, x_3, \dots\}$. Let $\{c_n\}$ be a sequence of positive numbers such that $\sum c_n$ converges; for instance $c_n = 2^{-n}$. Define
\[
f(x) = \sum_{x_n < x} c_n \qquad (a < x < b).
\]
The meaning of this summation is as follows: for $x \in (a,b)$, let $N(x) = \{n \in \mathbb{N} : x_n < x\}$ be the set of indices of those points of $E$ which lie to the left of $x$; then $f(x)$ is the sum of the terms $c_n$ with $n \in N(x)$. If there are no points $x_n$ to the left of $x$, the sum is empty, and $f(x) = 0$. In words: each point $x_n$ of $E$ carries a weight $c_n$, and $f(x)$ is the total weight lying strictly to the left of $x$.

\textbf{Facts about series used.} We use the following facts about a convergent series $\sum c_n$ of positive terms.
\begin{enumerate}[label=(S\arabic*)]
    \item For any set of indices $N' \subset \mathbb{N}$, the subseries $\sum_{n \in N'} c_n$ (the terms with $n \in N'$, taken in increasing order of $n$) converges, and $0 \leq \sum_{n\in N'} c_n \leq \sum_{n=1}^\infty c_n$. Indeed its partial sums increase and are bounded by $\sum c_n$.
    \item If $N' = N_1 \cup N_2$ with $N_1 \cap N_2 = \emptyset$, then $\sum_{n \in N'} c_n = \sum_{n \in N_1} c_n + \sum_{n\in N_2} c_n$. (Think of $\sum_{n\in N'} c_n$ as the full series in which the terms with $n \notin N'$ are replaced by $0$; inserting zeros does not change convergence or the sum, and the termwise sum of two convergent series converges to the sum of their sums.) More generally, a convergent series of positive terms is absolutely convergent, and the terms of an absolutely convergent series may be rearranged (for example, grouped according to where the points $x_n$ lie) without changing the sum. In particular, the order in which we add the terms $c_n$ with $x_n < x$ is immaterial, and $f$ is well defined.
    \item (Tails) Given $\varepsilon > 0$, there is $N$ such that $\sum_{n > N} c_n < \varepsilon$, since the partial sums $s_N$ converge to $s = \sum c_n$ and $\sum_{n>N} c_n = s - s_N$. Consequently, by (S1), $\sum_{n \in N'} c_n < \varepsilon$ for every set of indices $N'$ all of whose elements exceed $N$.
\end{enumerate}

\textbf{(a) $f$ is monotonically increasing.} Let $a < x < y < b$. Then $N(x) \subset N(y)$, and $N(y) \setminus N(x)$ consists of the indices $n$ with $x \leq x_n < y$. By (S2) and (S1),
\[
f(y) - f(x) = \sum_{x \leq x_n < y} c_n \geq 0. \tag{D}
\]
(If $E$ is dense, there is some $x_n \in (x,y)$, and then $f(y) - f(x) \geq c_n > 0$: $f$ is even strictly increasing.)

\textbf{A choice of $\delta$.} Fix $x \in (a,b)$ and $\varepsilon > 0$. By (S3), choose $N$ with $\sum_{n>N} c_n < \varepsilon$; if $x \in E$, say $x = x_k$, choose moreover $N \geq k$. Let
\[
\delta = \min\big( \{x - a,\ b - x\} \cup \{ |x - x_n| : 1 \leq n \leq N,\ x_n \neq x \} \big).
\]
This is a minimum of finitely many positive numbers, so $\delta > 0$. By this choice, $(x - \delta, x + \delta) \subset (a,b)$, and the interval $(x - \delta, x + \delta)$ contains \textit{none} of the finitely many points $x_1, \dots, x_N$, except possibly $x$ itself. Therefore: \textit{every point $x_n \neq x$ lying in $(x-\delta, x+\delta)$ has index $n > N$}, and the total weight of such points is less than $\varepsilon$.

\textbf{(b) $f(x-) = f(x)$ at every $x \in (a,b)$.} Let $x - \delta < t < x$. By (D),
\[
f(x) - f(t) = \sum_{t \leq x_n < x} c_n.
\]
Each $x_n$ occurring here lies in $[t, x) \subset (x - \delta, x)$ and differs from $x$, so $n > N$. By (S3), $0 \leq f(x) - f(t) < \varepsilon$. By Theorem \ref{thm:onesidedeps}, $f(x-) = f(x)$. So $f$ is continuous from the left everywhere.

\textbf{(c) Right-hand limits.} Let $x < t < x + \delta$. By (D),
\[
f(t) - f(x) = \sum_{x \leq x_n < t} c_n.
\]
The points $x_n$ occurring here are: possibly $x$ itself (if $x \in E$), and points of $(x, t) \subset (x, x+\delta)$, whose indices exceed $N$.
\begin{itemize}
    \item \textit{If $x \notin E$}, all the indices exceed $N$, so $0 \leq f(t) - f(x) < \varepsilon$. By Theorem \ref{thm:onesidedeps}, $f(x+) = f(x)$. Together with (b), $f(x-) = f(x) = f(x+)$, and $f$ is \textit{continuous at $x$} (Theorem \ref{thm:contonesided}).
    \item \textit{If $x = x_k \in E$}, the term with $n = k$ is present, and all other indices exceed $N$ (recall $N \geq k$). Hence, by (S2) and (S3),
    \[
    f(t) - f(x_k) = c_k + \sum_{x_k < x_n < t} c_n, \qquad 0 \leq f(t) - \big(f(x_k) + c_k\big) < \varepsilon.
    \]
    By Theorem \ref{thm:onesidedeps}, $f(x_k+) = f(x_k) + c_k$. Combined with (b),
    \[
    f(x_k+) - f(x_k-) = c_k > 0,
    \]
    so $f$ is \textit{discontinuous at $x_k$}: it jumps by exactly the weight $c_k$ sitting at $x_k$.
\end{itemize}

\textbf{Conclusion.} $f$ is monotonically increasing on $(a,b)$, discontinuous at every point of $E$ (with jump $c_k$ at $x_k$), and continuous at every point of $(a,b) \setminus E$. With $E = \mathbb{Q} \cap (a,b)$ and $c_n = 2^{-n}$, we obtain a strictly increasing function which is discontinuous at every rational point and continuous at every irrational point of $(a,b)$. Its set of discontinuities is countable, as Theorem \ref{thm:countable} demands, and yet dense: so Theorem \ref{thm:countable} cannot be improved to say that the discontinuities are isolated.
\end{example}

\subsection{Infinite Limits and Limits at Infinity}\label{sec:infinite}

\begin{remark}\label{rem:whynbhd}
Recall from Remark \ref{rem:contcompare} that the $\varepsilon$-$\delta$ definition of a limit is the same as the neighborhood definition: $\lim_{x\to p} f(x) = q$ if and only if for every neighborhood $U$ of $q$ there is a neighborhood $V$ of $p$ such that $f(x) \in U$ for all $x \in V \cap E$, $x \neq p$.

We would like to give a meaning to statements such as ``$f(t) \to +\infty$ as $t \to 0$'' or ``$f(t) \to 0$ as $t \to +\infty$''. The symbols $+\infty$ and $-\infty$ are not real numbers, and there is no distance between a real number and $+\infty$; so there are no $\delta$-neighborhoods $(x - \delta, x+\delta)$ of the points $\pm\infty$, and the $\varepsilon$-$\delta$ definition cannot even be written down. But the neighborhood formulation does not mention distances at all: it only needs to know what a \textit{neighborhood} of a point is. So, to define limits at $\pm\infty$ and limits equal to $\pm\infty$, it is enough to define neighborhoods of $+\infty$ and of $-\infty$. This is what we do now, working in the extended real number system $\overline{\mathbb{R}} = \mathbb{R} \cup \{+\infty, -\infty\}$.
\end{remark}

\begin{definition}\label{def:nbhdinfty}
For any real $c$, the set of real numbers $x$ such that $x > c$ is called a \textit{neighborhood of $+\infty$} and is written $(c, +\infty)$. Similarly, the set $(-\infty, c)$ is a \textit{neighborhood of $-\infty$}. As before, a neighborhood of a real number $x$ is an interval $(x - r, x + r)$ with $r > 0$.
\end{definition}

The idea: a real number $t$ is ``close to $+\infty$'' when it is large, that is, when it lies in $(c, +\infty)$ for a large $c$; the larger $c$, the smaller the neighborhood.

\begin{definition}\label{def:extlimit}
Let $f$ be a real function defined on $E \subset \mathbb{R}$, and let $x$ and $A$ be points of the extended real number system. Assume that every neighborhood of $x$ contains a point of $E$ different from $x$ (for real $x$ this says that $x$ is a limit point of $E$; for $x = +\infty$ it says that $E$ is not bounded above; for $x = -\infty$, that $E$ is not bounded below). We say that
\[
f(t) \to A \quad \text{as } t \to x,
\]
if for every neighborhood $U$ of $A$ there is a neighborhood $V$ of $x$ such that $V \cap E$ is not empty, and such that $f(t) \in U$ for all $t \in V \cap E$, $t \neq x$.
\end{definition}

\begin{remark}\label{rem:extconcrete}
The definition is the same as before; only the neighborhoods are new. Writing out the neighborhoods explicitly, we obtain the concrete forms of the definition.
\begin{enumerate}
    \item \textbf{$x$ and $A$ real.} Here $U = (A - \varepsilon, A + \varepsilon)$ and $V = (x - \delta, x + \delta)$, and $f(t) \in U$ means $|f(t) - A| < \varepsilon$, while $t \in V$, $t \neq x$ means $0 < |t - x| < \delta$. So the definition says: for every $\varepsilon > 0$ there is $\delta > 0$ with $|f(t) - A| < \varepsilon$ for $t \in E$, $0 < |t - x| < \delta$. This coincides with Definition \ref{def:limit} (with $X = Y = \mathbb{R}$).
    \item \textbf{$x = +\infty$, $A$ real.} For every $\varepsilon > 0$ there is a real $c$ such that $|f(t) - A| < \varepsilon$ for all $t \in E$ with $t > c$. For example, $1/t \to 0$ as $t \to +\infty$ (on $E = (0,\infty)$): given $\varepsilon$, take $c = 1/\varepsilon$.
    \item \textbf{$x$ real, $A = +\infty$.} For every real $M$ there is $\delta > 0$ such that $f(t) > M$ for all $t \in E$ with $0 < |t - x| < \delta$. For example, $1/t^2 \to +\infty$ as $t \to 0$: given $M > 0$, take $\delta = 1/\sqrt{M}$.
    \item \textbf{$x = +\infty$, $A = +\infty$.} For every real $M$ there is a real $c$ such that $f(t) > M$ for all $t \in E$ with $t > c$. For example, $t^2 \to +\infty$ as $t \to +\infty$ (and also as $t \to -\infty$).
\end{enumerate}
The cases involving $-\infty$ are obtained by replacing ``$t > c$'' by ``$t < c$'' and ``$f(t) > M$'' by ``$f(t) < M$''.

\textbf{A convenient phrase.} Let $x \in \overline{\mathbb{R}}$ be as in Definition \ref{def:extlimit}. We shall say that a statement $P(t)$ holds \textit{near $x$} if there is a neighborhood $V$ of $x$ such that $P(t)$ holds for all $t \in V \cap E$ with $t \neq x$. Then, by Definition \ref{def:extlimit} and the description of neighborhoods:
\begin{itemize}
    \item for real $A$: $f(t) \to A$ iff for every $\varepsilon > 0$, ``$|f(t) - A| < \varepsilon$'' holds near $x$ (every neighborhood of $A$ contains an interval $(A-\varepsilon, A+\varepsilon)$, and is one);
    \item $f(t) \to +\infty$ iff for every real $M$, ``$f(t) > M$'' holds near $x$;
    \item $f(t) \to -\infty$ iff for every real $M$, ``$f(t) < M$'' holds near $x$.
\end{itemize}
\end{remark}

\begin{lemma}\label{lem:nbhd}
Let $x, A, A', C \in \overline{\mathbb{R}}$.
\begin{enumerate}[label=(\alph*)]
    \item The intersection of two neighborhoods of $x$ is again a neighborhood of $x$. Consequently, if $P(t)$ holds near $x$ and $Q(t)$ holds near $x$, then ``$P(t)$ and $Q(t)$'' holds near $x$.
    \item If $A \neq A'$, there are a neighborhood $U$ of $A$ and a neighborhood $U'$ of $A'$ with $U \cap U' = \emptyset$.
    \item For a function $h$ on $E$: $h(t) \to C$ as $t \to x$ if and only if $-h(t) \to -C$ as $t \to x$ (with $-(+\infty) = -\infty$ and $-(-\infty) = +\infty$).
\end{enumerate}
\end{lemma}

\begin{proof}
(a) If $x$ is real: $(x - r_1, x + r_1) \cap (x - r_2, x + r_2) = (x - r, x + r)$ with $r = \min(r_1, r_2)$. If $x = +\infty$: $(c_1, +\infty) \cap (c_2, +\infty) = (\max(c_1,c_2), +\infty)$. If $x = -\infty$: $(-\infty, c_1) \cap (-\infty, c_2) = (-\infty, \min(c_1, c_2))$. For the consequence: if $P$ holds on $V_1 \cap E \setminus\{x\}$ and $Q$ on $V_2 \cap E \setminus\{x\}$, both hold on $(V_1 \cap V_2) \cap E \setminus\{x\}$, and $V_1 \cap V_2$ is a neighborhood of $x$.

(b) If $A < A'$ are both real, take radius $r = (A' - A)/2$: $(A - r, A + r)$ and $(A' - r, A' + r)$ are disjoint. If $A$ is real and $A' = +\infty$: take $(A - 1, A + 1)$ and $(A + 1, +\infty)$. If $A$ is real and $A' = -\infty$: take $(A-1, A+1)$ and $(-\infty, A - 1)$. If $A = -\infty$ and $A' = +\infty$: take $(-\infty, 0)$ and $(0, +\infty)$.

(c) The map $u \mapsto -u$ carries the neighborhoods of $C$ exactly onto the neighborhoods of $-C$: $(C - r, C + r) \mapsto (-C - r, -C + r)$, $(c, +\infty) \mapsto (-\infty, -c)$ and $(-\infty, c) \mapsto (-c, +\infty)$. And $h(t) \in U$ if and only if $-h(t) \in -U$. So the condition of Definition \ref{def:extlimit} for $h$ and $C$ is the same as that for $-h$ and $-C$.
\end{proof}

\textbf{Arithmetic in $\overline{\mathbb{R}}$.} We use the following conventions. For real $a$:
\[
a + (+\infty) = +\infty, \quad a + (-\infty) = -\infty, \quad (+\infty) + (+\infty) = +\infty, \quad (-\infty) + (-\infty) = -\infty,
\]
\[
a \cdot (+\infty) = +\infty, \ a \cdot (-\infty) = -\infty \ \text{ if } a > 0; \qquad a \cdot (+\infty) = -\infty, \ a \cdot (-\infty) = +\infty \ \text{ if } a < 0;
\]
\[
(+\infty)(+\infty) = (-\infty)(-\infty) = +\infty, \quad (+\infty)(-\infty) = -\infty, \quad \frac{a}{\pm\infty} = 0, \quad \frac{\pm\infty}{b} = (\pm\infty)\cdot\frac{1}{b} \ \ (b \text{ real}, b \neq 0).
\]
The expressions $\infty - \infty$, $0 \cdot \infty$, $\infty/\infty$ and $A/0$ are \textit{not defined}.

\begin{theorem}\label{thm:extalgebra}
Let $f$ and $g$ be defined on $E \subset \mathbb{R}$, and let $x \in \overline{\mathbb{R}}$ be as in Definition \ref{def:extlimit}. Suppose
\[
f(t) \to A, \qquad g(t) \to B \qquad \text{as } t \to x.
\]
Then
\begin{enumerate}[label=(\alph*)]
    \item $f(t) \to A'$ implies $A' = A$.
    \item $(f + g)(t) \to A + B$,
    \item $(fg)(t) \to AB$,
    \item $(f/g)(t) \to A/B$,
\end{enumerate}
provided the right members of (b), (c), and (d) are defined.
\end{theorem}

When $x$, $A$, $B$ are all real, (a)--(d) are Corollary \ref{cor:unique} and Theorem \ref{thm:limitalgebra}. The proofs below are similar to those, but now $x$ may be $\pm\infty$ and $A$, $B$ may be $\pm\infty$, and we treat all cases. We use the phrase ``near $x$'' of Remark \ref{rem:extconcrete} throughout; by Lemma \ref{lem:nbhd}(a), finitely many statements which each hold near $x$ hold simultaneously near $x$. This plays the role of ``take $\delta = \min(\delta_1, \delta_2)$'' in the earlier proofs.

\begin{proof}
\textbf{Part (a): Uniqueness.}\\
Suppose $f(t) \to A$ and $f(t) \to A'$ with $A \neq A'$. By Lemma \ref{lem:nbhd}(b), choose disjoint neighborhoods $U$ of $A$ and $U'$ of $A'$. Then ``$f(t) \in U$'' holds near $x$ and ``$f(t) \in U'$'' holds near $x$, so by Lemma \ref{lem:nbhd}(a) both hold for all $t \in V \cap E$, $t \neq x$, for a suitable neighborhood $V$ of $x$. By our assumption on $x$, $V$ contains a point $t \in E$ with $t \neq x$; for this $t$, $f(t) \in U \cap U' = \emptyset$, which is impossible. Hence $A' = A$.

\textbf{Part (b): Sums.} The sum $A + B$ is defined unless one of $A, B$ is $+\infty$ and the other is $-\infty$.

\textit{Case 1: $A$ and $B$ real.} Let $\varepsilon > 0$. Near $x$, $|f(t) - A| < \varepsilon/2$, and near $x$, $|g(t) - B| < \varepsilon/2$. Hence near $x$ both hold, and
\[
|(f+g)(t) - (A+B)| \leq |f(t) - A| + |g(t) - B| < \varepsilon.
\]
So $(f+g)(t) \to A + B$.

\textit{Case 2: $A = +\infty$ and $B \neq -\infty$ (so $B$ is real or $+\infty$).} Then $A + B = +\infty$. First, $g$ is bounded below near $x$: if $B$ is real, then $|g(t) - B| < 1$, hence $g(t) > B - 1$, near $x$; if $B = +\infty$, then $g(t) > 0$ near $x$. In either case there is a real $K$ with $g(t) > K$ near $x$. Now let $M$ be real. Since $f(t) \to +\infty$, $f(t) > M - K$ near $x$. Hence near $x$,
\[
f(t) + g(t) > (M - K) + K = M.
\]
So $(f+g)(t) \to +\infty$.

\textit{Case 3: $A = -\infty$ and $B \neq +\infty$.} By Lemma \ref{lem:nbhd}(c), $-f(t) \to +\infty$ and $-g(t) \to -B \neq -\infty$. By Case 2, $(-f) + (-g) = -(f+g) \to +\infty$, and by Lemma \ref{lem:nbhd}(c) again, $(f+g)(t) \to -\infty = A + B$.

\textit{Case 4: $A$ real and $B = \pm\infty$.} Interchange the roles of $f$ and $g$ and apply Case 2 or Case 3.

These cases exhaust all situations in which $A + B$ is defined.

\textbf{Part (c): Products.} The product $AB$ is defined unless one of $A, B$ is $0$ and the other is $\pm\infty$.

\textit{Case 1: $A$ and $B$ real.} As in Theorem \ref{thm:limitalgebra}(b): near $x$, $|g(t) - B| < 1$, so $|g(t)| < |B| + 1$. Given $\varepsilon > 0$, near $x$ we also have $|f(t) - A| < \frac{\varepsilon}{2(|B|+1)}$ and $|g(t) - B| < \frac{\varepsilon}{2(|A|+1)}$. Then near $x$,
\[
|f(t)g(t) - AB| \leq |f(t) - A|\,|g(t)| + |A|\,|g(t) - B| < \frac{\varepsilon}{2} + \frac{\varepsilon}{2} = \varepsilon.
\]

\textit{Case 2: $A = +\infty$ and $B > 0$ ($B$ a positive real number or $+\infty$).} Then $AB = +\infty$. First, $g$ is bounded below by a positive constant near $x$: if $B$ is real and positive, $|g(t) - B| < B/2$, hence $g(t) > B/2$, near $x$; if $B = +\infty$, $g(t) > 1$ near $x$. So there is $K > 0$ with $g(t) > K$ near $x$. Let $M > 0$. Near $x$, $f(t) > M/K$ (in particular $f(t) > 0$). Hence near $x$,
\[
f(t)g(t) > \frac{M}{K} \cdot K = M.
\]
(For $M \leq 0$ the inequality $f(t)g(t) > M$ follows from the case $M = 1$.) So $(fg)(t) \to +\infty$.

\textit{Case 3: the remaining cases.} Suppose $AB$ is defined and at least one of $A$, $B$ is infinite. Then neither $A$ nor $B$ is $0$. By interchanging $f$ and $g$ we may assume that $A$ is infinite. Let $\sigma = 1$ if $A = +\infty$ and $\sigma = -1$ if $A = -\infty$; let $\tau = 1$ if $B > 0$ and $\tau = -1$ if $B < 0$. By Lemma \ref{lem:nbhd}(c), $\sigma f(t) \to \sigma A = +\infty$ and $\tau g(t) \to \tau B > 0$. By Case 2,
\[
\sigma\tau\, f(t) g(t) = (\sigma f(t))(\tau g(t)) \to +\infty.
\]
Since $\sigma\tau = \pm 1$, Lemma \ref{lem:nbhd}(c) gives $f(t)g(t) \to \sigma\tau\cdot(+\infty)$, which, by the sign rules above, is exactly $AB$. For instance, if $A = -\infty$ and $B = 3$, then $(-f)g \to +\infty$, so $fg \to -\infty = (-\infty)\cdot 3$.

\textbf{Part (d): Quotients.} The quotient $A/B$ is defined when $B \neq 0$ and $A$, $B$ are not both infinite. We first show that $1/g(t) \to 1/B$, where $1/(\pm\infty) = 0$, and then use (c).

\textit{Step 1: $g(t) \neq 0$ near $x$.} If $B$ is real and nonzero, then near $x$, $|g(t) - B| < |B|/2$, so $|g(t)| > |B|/2 > 0$. If $B = \pm\infty$, then near $x$, $|g(t)| > 1$ (namely $g(t) > 1$, respectively $g(t) < -1$). Hence $f/g$ is defined at all $t \in V \cap E$, $t \neq x$, for some neighborhood $V$ of $x$. Let $E' = \{t \in E : g(t) \neq 0\}$ be the set where $f/g$ is defined. Every neighborhood $W$ of $x$ contains a point of $E'$ different from $x$: indeed $W \cap V$ is a neighborhood of $x$ (Lemma \ref{lem:nbhd}(a)), so it contains a point $t \in E$, $t \neq x$, and $g(t) \neq 0$. So $x$ satisfies the assumption of Definition \ref{def:extlimit} relative to $E'$. On $E'$ we still have $f(t) \to A$ and $g(t) \to B$, since the defining condition now has to be checked at fewer points.

\textit{Step 2: $1/g(t) \to 1/B$.} If $B$ is real and nonzero, then near $x$, as in Theorem \ref{thm:limitalgebra}(c),
\[
\Big|\frac{1}{g(t)} - \frac{1}{B}\Big| = \frac{|g(t) - B|}{|g(t)|\,|B|} < \frac{2}{|B|^2}|g(t) - B|,
\]
which is $< \varepsilon$ near $x$ once $|g(t) - B| < \varepsilon|B|^2/2$ near $x$. If $B = \pm\infty$, given $\varepsilon > 0$ we have $|g(t)| > 1/\varepsilon$ near $x$ (this is $g(t) > 1/\varepsilon$ if $B = +\infty$, and $g(t) < -1/\varepsilon$ if $B = -\infty$), hence $|1/g(t) - 0| < \varepsilon$ near $x$. So $1/g(t) \to 0 = 1/B$.

\textit{Step 3: Apply (c).} On $E'$, $f/g = f \cdot (1/g)$, and $f(t) \to A$, $1/g(t) \to 1/B$. We check that $A \cdot (1/B)$ is defined and equals $A/B$ in every case in which $A/B$ is defined:
\begin{itemize}
    \item $B$ real, $B \neq 0$: $1/B$ is a nonzero real number, so $A \cdot (1/B)$ is defined for every $A$, and it equals $A/B$ (for $A = \pm\infty$ this is our convention for $\pm\infty/b$).
    \item $B = \pm\infty$ and $A$ real: $A \cdot (1/B) = A \cdot 0 = 0 = A/B$.
\end{itemize}
(If $B = \pm\infty$ and $A = \pm\infty$, then $A/B$ is not defined, and correspondingly $A \cdot 0$ is not defined either.) By part (c), $(f/g)(t) \to A \cdot (1/B) = A/B$.
\end{proof}

\begin{remark}
The restrictions in Theorem \ref{thm:extalgebra} are genuine; in the undefined cases anything can happen. For instance, as $t \to +\infty$: with $f(t) = t + c$ and $g(t) = -t$ we have $f \to +\infty$, $g \to -\infty$ and $f + g = c$, which can be any number; with $f(t) = t^2$ and $g(t) = -t$, $f + g \to +\infty$; with $f(t) = t$ and $g(t) = -t^2$, $f + g \to -\infty$. Similarly, with $f(t) = 1/t$ and $g(t) = ct$ ($c > 0$), $f \to 0$, $g \to +\infty$ and $fg = c$. This is why $\infty - \infty$ and $0 \cdot \infty$ are left undefined.
\end{remark}

\end{document}
